\section{The two methods: precise transfers and repairs}
\label{sec:methods}

\subsection*{MT-A01. What the adopted Dirichlet theorem says adelically}
The supplied declarations \texttt{riemannHypothesis},
\texttt{LFunction\_ne\_zero\_of\_half\_lt\_re}, and
\texttt{LFunction\_eq\_zero\_re} in \texttt{ALPOGE\_Main.lean} have the
definitions printed in \texttt{ALPOGE\_Statements.lean}.  We use these as the
cited input $A_D$, as requested.  The supplied archive contains these exported
declarations and the scan report, but not the imported tree
\texttt{Supplements.Kernels.MainTheoremsAll}; a kernel build of that tree is
not a check performed in this run.  This distinction concerns our evidence,
not the stipulated status of the input theorem.

Every unitary idele-class character of $\mathbb Q$ has the form
$\chi_0|\cdot|^{it}$, where $\chi_0$ is a finite-order Dirichlet character
and $t\in\mathbb R$.  Consequently $A_D$ gives the rank-one Hecke GRH over
$\mathbb Q$ for these characters, not just its finite-order specialization.
Indeed the norm map has a continuous section and compact kernel, which for
$\mathbb Q$ is $\widehat{\mathbb Z}^{\times}$ after the sign convention at
infinity.  A continuous character of this profinite group into the circle has
finite image: its image is a compact quotient of a profinite group, hence
totally disconnected, whereas an infinite closed subgroup of the circle is
the circle.  It therefore factors through some
$(\mathbb Z/N\mathbb Z)^\times$.  The remaining unitary character of
$\mathbb R_{>0}$ is $x^{it}$.  In the convention
$L(s,\chi_0|\cdot|^{it})=L(s+it,\chi_0)$, zeros are translated only in
their imaginary coordinate.  This proves the assertion.  Geometrically,
the finite-order part is a rank-one representation of the arithmetic
fundamental group; the continuous norm twist is an analytic normalization,
not an additional finite-monodromy motive.  The result does not cover
arbitrary characters over a nontrivial number field.

\subsection*{MT-M01. A nonzero counterexample to the draft's general zero-count lemma}
Masullo, Lemma 9.3.2.4, PDF p.~425, claims a local $O(\log |T|)$ zero count
for the Mellin transform of every $\varphi\in\mathcal S_E(G)$.
We checked the image of the original PDF page, including the displayed
Jensen inequality and its last inference.  The quantifier is present in the
PDF; this is not an inference from the model transcription.

Take $G=\mathbb R_{>0}$, trivial character, Haar measure $dt/t$, and set
\[
 g_a(u)=\frac{1}{\sqrt{4\pi a}}\exp\left(-\frac{u^2}{4a}\right),
 \qquad a=1+i,\qquad
 \varphi(e^u)=\frac{g_{1+i}(u)-g_{1-i}(u)}{2i}.
\]
The square roots use the branch with positive real part.  This is a real,
nonzero smooth function.  For every real $b$ and every nonnegative integer
$j$, the $j$th derivative of $e^{bu}g_{1\pm i}(u)$ is a polynomial in $u$
times $e^{bu-u^2/(4(1\pm i))}$.  Its absolute value is bounded by a
polynomial times $e^{bu-u^2/8}$, hence is rapidly decreasing.  Thus
$t^b\varphi(t)$ is Schwartz in the logarithmic coordinate for every $b$;
this is exactly the condition of Definition D.3.3.1, PDF p.~804 (also
\S1.1.2, PDF p.~11).  In particular $\varphi\in\mathcal S_E(G)$.

For $\operatorname{Re}a>0$, the Gaussian integral and analytic continuation
in $a$ give
\[
 \int_{\mathbb R}g_a(u)e^{su}\,du=e^{as^2},\qquad s\in\mathbb C.
\]
For completeness, the identity first follows by completing the square for
positive real $a$ and real $s$; both sides are entire in $s$ by Gaussian
domination on compact sets, and holomorphic in $a$ on
$\operatorname{Re}a>0$, where $\operatorname{Re}(1/a)>0$.  The identity
theorem proves the claimed continuation.  Therefore
\[
 \mathcal M\varphi(s)=e^{s^2}\sin(s^2).
\]
Its zeros on the positive imaginary axis are
$i\sqrt{k\pi}$, $k=1,2,\ldots$, and are simple.  For $I=[-1,1]$ and
$T\ge2$, their number in $I+i[T,T+1]$ is
\[
 \#\{k:T^2/\pi\le k\le(T+1)^2/\pi\}
   =\frac{2T+1}{\pi}+O(1).
\]
This is not $O(\log T)$.  The function is nonzero, so removing the zero
function from the lemma does not repair it.  The proof on p.~425 needs a
\emph{lower} bound for the modulus at a Jensen centre, whereas rapid decay
supplies upper bounds.  No statement about the zeros of a Hecke L-function
is contradicted by this example.

\subsection*{MT-M02. The displayed approximate identities are proper kernel projections}
Theorem 9.3.2.6, equation (9.3.2.10), PDF p.~426, defines subspaces by
intersecting kernels of all jet evaluations up to height $T$, then claims
their orthogonal projections tend strongly to the identity.  The printed
formula and the asserted inclusion in its proof have both been checked on
the page image.

Here is a counterexample even with a single zero, so the zero-count issue
in MT-M01 is irrelevant.  Put
\[
 g(u)=(4\pi)^{-1/2}e^{-u^2/4},\qquad
 \psi(e^u)=-g'(u).
\]
Then $\psi\in\mathcal S_E(\mathbb R_{>0})$, and integration by parts gives
$\mathcal M\psi(z)=z e^{z^2}$.  Its only zero is $0$, with multiplicity one,
strictly inside the strip for $I=[-1,1]$.  Choose Sobolev index $0$.
The draft's Hilbert norm is then
\[
 \|f\|_H^2=\int_{\mathbb R}|f(e^u)|^2e^{2|u|}\,du.
\]
The functional $\delta_0(f)=\int_{\mathbb R}f(e^u)\,du$ is continuous,
because Cauchy--Schwarz bounds its modulus by
$\|f\|_H(\int e^{-2|u|}\,du)^{1/2}=\|f\|_H$.
It is nonzero, since $\delta_0(g)=1$.  For every $T>0$ the subspace displayed
in (9.3.2.10) is the same proper closed subspace $\ker\delta_0$.
Its orthogonal projector cannot converge to the identity.  In particular,
the Riesz vector $r(e^u)=e^{-2|u|}$ has $\|r\|_H=1$ and is killed by every
one of these projectors.  It belongs to the indicated completion because
$C_c^\infty(\mathbb R)$ is contained in $\mathcal S_E$ in logarithmic
coordinates and is dense in this weighted $L^2$ space: truncate first and
then smooth on compact intervals, where the positive weight is bounded
above and below.  The ranges decrease when more evaluation
conditions are added, whereas the proof on p.~426 uses the opposite
inclusion.

\subsection*{MT-M03. The exact projection repair}
Let $H$ be a Hilbert space, let $(\ell_j)_{j\ge1}$ be continuous linear
functionals, and put
\[
 V_N=\bigcap_{j\le N}\ker\ell_j,\qquad V=\bigcap_{j\ge1}\ker\ell_j.
\]
Then $P_{V_N}\to P_V$ strongly.  The complementary projectors converge
strongly to $P_{V^\perp}$, the projector onto the closed span of the Riesz
vectors of the $\ell_j$.  Thus they give an approximate identity on the
spectral quotient $H/V\simeq V^\perp$.

Here is a full proof.  For fixed $x\in H$, put $y_N=P_{V_N}x$.  For
$M\ge N$, nestedness gives $P_{V_M}y_N=y_M$ and hence
\[
 \|y_N-y_M\|^2=\|y_N\|^2-\|y_M\|^2.
\]
The nonincreasing nonnegative sequence $\|y_N\|^2$ converges, so $(y_N)$ is
Cauchy.  Write its limit as $y$.  For each fixed $N$, all sufficiently late
$y_M$ lie in the closed space $V_N$, so $y\in V$.  Also
$x-y_N\perp V$ for every $N$; passing to the limit gives $x-y\perp V$.
Therefore $y=P_Vx$.  Taking complements proves the second assertion.
Finally the orthogonal complement of the closed span of the Riesz vectors
is precisely $V$, which proves the last identification.  Exhaustion by
height instead of an integer enumeration gives the same proof.

This repairs a functional-analytic direction error.  It does not identify
the common kernel with the closed adelic theta image, prove a Schauder
basis, establish nuclearity of the actual convolution operators, or justify
the geometric-kernel decomposition.  Those are separate statements and
must be checked in the quotient on which they are used.  In Grothendieck's
trace language this is the distinction between an annihilator and its
spectral quotient, before any trace or Fredholm determinant is taken.

\subsection*{MT-M04. An arithmetic replacement for the false general lemma}
For every fixed unitary Hecke character, the number of nontrivial zeros
in a unit-height window is $O_{K,\eta}(\log(|T|+3))$, unconditionally.
The full proof for arbitrary infinity type is NT-Z10 above: polynomial
vertical growth and a uniform lower bound at the Euler-product centre
give Jensen's estimate.  It includes pure norm twists, whose single
translated pole is excluded before applying Jensen at large height.
For a fixed primitive finite-order character a conductor-explicit
primary locator is Greni\'e--Molteni,
\emph{An explicit Chebotarev density theorem under GRH}, Appendix A,
Theorem A.1, author PDF p.~26: that appendix's zero count is unconditional.
Subtract the counts at consecutive heights.  The main term has derivative
$O(\log(|T|+3))$ for a fixed field and character and each remainder has the
same bound.  The arithmetic restriction therefore closes the local
zero-count gate, including infinite-order characters, without using GRH.
It supplies a route around MT-M01 for arithmetic divisors, but does not
establish Lemma 9.3.2.5 for arbitrary Schwartz convolution ideals.  The
constant in the infinite-order statement is fixed-character, not a
claimed uniform archimedean-conductor constant.

\subsection*{MT-ART01. A nonmonomial GL\texorpdfstring{$_2$}{2} consequence: odd Artin forms}
Assume $M_H$.  Let $\rho:G_{\mathbb Q}\to\operatorname{GL}_2(\mathbb C)$
be continuous, irreducible, odd and of finite image.  Then its Artin
L-function is entire and all of its nontrivial zeros lie on
$\operatorname{Re}s=1/2$.

The holomorphy input is already a theorem: Khare--Wintenberger,
\emph{Serre's modularity conjecture (I)}, Corollary 10.2(ii), author PDF
p.~21 (with Theorem 10.1(ii), pp.~20--21), associates $\rho$ to a weight-one
cuspidal newform.  Its standard L-function is entire.  Let $L$ be the finite
Galois extension cut out by $\rho$.  Brauer induction writes its character
as an integer linear combination of characters induced from
one-dimensional characters of subgroups of $\operatorname{Gal}(L/\mathbb Q)$.
Artin induction and class field theory therefore express $L(s,\rho)$ as
a finite product of integer powers of finite-order Hecke L-functions over
the corresponding fixed fields.  In the open critical strip, $M_H$ makes
each factor holomorphic and nonzero away from the critical line.  The
product, including its negative exponents, is consequently holomorphic and
nonzero there.  Combining this with the independent holomorphy theorem
proves GRH for $L(s,\rho)$.

This includes projective image $A_5$: no monomial expression with only
nonnegative exponents is required.  It is a new transfer in this catalogue,
not a new proof of the classical modularity theorem.  In motivic terms it
concerns odd rank-two Artin motives and their weight-one automorphic
realizations.  It gives no GRH for the weight-two L-function of an arbitrary
non-CM elliptic curve.  Nor does it prove holomorphy of arbitrary even
icosahedral or higher-dimensional Artin representations.  The inherited
Brauer divisor-support statements in RETURN\_ZETA\_TWO are the general
input behind this specialization.

\subsection*{MT-ART02. The missing positivity is a finite matrix at each zero}
Let $L/K$ be a finite Galois extension with group $G$, and let
$0<\operatorname{Re}s_0<1$.  By Brauer meromorphy define
\[
 n_\chi=\operatorname{ord}_{s=s_0}L(s,\chi,L/K)\in\mathbb Z,
 \qquad\Theta_{s_0}=\sum_{\chi\in\operatorname{Irr}(G)}n_\chi\chi.
\]
Form the $|G|$ by $|G|$ matrix
\[
 K_{g,h}=\Theta_{s_0}(g h^{-1}).
\]
Then all Artin L-functions for $L/K$ are holomorphic at $s_0$ if and only
if $K$ is positive semidefinite.  The matrix is Hermitian because
$\chi(g^{-1})=\overline{\chi(g)}$ and $n_\chi$ is real.

To prove the equivalence, consider the central convolution operator
$T=|G|^{-1}K$ on $\mathbb C[G]$.  On the matrix-coefficient space of an
irreducible character $\chi$, Schur orthogonality makes $T$ scalar with
eigenvalue $n_\chi/\chi(1)$ (interchanging $\chi$ and its dual under the
opposite convolution convention only permutes the list).  Each such space
has dimension $\chi(1)^2$ and together they exhaust the regular
representation.  Thus $K\succeq0$ is equivalent to $n_\chi\ge0$ for every
$\chi$, which is exactly absence of a pole of every factor.

The restrictions of $\Theta_{s_0}$ to cyclic subgroups are already actual
characters.  Indeed, for a cyclic subgroup $C$ and a one-dimensional
character $\eta$ of $C$, Frobenius reciprocity and Artin induction give
\[
 \langle\Theta_{s_0}|_C,\eta\rangle_C
   =\operatorname{ord}_{s=s_0}L(s,\eta,L/L^C)\ge0,
\]
because the right-hand side is an order of a holomorphic Hecke function in
the strip.  This proves positivity on every such restriction; it does not
prove the full matrix is positive.  Also
\[
 \Theta_{s_0}(1)=\sum_\chi n_\chi\chi(1)
   =\operatorname{ord}_{s=s_0}\zeta_L(s).
\]
The point of the criterion is the exact remaining object: a central
positive kernel in the Artin character ring.  $M_H$ removes all tests off
the critical line, while neither zero location nor cyclic positivity
supplies the full critical-line kernel.  An Alp\"oge-type positivity
argument would have to produce this kernel or an equivalent positive
decomposition, with the actual Artin orders as coefficients.

\subsection*{MT-ART04. All induced-linear positivity tests still leave a gap}
There is a virtual character which passes every positivity inequality
obtained by inducing a one-dimensional character of a subgroup, but is not
an actual character.  This refutes an abstract reconstruction rule; we do
not claim that this virtual character occurs as the orders of actual
Artin L-functions at a point.

Let $G=A_5$, and let $\chi$ be its four-dimensional augmentation character
for the natural action on five letters.  Define
\[
             \Theta=\operatorname{reg}_G-5\chi.
\]
The natural permutation action is doubly transitive, so its permutation
character has inner-product norm squared two.  Its invariant subspace is
one-dimensional; the orthogonal complement therefore has character
$\chi$ with norm squared one, proving that $\chi$ is irreducible.
The coefficient of $\chi$ in $\Theta$ is $4-5=-1$, so $\Theta$ is not an
actual character and its full matrix from MT-ART02 is not positive.

Nevertheless, for every $H\le G$ and every linear character $\eta$ of $H$,
\[
       \left\langle\Theta,\operatorname{Ind}_H^G\eta\right\rangle_G
                    \ge0.                                      \tag{MT.1}
\]
Here are all the details.  The conjugacy classes of $A_5$ have sizes
$1,15,20,12,12$, as follows by counting the double transpositions,
three-cycles and the two five-cycle classes.  A normal subgroup is a union
of these classes containing the identity, and its size divides $60$.
No proper nontrivial sum of these class sizes containing $1$ divides $60$;
thus $A_5$ is simple.  An action on cosets of a proper subgroup of index
$2,3,$ or $4$ would be an injective homomorphism into the corresponding
symmetric group, which is impossible by order.  Lagrange's theorem leaves
only the possible proper indices
\[
                    k\in\{5,6,10,12,15,20,30,60\}.
\]
Put $m=\langle\chi|_H,\eta\rangle_H$.  Since $\chi|_H$ has dimension four,
$0\le m\le4$.  Frobenius reciprocity says that
$\operatorname{Ind}_H^G\eta$, of dimension $k$, contains $m$ copies of
$\chi$, so $4m\le k$.  These two bounds imply $k-5m\ge0$ for every index
in the list except possibly $k=12$.  At $k=12$, $|H|=5$, and $H$ acts
regularly on the five letters.  The restriction of $\chi$ is the sum of
the four distinct nontrivial characters of $C_5$, so $m\le1$ and again
$k-5m\ge0$.  Finally
\[
 \left\langle\operatorname{reg}_G,
       \operatorname{Ind}_H^G\eta\right\rangle=k,
\]
proving (MT.1) for proper $H$.  For $H=G$, simplicity and nonabelianity
make the trivial character the only linear character, and its coefficient
in $\Theta$ is one.  This proves (MT.1) in all cases.

For comparison with the methods, Artin induction and Hecke holomorphy
give exactly inequalities of the form (MT.1) for an actual local Artin
order character.  The example proves that these inequalities alone do not
force a positive character, even before any question of zero location.
Thus a transfer of determinant positivity from the Alp\"oge method would
need additional arithmetic information, not just all the existing
induced-linear nonnegative orders.  The elementary exact computation
\texttt{checks/artin\_cone\_certificate.py} checks the negative matrix
quadratic form: with $v(g)=\chi(g)$ one has $v^*Kv=-900$.

\subsection*{MT-ART03. The classical simple-zero reduction}
Stark, \emph{Some effective cases of the Brauer--Siegel theorem},
Inventiones Mathematicae \textbf{23} (1974), Theorem 3, p.~144, proves
that if $\operatorname{ord}_{s_0}\zeta_L\le1$ and $s_0\ne1$, then all
Artin factors for $L/K$ are holomorphic at $s_0$.  We cite this classical
result rather than count another proof as a finding.  Applied with the
divisor-support theorem under $M_H$, it confines any still-possible pole
inside the strip to a multiple zero of the top-field zeta on the critical
line.  No simplicity assertion for these Dedekind zetas is being made.

\subsection*{MT-Q01. Exact path quadrics in every characteristic}
For $k\ge2$ let
\[
 V_k=\left\{(x_1,\ldots,x_k):\sum_{i=1}^{k-1}x_ix_{i+1}=0\right\}
       \subset\mathbb A^k_{\mathbb Z},\qquad
 r=\lfloor k/2\rfloor,\quad d=\lceil k/2\rceil.
\]
Then for every finite field $\mathbb F_q$, including characteristic two,
\[
 \#V_k(\mathbb F_q)=q^{k-1}+(q-1)q^{d-1}.                 \tag{MT.2}
\]
In the relative Grothendieck ring of varieties over $\mathbb Z$ one has
\[
                 [V_k]=\mathbb L^{k-1}+\mathbb L^d-\mathbb L^{d-1}.
                                                               \tag{MT.3}
\]
Here $\mathbb L=[\mathbb A^1_{\mathbb Z}]$; the identity is an identity of
classes, not a claim of a direct-sum splitting of all integral realizations.

To prove these formulas, group the odd coordinates into $x\in\mathbb A^d$
and the even coordinates into $y\in\mathbb A^r$.  The quadratic form
becomes $x^{\mathsf T}Ay$, where the integral $d$ by $r$ matrix $A$ has an
$r$ by $r$ triangular minor with diagonal entries one.  It has a left
inverse over $\mathbb Z$ and its columns extend to a unimodular basis of
$\mathbb Z^d$.  Thus after an invertible integral change of the $x$
coordinates the equation is $\sum_{j=1}^r x_jy_j=0$, with $d-r$ free
$x$ coordinates.  Over $y=0$ its fibre is $\mathbb A^d$.  Over
$\mathbb A^r\setminus\{0\}$, the map of vector bundles
$\mathcal O^d\to\mathcal O$ given by $(y_1,\ldots,y_r,0,\ldots,0)$ is
surjective, and its kernel is a vector bundle of rank $d-1$.
It is Zariski locally trivial on the opens where a coordinate $y_j$ is
invertible.  Additivity and the vector-bundle relation give
\[
 [V_k]=\mathbb L^d+(\mathbb L^r-1)\mathbb L^{d-1},
\]
which is (MT.3).  Counting the same fibres over $\mathbb F_q$ proves
(MT.2) directly, without a purity theorem.

Consequently the exact error term in the unweighted orthogonality identity
of the supplied Fable sketch is
\[
             q\,\#V_k(\mathbb F_q)=q^k+(q-1)q^d.
\]
For a fixed nonzero additive-character parameter its quadratic sum is
$q^d$, by summing first over $x$.  This gives square-root size when $k$ is
even; for odd $k$ the extra factor $q^{1/2}$ reflects the one-dimensional
radical.  It explains the sign and the exact constant for this family.
The general weighted operator estimate in
\texttt{OURS\_FABLE\_PROGRAM\_SUGGESTIONS\_20261120\_EXPLAINED.md},
\S3, remains a separate inherited result.  This calculation does not prove
the sign or norm assertions for arbitrary sheaves in Alp\"oge's steps 3
and 5, nor identify this elementary model with all the constructions in
the omitted Lean dependency tree.

\subsection*{MT-Q02. Its arithmetic zeta function and shifted zero lines}
The exact all-prime formula in MT-Q01 gives
\[
 Z(V_k/\mathbb F_q,T)
     =\frac{1-q^{d-1}T}{(1-q^{k-1}T)(1-q^dT)},
 \qquad
 \zeta_{V_k}(z)
     =\frac{\zeta(z-k+1)\zeta(z-d)}{\zeta(z-d+1)}.       \tag{MT.4}
\]
Indeed apply (MT.2) over every $\mathbb F_{q^m}$ in
$Z=\exp(\sum_{m\ge1}\#V_k(\mathbb F_{q^m})T^m/m)$ and use
$-\log(1-aT)=\sum_{m\ge1}a^mT^m/m$.  Multiplying the resulting factors
over rational primes proves the second identity first in a right
half-plane of absolute convergence, and then meromorphically.
Under $A_D$ the nontrivial divisor of this function is supported on
\[
    \operatorname{Re}z=k-\tfrac12,\qquad
    \operatorname{Re}z=d+\tfrac12,\qquad
    \operatorname{Re}z=d-\tfrac12.
\]
Divisor orders are signed and cancellations are kept exactly as in
(MT.4); the denominator's zeros are possible poles, not additional zeros.
The shifted pole of $\zeta$ at $1$ and shifted trivial zeros are recorded
by the same explicit formula, separately from this nontrivial-divisor
assertion.  This is an all-characteristic, concrete Artin--Tate example
linking the path model of the method with arithmetic zeta functions and
the Grothendieck ring.  It illustrates why a mixed-weight scheme has
several shifted lines rather than a single universal critical line.

\subsection*{MT-L01. Damped Euler products and Adams operations}
Let $\rho$ be a finite-image complex representation of $G_K$ with no
invariant vectors, of dimension $d$.  At each finite place $v$ put
$A_v=\rho(\operatorname{Frob}_v)|_{V^{I_v}}$.  Its eigenvalues have
absolute value one.  For $|c|\le1$ define, initially for $\operatorname{Re}s>1$,
\[
 D_c(s,\rho)=\prod_v\det(1-cA_v(Nv)^{-s})^{-1}.
\]
If $L(s,\rho)$ is holomorphic and nonzero on
$H=\{\operatorname{Re}s>1/2\}$, then $D_c$ has a holomorphic nonvanishing
continuation to $H$, given by
\begin{align*}
 D_c(s,\rho)&=\exp\{c\log L(s,\rho)+E_c(s)\},\\
 E_c(s)&=\sum_v\sum_{k\ge2}
   \frac{c^k-c}{k}\operatorname{tr}(A_v^k)(Nv)^{-ks}.
\end{align*}
The branch of $\log L$ is fixed by the Euler series in
$\operatorname{Re}s>1$.

For a full proof, expand $-\log\det(1-X)$ in powers of $X$ in the initial
half-plane.  The terms with $k=1$ cancel in
$\log D_c-c\log L$, giving the displayed $E_c$.  On a compact subset with
$\operatorname{Re}s\ge1/2+\varepsilon$, its terms in absolute value sum to
at most
\[
 2d\sum_v\sum_{k\ge2}(Nv)^{-k(1/2+\varepsilon)}
 \le \frac{2d}{1-2^{-1/2-\varepsilon}}
       \sum_v(Nv)^{-1-2\varepsilon}<\infty.
\]
Normal convergence, also uniform in $|c|\le1$, proves that $E_c$ is
holomorphic in $s$.  A nonvanishing holomorphic function on the simply
connected half-plane $H$ has a holomorphic logarithm, uniquely normalized
by its existing branch on the right.  The exponential above is therefore
holomorphic, nonzero, and agrees with the product where it converges.
This proves the continuation.

Under $M_H$, Brauer induction supplies the needed holomorphic nonvanishing
on $H$ for $\rho$ with no invariants: all Hecke factors have no zeros
there, and their orders at $1$ combine to
$-\dim V^{G_K}=0$.  This does not require Artin holomorphy on the boundary
line.  Under $A_D$ the same argument applies to sums of nonprincipal
Dirichlet characters over $\mathbb Q$.

The traces $\operatorname{tr}(A_v^k)$ are the local Adams-power characters;
the proof separates the first prime term from all higher Adams terms,
whose series converge in the larger half-plane.  This is a concrete
$\lambda$-ring link, rather than a claim that positivity of coefficients
follows from the formal operations.  For $c=1/m$ the phenomenon is closely
related to Gun--Hazra--Sahu, \emph{On holomorphy and non-vanishing of Artin
L-functions} (2025), Theorem 5.  The displayed formula covers all complex
$|c|\le1$ and makes its hypotheses explicit.

\subsection*{MT-L02. One irrational parameter detects the zero-free half-plane}
Fix an irrational real number $0<c<1$ and a nontrivial irreducible Artin
representation $\rho$.  Then the following are equivalent:
\begin{enumerate}
 \item $L(s,\rho)$ is holomorphic and nonzero on $H$;
 \item the Euler product $D_c(s,\rho)$ admits a single-valued meromorphic
       continuation to $H$.
\end{enumerate}
In the second assertion meromorphic is a global, single-valued condition;
a branch on a slit domain is not enough.

MT-L01 proves the forward implication.  Conversely $L$ is meromorphic by
Brauer induction and $E_c$ is holomorphic on $H$ by the same normally
convergent series without any hypothesis about the zeros of $L$.
Logarithmic differentiation in $\operatorname{Re}s>1$ gives
\[
 \frac{D_c'}{D_c}=c\frac{L'}L+E_c'.
\]
Both sides are meromorphic on connected $H$ if the asserted continuation
of $D_c$ exists, so the identity holds on $H$.  Taking residues at any
$s_0\in H$ gives
\[
 \operatorname{ord}_{s_0}D_c
      =c\operatorname{ord}_{s_0}L.
\]
Both orders are integers, and $c$ is irrational.  Both must be zero.
Thus $L$ is holomorphic and nonvanishing on $H$, proving the equivalence.
If instead $c=a/b$ in lowest terms, the same equality proves
$b\mid\operatorname{ord}_{s_0}L$ everywhere in $H$.

GRH in the meromorphic-divisor sense for the nontrivial irreducible Artin
factors is therefore equivalent, by the functional equation, to these
irrational-parameter continuation tests.  To include the trivial factor one
must separately add RH for $\zeta_K$; that input is already part of $M_H$.
This is an analytic multiplicity test.  It neither constructs the
continuation from arithmetic lattices nor removes poles on the critical
line.  With a trivial summand of multiplicity $r$, the known pole at $1$
contributes local monodromy $e^{-2\pi i c r}$; this is why the no-invariants
hypothesis was retained rather than suppressed.

\subsection*{MT-ART05. The exact induced-linear cone for \texorpdfstring{$A_5$}{A5}}
The coarse counterexample MT-ART04 can be sharpened to a degree-three
virtual character. More usefully, the same calculation gives explicit
lower bounds for any pole of an actual $A_5$ Artin function. These are
group-theoretic deductions, not realizations of the virtual examples as
arithmetic divisor orders.

Write the irreducible characters in the order
$1,\chi_3,\chi'_3,\chi_4,\chi_5$, and put
\[
 \Theta=a+b\chi_3+c\chi'_3+d\chi_4+e\chi_5,
 \qquad a,b,c,d,e\in\mathbb Z.
\]
Then the inequalities
$\langle\Theta,\operatorname{Ind}_H^{A_5}\eta\rangle\ge0$ for
every subgroup $H$ and every linear character $\eta$ are equivalent to
the following seven inequalities:
\[
 a\ge0,\quad e\ge0,\quad a+d\ge0,\quad b+c\ge0,\quad
 b+d+e\ge0,\quad c+d+e\ge0,\quad b+c+d\ge0.       \tag{MT-A5-cone}
\]
If $r=\Theta(1)=a+3b+3c+4d+5e$, then $d<0$ forces $r\ge3$,
and $b<0$ or $c<0$ forces $r\ge4$. These thresholds are sharp for
this abstract cone: respectively
\[
       1+\chi_3+\chi'_3-\chi_4,
       \qquad -\chi_3+\chi'_3+\chi_4               \tag{MT-A5-ex}
\]
have degrees three and four and satisfy every inequality.

\begin{proof}
We give the character and subgroup details so the finite assertion has
no unreported character-table oracle. The five conjugacy classes have
orders $1,2,3,5,5$ and sizes $1,15,20,12,12$. The natural four-dimensional
augmentation character has values $(4,0,1,-1,-1)$ and norm one, as in
MT-ART04. Its real representation is oriented because all permutations
are even. The Hodge-star decomposition of its exterior square has two
invariant real three-dimensional spaces. The exterior-square character
has values $(6,-2,0,1,1)$ by
\[
 \chi_{\wedge^2 V}(g)=\tfrac12(\chi_V(g)^2-\chi_V(g^2)).
\]
Its norm squared is two and it has no invariant vectors, so the two
three-dimensional summands are distinct irreducibles. Simplicity of
$A_5$ makes each faithful and its determinant trivial. An order-two or
order-three element in either corresponding $\mathrm{SO}_3$ representation
has trace $-1$ or $0$. At order five the two traces are
$\varphi=(1+\sqrt5)/2$ and $\bar\varphi=(1-\sqrt5)/2$; their sum is
one and each is the trace of a nontrivial order-five rotation.

There are six Sylow-five subgroups. Their normalizers have order ten
and are dihedral, since the centralizer of a five-cycle is its own
cyclic subgroup. The permutation action on the six subgroups has
values $(6,2,0,1,1)$: an order-five element fixes its unique Sylow
subgroup, an order-three element fixes none, and double counting gives
$6\cdot5/15=2$ fixed subgroups for an involution. Removing the constants
gives an irreducible character $\chi_5=(5,1,-1,0,0)$, of norm one.
The squared dimensions $1+9+9+16+25=60$ exhaust the character table.

Every proper subgroup is contained in an $A_4$, a $D_5$ of order ten,
or an $S_3$ of order six. Here is an elementary verification. Simplicity
excludes proper indices two, three and four. The remaining possible
proper subgroup orders divide sixty and belong to
$\{1,2,3,4,5,6,10,12\}$. There are no elements of order four, six or
ten. Thus subgroups of orders four, six and ten have types
$V_4,S_3,D_5$. The centralizer of an involution is its unique $V_4$,
whose normalizer is $A_4$. In a subgroup of order twelve, a normal
Sylow-three subgroup is impossible since the ambient normalizer of
$C_3$ has order six. There are therefore four Sylow-three subgroups,
contributing eight nonidentity elements. Three Sylow-two subgroups
would contribute nine further involutions, because distinct $V_4$'s
have disjoint nonidentity sets. Hence the Sylow-two subgroup is normal
and the order-twelve subgroup is its $A_4$ normalizer. The smaller
cyclic and Klein subgroups are contained in these listed groups.

Let $u,\bar u$ be the nontrivial linear characters of $A_4$,
$v$ its irreducible character of degree three, and let $\epsilon$
denote a sign character of either $D_5$ or $S_3$. The restriction
decompositions obtained from the displayed character values are
\[
\begin{array}{c|ccc}
 & A_4 & D_5 & S_3\\\hline
1&1&1&1\\
\chi_3&v&\epsilon+v_1&\epsilon+v_2\\
\chi'_3&v&\epsilon+v'_1&\epsilon+v_2\\
\chi_4&1+v&v_1+v'_1&1+\epsilon+v_2\\
\chi_5&u+\bar u+v&1+v_1+v'_1&1+2v_2
\end{array}                                                   \tag{MT-A5-res}
\]
Here $v_1,v'_1$ are the two dihedral characters of degree two and
$v_2$ is the standard $S_3$ character. The nonlinear characters in
these three groups are themselves induced from linear characters:
$v$ from a nontrivial character of $V_4$, $v_1,v'_1$ from $C_5$,
and $v_2$ from $C_3$. Transitivity of induction and Frobenius reciprocity
therefore show that all induced-linear tests imply nonnegative
multiplicities in every column of (MT-A5-res). Together with the
global trivial-character test $a\ge0$, these are exactly the seven
inequalities (MT-A5-cone), after deleting redundant inequalities.
Conversely those seven inequalities make every column an actual
character. Restricting further gives actual characters on all proper
subgroups, hence every required linear inner product is nonnegative.
The global linear test is again $a\ge0$. This proves the equivalence.

For the degree bounds, first put $d=-k<0$. Then $a\ge k$,
$b+c\ge k$, and $b,c\ge k-e$. If $e=0$, these imply
$r\ge k+6k-4k=3k\ge3$. If $e\ge1$, the first two inequalities give
$r\ge k+3k-4k+5e\ge5$. Next put $b=-k<0$. If $d\ge0$, we have
$c\ge k$, $d+e\ge k$, whence
$r\ge4d+5e\ge4k$. If $d=-j<0$, then
$a\ge j$, $c\ge k+j$, and $e\ge k+j$, giving
$r\ge j+3j-4j+5(k+j)\ge4k$. The argument for $c<0$ is identical.
Finally direct substitution verifies both examples (MT-A5-ex).
\end{proof}

\subsection*{MT-ART06. Pole multiplicity gates for actual \texorpdfstring{$A_5$}{A5} extensions}
Let $L/K$ be an $A_5$ extension and $s_0\ne1$. Its Heilbronn character
has coefficients equal to the local Artin orders. Hecke holomorphy and
induction give all the tests in MT-ART05, unconditionally. Hence every
Artin factor is holomorphic at a zero of $\zeta_L$ of order at most two.
A pole of the four-dimensional factor requires order at least three
and a zero of $\zeta_K$ at the same point. A pole of either
three-dimensional factor requires order at least four. Conditional on
$M_H$, all remaining nontrivial candidate poles lie on the critical line.

At the central point the lower bounds strengthen to six and eight,
respectively. To prove this final refinement, the permutation actions
of degrees five and six give
\[
 L(s,\chi_4)=\frac{\zeta_{L^{A_4}}(s)}{\zeta_K(s)},\qquad
 L(s,\chi_5)=\frac{\zeta_{L^{D_5}}(s)}{\zeta_K(s)}.
\]
The completed Dedekind zeta functional equation has sign $+1$, and
its completion factors are finite and nonzero at $1/2$. Consequently
the central orders $a,d,e$ are even integers, with $a,e\ge0$ by
(MT-A5-cone). If $d=-k<0$, then $k\ge2$; the preceding proof gives
$r\ge3k\ge6$ when $e=0$ and $r\ge5e\ge10$ otherwise.
If $b=-k<0$ and $d\ge0$, the nonnegative even integers $d,e$
satisfy $d+e\ge k\ge1$, so $r\ge4d+5e\ge8$.
If $d=-j<0$, then $j\ge2$ and the preceding bound
$r\ge5(k+j)$ is stronger. The same holds for $c<0$.
This proves all assertions; no assumption on simplicity of zeros has
been introduced. The degree-three and degree-four virtual examples
only show the limits of the indicated character tests away from the
additional central parity constraints.

These refinements put the residual Artin question into explicit local
inequalities on a finite Galois quotient, within the SGA 1 V \S7
framework. They sharpen ORDER item 42 and ZETA\_TWO ZT02 without
claiming general Artin holomorphy or constructing actual negative
Artin divisor orders.

\subsection*{MT-H01. A unitary Hecke family outside the algebraic-CM normalization}
Let $K=\mathbb Q(\sqrt5)$, $\epsilon=(1+\sqrt5)/2$, and let $\sigma$
be its nontrivial automorphism. Put $t_0=\pi/\log\epsilon$.
For each nonzero integer $k$ there is a conductor-one unitary Hecke
character defined on fractional ideals by
\[
       \eta_k((\alpha))=
       \left|\frac{\alpha}{\sigma(\alpha)}\right|^{ik t_0}.
                                                               \tag{MT-H1}
\]
It has infinite order, $\eta_k^\sigma=\eta_{-k}$, and is trivial
on the rational ideles. No positive power of $\eta_k$ becomes an
algebraic Hecke character after multiplying by a complex norm power
and a finite-order character. Thus the full unitary-Hecke input $M_H$
has an explicit scope beyond the algebraic Hecke characters used in
the potentially CM motive transfer.

\begin{proof}
The ring of integers is $\mathbb Z[\epsilon]$, with discriminant five.
Minkowski's ideal-class bound is $\sqrt5/2<2$, so every fractional
ideal is principal. An exact classical locator for the bound is
Milne, \emph{Algebraic Number Theory}, Theorem 4.3 (v3.08), PDF p.~72.
The units are $\{\pm\epsilon^j:j\in\mathbb Z\}$. For a direct
verification, a positive unit $u$ with $1<u<\epsilon$ would have
integer trace $u+u^{-1}$ in $(2,\sqrt5)$ if its norm were one, or
$u-u^{-1}$ in $(0,1)$ if its norm were minus one; neither interval
contains an integer. Multiplying an arbitrary positive unit by a
power of $\epsilon$ reduces it to $[1,\epsilon)$, so it is a power
of $\epsilon$, and signs give the assertion.

Multiplying $\alpha$ by a unit changes the logarithm of the absolute
ratio in (MT-H1) by $2j\log\epsilon$. Its phase therefore changes by
$2\pi i k j$, so (MT-H1) is well-defined and multiplicative. It is
unitary. The classical ideal/infinity-type description of Hecke
characters makes it an unramified finite-place character, with
archimedean exponents $\pm ikt_0$ (their simultaneous sign depends
on the ideal-to-idele convention). The unit compatibility just checked
is exactly what is required for this description. Rational ideles
have equal archimedean components and equal valuations at conjugate
places; their images are trivial. For a ramified rational place this
also follows here from the generator $\sqrt5$, whose absolute ratio
is one. Conjugating (MT-H1) inverts the ratio.

On the positive archimedean subgroup the exponents of any algebraic
Hecke character are integers, up to an overall convention sign.
A norm power adds the same complex number to the two exponents;
a finite-order character is trivial on this connected subgroup.
If a positive power $m$ of (MT-H1) had the asserted algebraic
normalization, subtraction of its two exponents would give
$2mikt_0$ equal to the difference of two integers. This is impossible.
The same nonconstant archimedean restriction proves infinite order.
\end{proof}

This does not assert a transcendence theorem for individual Hecke
values, or rule out every conceivable generalized motive. It precisely
separates the character class covered by CM01 from the broader unitary
character class quantified in Masullo's endpoint. A proposed transport
of the CM theta-value and algebraic-height steps must therefore address
this extra archimedean data; norm twisting alone cannot do so.

\subsection*{MT-H02. An explicit nonholomorphic \texorpdfstring{$\mathrm{GL}_2$}{GL2} family}
For $k\ge1$, the classical Hecke--Maa\ss\ correspondence associates to
$\eta_k$ a cuspidal dihedral Maa\ss\ newform $g_k$ of weight zero,
level five and nebentypus $\chi_5$. Its spectral parameter is $kt_0$,
its Laplace eigenvalue is $1/4+k^2t_0^2$, and
\[
       L(s,g_k)=L_K(s,\eta_k).
                                                               \tag{MT-H2}
\]
Consequently $M_H$ gives GRH for every member of this family. Its
entireness is already the classical Hecke theorem. The exact primary
locator for the correspondence, level, spectral parameter and coefficient
identity is Humphries--Khan, \emph{On the Random Wave Conjecture for
Dihedral Maa\ss\ Forms}, GAFA 30 (2020), Appendix A.1, printed pp.~106--107 (PDF pp.~73--74), the displayed
formula for $\psi((\alpha))$ and the subsequent formulas for
$t_g$ and $\lambda_{g_\psi}(n)$.

For clarity, the transfer proof is short and complete. MT-H01 gives
a non-norm character: its two archimedean exponents are unequal.
The cited correspondence is therefore cuspidal, and its formula
$\lambda_{g_k}(n)=\sum_{N\mathfrak a=n}\eta_k(\mathfrak a)$ identifies
the absolutely convergent Dirichlet series with (MT-H2), including all
finite primes. Analytic continuation gives the equality everywhere.
The character is unitary, so its nontrivial critical line is
$\operatorname{Re}s=1/2$, which is exactly the assertion supplied by
$M_H$. This proves the stated implication without a generic
non-dihedral Maa\ss\-form hypothesis.

These forms live on the analytic modular curve of level five and in
the automorphic-induction part of the adelic framework. They are an
explicit additional $\mathrm{GL}_2$ branch beside the holomorphic
CM and finite-image weight-one branches, not the missing general
non-CM $\mathrm{GL}_2$ theorem in ORDER item 42.

\subsection*{MT-H03. Exact tensor and symmetric-power products in that family}
Write $F_j(s)=L_K(s,\eta_j)$ for $j\ne0$ and
$F_0(s)=\zeta_K(s)=\zeta(s)L(s,\chi_5)$. Then $F_{-j}=F_j$.
For the actual local two-dimensional parameters $V_j$ of $g_j$, the
following finite Euler-product identities hold:
\[
 L(s,V_m\otimes V_n)=F_{m+n}(s)F_{m-n}(s),\qquad m,n\ge1,
                                                               \tag{MT-H3}
\]
and, for $n\ge0$ and $k\ge1$,
\[
 L(s,\operatorname{Sym}^n V_k)=
 \prod_{j=0}^{\lfloor(n-1)/2\rfloor}F_{(n-2j)k}(s)
 \begin{cases}\zeta(s),&n\text{ even},\\1,&n\text{ odd}.
 \end{cases}                                                   \tag{MT-H4}
\]
The product is empty when $n=0$. Conditional on $M_H$, all nontrivial
zeros of these products lie on $\operatorname{Re}s=1/2$. They are
entire except for a simple pole at one when $m=n$ in (MT-H3), or
when $n$ is even in (MT-H4). This is a covered arbitrarily large
degree family, not a use of an unknown general symmetric-power lift.

\begin{proof}
Conjugating ideals proves $F_{-j}=F_j$. At a split rational prime
the two local parameters of $V_k$ are $z^k,z^{-k}$, where
$z=\eta_1(\mathfrak p)$. The tensor weights are
$z^{m+n},z^{m-n},z^{-m+n},z^{-m-n}$, which gives (MT-H3).
The symmetric-power weights are $z^{(n-2j)k}$ for $0\le j\le n$;
pairing opposite weights leaves one middle weight $1$ precisely for
even $n$, proving (MT-H4) there.

At an inert prime the unique prime ideal is $(p)$ and has character
value one. The two induced Frobenius eigenvalues are $1,-1$ for
every nonzero $k$. Tensoring gives two of each sign. On the
symmetric power, the even-index monomials give eigenvalue $1$ and
the odd-index monomials give $-1$. If $n=2r$, these multiplicities
are $r+1,r$; if $n=2r+1$, they are $r+1,r+1$. These are exactly
the factors on the right sides. If $m=n$, the factor $F_0$ at an
inert prime is $(1-p^{-2s})^{-1}$, so the same check includes that
case.

The only ramified prime is five. The ideal $(\sqrt5)$ has
$\eta_k$-value one and the character is unramified at every finite
place. Thus its local character at that quadratic completion is
trivial, and its induced Weil parameter is $1\oplus\chi_{5,5}$.
The inertia-fixed dimension is one. A tensor product has two
trivial constituents and two $\chi_{5,5}$ constituents; a symmetric
power has respectively $r+1,r$ or $r+1,r+1$ of them. The latter
constituents have no inertia invariants. Hence the ramified Euler
factors also agree exactly, rather than only away from five.
Equivalently the proof is the weight-pair decomposition for diagonal
matrices $\operatorname{diag}(z,z^{-1})$ and an outside coset
acting by interchanging the two basis vectors.

Absolute convergence initially gives the identities for
$\operatorname{Re}s>1$. The right sides give their meromorphic
continuations. All $F_j$, $j\ne0$, are entire Hecke functions with
nonzero boundary value at one; the latter is Ramakrishnan--Valenza,
Theorem 7-28, printed p.~289. The same holds for $L(1,\chi_5)$.
Thus the displayed zeta factors give exactly the asserted simple
poles, with no cancellation at one. Finally $M_H$ locates all the
nontrivial zeros of each unitary Hecke factor. The identities contain
no negative exponents, so they introduce no hidden poles on the
critical line. This proves every assertion.
\end{proof}

The middle summand in (MT-H4) is the trivial character for this
specific family. In a differently normalized dihedral family it
depends on the central character; replacing it universally by either
$1$ or $\chi_K$ would be incorrect. This is the representation-ring
and lambda-operation part of the adelic framework, beside SGA 6 V
\S\S2,7, with the explicit analytic parameters of MT-H01 retained.

\subsection*{MT-ART07. The entire non-real $A_5$ Artin category}
Let $L/\mathbb Q$ be Galois with group $A_5$, and suppose $L$ is not
totally real. Every nontrivial irreducible Artin L-function for this
extension is entire, unconditionally. Under $M_H$ each satisfies GRH.
This closes the non-real $A_5$ specialization of the holomorphy gate;
MT-ART05/06 remain useful for totally real extensions and other base fields.

The holomorphy assertion is classical, not a new finding. An exact primary
locator is Calegari, \emph{The Artin conjecture for some $S_5$-extensions},
\S1, publisher-author PDF p.~2, the paragraph immediately before
Theorem 1.1: it states the stronger automorphy assertion for $A_5$ when
complex conjugation is nontrivial, citing Khare--Wintenberger and Kim's
icosahedral functorial constructions. The two conditions are equivalent
here because a Galois extension has a real embedding exactly when every
embedding is real, equivalently complex conjugation acts trivially.
We cite that classical result rather than repeat its proof.

Here is the complete GRH transfer and its tensor-category form.
For every finite-dimensional representation $V$ of $A_5$, put
$m=\dim V^{A_5}$. Complete reducibility gives
\[
 L(s,V)=\zeta(s)^m
          \prod_{\chi\ne1}L(s,\chi)^{\langle V,\chi\rangle},
 \qquad \langle V,\chi\rangle\in\mathbb Z_{\ge0}.
\]
The product after removing $\zeta^m$ is entire by the cited theorem.
For each irreducible factor, Brauer induction expresses its L-function
as a finite product of integer powers of finite-order Hecke L-functions.
At a point of the open critical strip off its middle line, $M_H$ makes
every such factor finite and nonzero. The product is consequently nonzero
there, including when some Brauer exponents are negative. This proves
the claimed GRH without making a positivity assumption on those exponents.

At $s=1$ every nontrivial irreducible factor is also nonzero. Indeed in a
Brauer expression $\chi=\sum_j a_j\operatorname{Ind}_{H_j}^{A_5}\eta_j$,
the order there is $-\sum_{\eta_j=1}a_j$, by the simple Dedekind pole and
Hecke boundary nonvanishing. Frobenius reciprocity identifies the sum with
$\langle\chi,1\rangle=0$. Thus $L(s,V)$ has a pole of order exactly $m$
at $1$, and has no other poles. Tensor products, duals, exterior and
symmetric powers, and every Schur functor produce representations of the
same finite group, so the same conclusions hold for them after removing
their explicitly determined invariant summands. This is a statement about
the whole Artin tensor category attached to this cover, not about a
general positive-weight motive. The companions in Grothendieck's language
are SGA~1, Expos\'e~V, \S7 (finite monodromy), and SGA~6, Expos\'e~V,
\S\S2,7 (representation-ring operations), together with AG-LAMBDA.

\subsection*{MT-ART08. A finite exceptional set is enough for a holomorphy transfer}
Here is a useful precise comparison lemma. Let $V$ be a self-dual,
finite-image complex Galois representation over a number field $K$.
Suppose an Euler product $F(s)$ converges in a right half-plane and
continues to an entire function, with the same finite local factors as $L(s,V)$ outside a
finite set $S$. Suppose also that, for each $v\in S$, the reciprocal of
the local factor of $F$ is a polynomial $Q_v((Nv)^{-s})$.
Then $L(s,V)$ is entire. No equality of ramified local parameters is
needed for this conclusion.

To prove this, write the Artin factor at $v$ as
$P_v((Nv)^{-s})^{-1}$. Every eigenvalue of Frobenius on the inertia
invariants is a root of unity, since $V$ has finite image. Hence the zeros
of $P_v((Nv)^{-s})$ all lie on $\operatorname{Re}s=0$.
Equality in a right half-plane, followed by meromorphic continuation,
gives
\[
 L(s,V)=F(s)\prod_{v\in S}
             \frac{Q_v((Nv)^{-s})}{P_v((Nv)^{-s})}.
\]
The right side is holomorphic on $\operatorname{Re}s>0$.
The classical completed Artin functional equation is
\[
 A^{s/2}\Gamma_V(s)L(s,V)
   =W A^{(1-s)/2}\Gamma_V(1-s)L(1-s,V),
\]
with $A>0$ and $W\ne0$; the gamma factor is a product of
$\Gamma_{\mathbb R}(s)$, $\Gamma_{\mathbb R}(s+1)$ and
$\Gamma_{\mathbb C}(s)$ with nonnegative integral multiplicities.
An exact locator for these formulas is Prasad--Yogananda,
\emph{A Report on Artin's holomorphy conjecture}, \S3.3,
PDF pp.~8--9, with Brauer induction in Theorem 3.3.
For $\operatorname{Re}s<1$, the numerator
$\Gamma_V(1-s)L(1-s,V)$ is holomorphic, since
$\operatorname{Re}(1-s)>0$. The reciprocal $1/\Gamma_V(s)$ is entire.
The functional equation therefore makes $L(s,V)$ holomorphic on this
second half-plane too. The two half-planes cover $\mathbb C$, proving
the assertion. For a non-self-dual $V$ the identical argument works if
the hypotheses are also supplied for $V^\vee$.

The finiteness of $S$ is an actual hypothesis. Agreement at a set of
places of density one does not give the finite product in the proof.
In particular this lemma cannot by itself upgrade the weak
five-dimensional correspondence in Calegari's Theorem 1.1.
In geometric terms this separates modification at finitely many closed
points of an arithmetic base from a density statement about infinitely
many Frobenius fibres.

\subsection*{MT-S501 and MT-S502. The $S_5$ transfer and its one-value gate}
Let $L/\mathbb Q$ be an $S_5$ extension, with complex conjugation of cycle
type $(12)(34)$. Assume $5$ is unramified and its Frobenius also has this
cycle type. Then the irreducible factors of dimensions $4$ and $6$ are
entire and, under $M_H$, satisfy GRH (MT-S501).

Let $E=L^{A_5}$, let $F\subset L$ have degree $6$ over $\mathbb Q$, and
put $H=EF$. Under $M_H$, the additional single condition
\[
                         \zeta_H(1/2)\ne0
\]
implies that every nontrivial irreducible factor is entire and satisfies
GRH (MT-S502). This is a proved reduction with the displayed extra input,
not a proof of that central nonvanishing.

The classical inputs are Calegari's Theorems 1.1 and 1.2,
publisher-author PDF pp.~2--3; the definition $H=EF$ is explicit in
Theorem 1.2 of arXiv:1112.1152v1, p.~2. The former supplies the
dimension-$4$ and dimension-$6$ holomorphy, and the latter supplies the
remaining case if $\zeta_H$ has no zero in the real interval $(0,1)$.
Under $M_H$ the only possible real zero in that interval is $1/2$.
Thus its interval condition is exactly equivalent to the displayed
single-value condition. Once holomorphy is available, the Brauer
zero-location argument of MT-ART07 proves GRH for every factor.
This is the full proof of both transfers. The one-dimensional factors
are Dirichlet functions and already fall under $A_D$.

The totally real quadratic field $E$ is the base for the odd projective
two-dimensional representations used in the cited theorem. Accordingly
this is a higher-rank consequence of a specific automorphic descent
theorem. It is not obtained by treating every $S_5$ factor as a single
Hecke character. In the Artin-motive programme, the remaining gate in
this family is one central Dedekind value. AG-SIGN controls a nonzero
leading coefficient after a central zero is removed; it does not prove
that the central order is zero.

\subsection*{MT-S503. One unconditional holomorphy example becomes a full GRH example}
Take the splitting field $L$ of
\[
                        x^5-x^3-x^2+x+1.
\]
Calegari, Theorem 5.1, publisher-author PDF p.~15, proves automorphy of
every representation of $\operatorname{Gal}(L/\mathbb Q)=S_5$;
the quintic field has discriminant $1609$. Consequently under $M_H$
every nontrivial irreducible Artin L-function of this explicit field
satisfies GRH, by the same signed Brauer argument. This example uses
that cited theorem, not the local-at-$5$ hypotheses of MT-S501, which
the source explicitly says this polynomial does not satisfy.
The historical bounded-height computation in Theorem 5.3 is part of
the classical proof's evidence; it is not a computation performed in
this run. Geometrically this supplies a concrete finite \(S_5\)-cover
of the arithmetic base whose complete Artin category is covered.

\subsection*{MT-M05. The stated Schauder quotient loses more than its zero evaluations see}
Lemma 9.3.2.5 on original PDF pp.~425--426 asserts a basis for the
quotient by the actual convolution ideal $\varphi*\mathcal S_E(G)$,
with coefficient functionals the Mellin jets at the zeros of
$\mathcal M\varphi$. We inspected the original p.~426 image: there
is no closure over this ideal. The assertion is false even when the
Mellin divisor consists of one simple zero. Thus the counterexample
does not rely on the zero-count defect in MT-M01.

Use $G=\mathbb R_{>0}$, trivial character and $u=\log t$; convolution
becomes ordinary additive convolution. Write
\[
 g_a(u)=(4\pi a)^{-1/2}e^{-u^2/(4a)},\qquad a>0.
\]
Choose $0<b<a$ and set $\varphi(e^u)=-g_a'(u)$,
$h(e^u)=-g_b'(u)$. Both functions belong to the exponential-Schwartz
space, and
\[
 \mathcal M\varphi(s)=s e^{as^2},\qquad
 \mathcal Mh(s)=s e^{bs^2}.
\]
The first function has just the simple zero $0$, lying in the interior
of the strip for $I=[-1,1]$. The only claimed coefficient functional is
$\delta_0$, and $\delta_0(h)=0$.
But $h\notin\varphi*\mathcal S_E$: if $h=\varphi*f$, multiplicativity
of the absolutely convergent Mellin transform gives
$\mathcal Mf(s)=e^{(b-a)s^2}$ for $s\ne0$, and then everywhere by
analyticity. On the imaginary axis this has modulus
$e^{(a-b)\tau^2}$. In contrast, $f(e^u)$ is integrable and therefore
$|\mathcal Mf(i\tau)|\le\int|f(e^u)|\,du$, a contradiction.
Thus the nonzero quotient class of $h$ has every asserted coefficient
zero, which is incompatible with uniqueness in a Schauder expansion.
The still simpler choice $\varphi=g_a$, $h=g_b$ has no Mellin zeros
and a nonzero quotient, so even the empty-divisor case fails.

\subsection*{MT-M06. The closed-ideal repair for every polynomial Gaussian divisor}
In logarithmic coordinates put
\[
 \mathcal E=\{f\in C^\infty(\mathbb R,\mathbb C):
   p_{N,j}(f):=\sup_x e^{N|x|}|f^{(j)}(x)|<\infty
   \text{ for every integer }N,j\ge0\}.
\]
This is $\mathcal S_E(\mathbb R_{>0})$ with its exponential-Schwartz
Fr\'echet topology. Polynomial factors in the usual Schwartz seminorms
are absorbed by increasing $N$. Let
\[
 P(s)=\prod_{\nu=1}^r(s-\rho_\nu)^{m_\nu},\quad
 d=\sum_\nu m_\nu,\qquad
 \varphi=P(-D)g_a,\quad D=d/dx.
\]
Distinct $\rho_\nu\in\mathbb C$, positive integral $m_\nu$, and
$a>0$ are arbitrary; $P=1,d=0$ is allowed. Then
\[
 \overline{\varphi*\mathcal E}
   =J_P:=\{h\in\mathcal E:
          (\mathcal Mh)^{(j)}(\rho_\nu)=0,
          \ 0\le j<m_\nu\},                         \tag{MT-M6}
\]
and the quotient by this \emph{closed} ideal is topologically
$\mathbb C^d$, with exactly the displayed Mellin jets as coordinates.
In particular its coordinate basis is absolute. This repairs the
finite-divisor case after making the closure operation explicit.

\begin{proof}
First, $g_a*\mathcal E$ is dense in $\mathcal E$. We give the full
argument, since density is distinct from the false surjectivity in
MT-M05. Convolution is continuous on $\mathcal E$: the inequality
$e^{N|x|}\le e^{N|x-y|}e^{N|y|}$ bounds each differentiated
convolution by an exponential supremum seminorm times an exponential
$L^1$ seminorm; the latter is bounded by a constant times
$p_{N+1,0}$.

Suppose a continuous linear functional $T$ vanishes on
$g_a*\mathcal E$. For every real $y$ and $\epsilon>0$ the function
$g_{a+\epsilon}(x-y)=g_a*g_\epsilon(\,\cdot-y)$ lies in this range.
It converges to $g_a(x-y)$ in every $p_{N,j}$ as
$\epsilon\downarrow0$, by uniform Gaussian estimates with the
parameter in a compact subinterval of $(0,\infty)$. Hence
$T(g_a(\,\cdot-y))=0$. The function
\[
 B(z)=T\bigl(e^{-x^2/(4a)}e^{zx/(2a)}\bigr)
\]
is entire: on every compact set of $z$, all its proposed complex
derivatives converge in every $p_{N,j}$ by Gaussian domination.
Completing the square gives $B(y)=0$ for real $y$, so the identity
theorem gives $B=0$.

For any Schwartz function $\psi$, Fourier inversion now gives
\[
 T(e^{-x^2/(4a)}\psi(x))
  =\frac1{2\pi}\int_{\mathbb R}\widehat\psi(\xi)
        B(2ai\xi)\,d\xi=0.
\]
This passage under $T$ is justified in the Fr\'echet space itself:
the $p_{N,j}$ seminorm of
$e^{-x^2/(4a)}e^{i\xi x}$ is at most
$C_{N,j,a}(1+|\xi|)^j$, whereas $\widehat\psi$ is rapidly
decreasing. For $h\in C_c^\infty$, take
$\psi=e^{x^2/(4a)}h$, which is again compactly supported and smooth.
Then $T(h)=0$. Finally $C_c^\infty$ is dense in $\mathcal E$:
multiply by a smooth cutoff equal to one on $[-R,R]$ and supported
in $[-2R,2R]$; the product rule and the bound from $p_{N+1,j}$ make
every $p_{N,j}$ of the discarded tail tend to zero. Thus $T=0$.
Hahn--Banach separation proves the asserted density.

Second, the range of $P(-D)$ on $\mathcal E$ is exactly $J_P$.
Integration by parts gives $\mathcal M(-Df)(s)=s\mathcal Mf(s)$,
so inclusion in $J_P$ follows. For the converse, if
$\mathcal Mh(\rho)=0$, define
\[
 f(x)=e^{-\rho x}\int_{-\infty}^x e^{\rho t}h(t)\,dt
     =-e^{-\rho x}\int_x^\infty e^{\rho t}h(t)\,dt.
\]
The equality is precisely the vanishing of the Mellin value.
The first expression controls the negative tail and the second the
positive tail. For any prescribed $N$, using an exponential bound
on $h$ with exponent larger than $N+|\operatorname{Re}\rho|+1$
proves exponential decay of $f$ on both tails; the equation
$f'=h-\rho f$ then gives the same bounds for every derivative.
Therefore $f\in\mathcal E$, and $(D+\rho)f=h$.
Equivalently $-f$ solves $(-D-\rho)(-f)=h$, and its Mellin transform
is $\mathcal Mh(s)/(s-\rho)$. Repeating this division through all
roots, including multiplicities, proves $P(-D)\mathcal E=J_P$.

The Mellin jet functionals are continuous, so $J_P$ is closed.
Since $g_a*\mathcal E$ is dense and $P(-D)$ is continuous,
$P(-D)(g_a*\mathcal E)$ is dense in $P(-D)\mathcal E=J_P$.
Differentiation commutes with convolution, identifying that range
with $\varphi*\mathcal E$. This proves (MT-M6).

For $d=0$ the quotient is zero and no interpolation is required.
For $d>0$, it remains to verify that no coordinate is missing. Given arbitrary
$d$ target jets, ordinary polynomial Hermite interpolation constructs
a polynomial $Q$ of degree below $d$ for which
$Q(s)e^{bs^2}$ has exactly those jets, with any fixed $b>0$.
Indeed multiplication by the nowhere-zero entire function $e^{bs^2}$
is invertible on each finite jet space, so interpolate the resulting
jets of $Q$. The function $Q(-D)g_b$ belongs to $\mathcal E$ and
realizes them. Choosing the targets to be the coordinate vectors
gives a continuous linear section from the finite-dimensional space
$\mathbb C^d$ to $\mathcal E$. The jet map consequently induces
a continuous isomorphism of the quotient with $\mathbb C^d$,
with continuous inverse. Finite-dimensional coordinate expansions are
absolute, completing the proof.
\end{proof}

For $P=s$, MT-M05 therefore exhibits a proper, nonclosed convolution
ideal whose closure is $\ker\delta_0$. The closed quotient really
does have the one coordinate the draft intended. For $P=1$ the
principal ideal is dense but proper, and the closed quotient is zero.
This is a concrete repair rather than a certification of the general
infinite-divisor statement. For arithmetic Hecke divisors the remaining
gate is controlled infinite interpolation and the claimed absolute
basis/nuclear estimates on the correct separated quotient. A local
$O(\log T)$ zero count alone does not prove those estimates.
In cote~1 trace language, Hausdorff completion and the identification
of the spectral quotient precede taking its nuclear trace.

\subsection*{MT-M07. An infinite closed quotient with an absolute Mellin basis}
The closed-ideal repair also admits a complete infinite-divisor example.
Fix $a,\ell>0$ in the logarithmic space $\mathcal E$ of MT-M06, and put
\[
 \varphi(x)=g_a(x-\ell)-g_a(x),\qquad
 \rho_n=\frac{2\pi i n}{\ell}\quad(n\in\mathbb Z).
\]
Then $\mathcal M\varphi(s)=e^{as^2}(e^{\ell s}-1)$, with precisely
the simple zeros $\rho_n$. Let
\[
 (\mathcal P_\ell f)(t)=\sum_{k\in\mathbb Z}f(t+k\ell).
\]
The quotient $\mathcal E/\overline{\varphi*\mathcal E}$ is
topologically $C^\infty(\mathbb R/\ell\mathbb Z)$ under
$\mathcal P_\ell$, and has an absolute Schauder basis whose coefficient
functionals are exactly $f\mapsto\mathcal Mf(\rho_n)$.

\begin{proof}
The periodization series and every derivative converge uniformly on a
fundamental interval, by exponential decay. This gives a continuous map
to the indicated smooth periodic space. Choose
$\chi\in C_c^\infty(\mathbb R)$ nonnegative and positive on
$[-\ell,\ell]$, and set
$\theta=\chi/\mathcal P_\ell\chi$. The denominator is positive,
smooth and periodic, so $\theta$ is smooth, compactly supported and
$\mathcal P_\ell\theta=1$. The map
$R(F)(x)=\theta(x)F(x)$ is a continuous section of periodization,
proving surjectivity.

Let $\Delta_\ell h(x)=h(x-\ell)-h(x)$. Telescoping shows that its
range lies in $\ker\mathcal P_\ell$. Conversely, if
$\mathcal P_\ell f=0$, define
\[
 h(x)=\sum_{k\ge1}f(x+k\ell)
     =-\sum_{k\le0}f(x+k\ell).
\]
These locally normally convergent series define a smooth function.
For $x\ge0$ use the first expression and for $x\le0$ the second.
For every positive integer $N$ and every $j$, summing the geometric
series of bounds $p_{N,j}(f)e^{-N|x+k\ell|}$ proves
$p_{N,j}(h)<\infty$; the $N=0$ case follows from $N=1$.
Thus $h\in\mathcal E$, and direct index shifting gives
$\Delta_\ell h=f$. The range of $\Delta_\ell$ is therefore the
closed space $\ker\mathcal P_\ell$.

By MT-M06, $g_a*\mathcal E$ is dense in $\mathcal E$.
Continuity of $\Delta_\ell$ and commutation with convolution imply
\[
 \overline{\varphi*\mathcal E}
 =\overline{\Delta_\ell(g_a*\mathcal E)}
 =\ker\mathcal P_\ell.
\]
The continuous section $R$ identifies the quotient topology with
the usual smooth periodic topology, rather than merely giving a
vector-space identification.

The Fourier expansion in the convention used here is
\[
 \mathcal P_\ell f(t)=\frac1\ell
       \sum_{n\in\mathbb Z}\mathcal Mf(\rho_n)e^{-\rho_n t}.
\]
Its coefficients are obtained by integrating over a fundamental
interval and summing translates, justified by absolute convergence.
They decrease faster than every power of $|n|$: integrate by parts
arbitrarily many times and use the integrability of every derivative
of $f$. The Fourier series consequently converges absolutely in
every smooth seminorm and equals $\mathcal P_\ell f$ by Fourier
uniqueness. Representatives for the basis are
\[
                     b_n(x)=\ell^{-1}\theta(x)e^{-\rho_n x}.
\]
Indeed $\mathcal Mb_n(\rho_m)=\delta_{m,n}$, by
$\mathcal P_\ell\theta=1$ and integration on one period. Moreover,
$p_{N,j}(b_n)\le C_{N,j}(1+|n|)^j$ since $\theta$ has compact
support. Hence the expansion
$\sum_n\mathcal Mf(\rho_n)b_n$ is absolutely convergent even in
$\mathcal E$, and represents the class of $f$ in the quotient.
The jets determine that class uniquely by the Fourier identification.
Ordering the indices as $0,1,-1,2,-2,\ldots$ gives an absolute
Schauder basis, as asserted.
\end{proof}

The actual ideal still cannot replace its closure: with $0<b<a$,
$h=\Delta_\ell g_b$ has all these Mellin values zero but is not
$\varphi*f$ for any $f\in\mathcal E$, by exactly the Fourier-growth
division argument of MT-M05. Thus this example tests an infinite
divisor with controlled spacing, not just the empty or finite case.

\subsection*{MT-M08. The corresponding trace and positive convolution formula}
For $f\in\mathcal E$, convolution by $f$ descends to the quotient of
MT-M07. On $L^2(\mathbb R/\ell\mathbb Z,dt)$ it is trace class and
\[
 \operatorname{Tr}C_f
   =\sum_{n\in\mathbb Z}\mathcal Mf(\rho_n)
   =\ell\sum_{k\in\mathbb Z}f(k\ell).                 \tag{MT-M8}
\]
All sums converge absolutely. If
$f=h*\widetilde h$ with $h\in\mathcal E$ and
$\widetilde h(x)=\overline{h(-x)}$, this operator is positive and
\[
                  \operatorname{Tr}C_f
                   =\sum_n|\mathcal Mh(\rho_n)|^2\ge0.
\]

To prove the assertions, first use absolute convergence and changes of
variables to periodize a convolution:
\[
 \mathcal P_\ell(f*g)(t)
  =\int_0^\ell(\mathcal P_\ell f)(t-u)
                     (\mathcal P_\ell g)(u)\,du.
\]
This proves descent and identifies the smooth kernel
$K_f(t,u)=(\mathcal P_\ell f)(t-u)$.
The orthonormal Fourier vectors
$\ell^{-1/2}e^{-\rho_n t}$ are eigenvectors with eigenvalues
$\mathcal Mf(\rho_n)$. These eigenvalues decrease faster than every
power, as proved in MT-M07, so their absolute sum is finite.
The diagonal expansion is thus a trace-class expansion and gives
the first trace sum. Evaluating the absolutely convergent Fourier
series of $\mathcal P_\ell f$ at zero gives the second equality
in (MT-M8), including its factor $\ell$. Finally the Mellin transform
of $\widetilde h$ at a purely imaginary argument is the conjugate
of that of $h$; the eigenvalues are therefore the displayed squares.

There is also a direct nuclear realization on the smooth quotient.
Explicitly, put $c_n(F)=\ell^{-1}\int_0^\ell F(t)e^{\rho_nt}\,dt$
and $e_n(t)=e^{-\rho_nt}$. The $c_n$ are an equicontinuous family,
bounded by the uniform norm. In
\[
 C_f=\sum_n(1+|n|)^{-2}\,
       c_n\otimes\bigl((1+|n|)^2\mathcal Mf(\rho_n)e_n\bigr),
\]
the scalars are summable and the vector family is bounded in every
$C^\infty$ seminorm. This is a nuclear decomposition. Fourier partial
sums converge uniformly on bounded sets of $C^\infty$ (apply a
uniform higher-derivative bound to the coefficient tail), so the
finite-rank trace limit is the same absolute sum above.

This supplies a fully checked spectral-quotient, basis and trace
model beside cote~1's nuclear trace language. It does not identify
an arithmetic prime distribution with this lattice or prove the
corresponding basis theorem for the actual Hecke zero set. In that
case the exact infinite interpolation and operator estimates remain
to be supplied on the correct closed quotient.

\subsection*{MT-M09. Even a bounded local zero count does not give the asserted point basis}
Passing to the closed ideal is necessary, but does not prove the
general basis claim. There exists $\varphi\in\mathcal E$ whose Mellin
transform has only simple zeros on the imaginary axis and only
$O(1)$ zeros in every unit-height interval, for which
\[
        Q=\mathcal E/\overline{\varphi*\mathcal E}
\]
has no Schauder basis with the individual zero evaluations as its
coefficient functionals. This holds for every enumeration of the
zeros, and therefore also rules out an absolute Schauder basis of
that specified kind.

\begin{proof}
Put $\epsilon_n=e^{-n^4}$ for $n\ge1$ and define
\[
 H(s)=\prod_{n=1}^{\infty}
       \left(1+\frac{s^2}{n^2}\right)
       \left(1+\frac{s^2}{(n+\epsilon_n)^2}\right),
 \qquad F(s)=e^{as^2}H(s),\quad a>0.
\]
The product converges uniformly on compact sets because the sum of
the absolute deviations of its factors from one does. Away from the
listed zeros the tail logarithms converge locally uniformly, so
there are no additional zeros. Its zeros are exactly
$\pm in,\ \pm i(n+\epsilon_n)$, all simple. Since
$0<\epsilon_n<1/2$, each unit-height interval contains a bounded
number of them.

For $r=|s|$ the product satisfies
\[
 \log|H(s)|\le2\sum_{n\ge1}\log(1+r^2/n^2)\le2\pi r.
\]
For the last bound compare the decreasing function of $x$ with its
integral on $(0,\infty)$; integration by parts gives
$\int_0^\infty\log(1+r^2/x^2)\,dx=\pi r$.
Consequently, on every strip $|\sigma|\le A$,
\[
 |F(\sigma+it)|
 \le \exp\{aA^2-at^2+2\pi\sqrt{A^2+t^2}\}.
\]
Define
\[
 \varphi(x)=\frac1{2\pi}\int_{\mathbb R}F(it)e^{-itx}\,dt.
\]
The Gaussian bound permits every derivative in $x$ under the integral.
Shifting the contour to $\operatorname{Re}s=\sigma$ gives, for
every real $\sigma$ and every $j\ge0$,
\[
 \varphi^{(j)}(x)
   =\frac{e^{-\sigma x}}{2\pi}
      \int_{\mathbb R}(-\sigma-it)^jF(\sigma+it)e^{-itx}\,dt.
\]
The horizontal sides of the contour tend to zero by the displayed
Gaussian estimate, uniformly for $\sigma$ in the finite interval
used in that shift. Choosing $\sigma=N+1$ for $x\ge0$ and
$\sigma=-N-1$ for $x\le0$ proves every seminorm $p_{N,j}$ is finite.
Thus $\varphi\in\mathcal E$. Cauchy estimates applied on slightly
larger strips show that $F(it)$ and all its derivatives are Schwartz.
Fourier inversion gives $\mathcal M\varphi(it)=F(it)$;
both sides are entire, so $\mathcal M\varphi=F$ everywhere.

Every evaluation at a listed zero annihilates the convolution ideal
and its closure, and therefore descends continuously to $Q$.
Let $p=p_{1,0}$ and let
$\overline p([f])=\inf_{j\in\overline{\varphi*\mathcal E}}p(f+j)$
be the corresponding quotient seminorm. For any $f\in\mathcal E$,
\[
 |(\delta_{in}-\delta_{i(n+\epsilon_n)})(f)|
  \le\epsilon_n\int_{\mathbb R}|x||f(x)|\,dx
  \le2\epsilon_n p(f).
\]
The first inequality follows from
$|e^{i\epsilon_n x}-1|\le\epsilon_n|x|$.
Because the difference functional annihilates the closed ideal,
the same bound holds with $p$ replaced by $\overline p$ on $Q$.

Suppose the claimed Schauder basis existed, and denote the basis
vector associated to $\delta_{in}$ by $b_n$. Biorthogonality of
basis vectors and coefficient functionals gives
$\delta_{in}(b_n)=1$ and
$\delta_{i(n+\epsilon_n)}(b_n)=0$. Hence
\[
                 \overline p(b_n)\ge(2\epsilon_n)^{-1}.
\]
Take the single test function $g_b\in\mathcal E$ for any fixed
$b>0$. Its coefficient at $\delta_{in}$ is $e^{-bn^2}$.
The corresponding terms of its proposed basis expansion satisfy
\[
 \overline p\bigl(e^{-bn^2}b_n\bigr)
        \ge\frac12 e^{n^4-bn^2}\longrightarrow\infty.
\]
Terms of a convergent series in a locally convex space must tend
to zero. In every enumeration these infinitely many distinct terms
occur arbitrarily far along the series, giving a contradiction.
This proves the assertion.
\end{proof}

This is a counterexample for general exponential-Schwartz Mellin
divisors, not a claim that such clustered zeros occur for an actual
Hecke L-function. It pinpoints why NT-Z10/MT-M04 cannot by themselves
establish the arithmetic Schauder step: even an $O(1)$ local count
does not control the norm of individual interpolation vectors.
The next possible repair must incorporate quantitative separation,
or group nearby evaluations into divided-difference blocks and prove
convergence of the resulting block expansion. Neither property is
supplied by a zero-location hypothesis alone in the present argument.

\subsection*{MT-M10. A point-jet basis forces a polynomial cluster-product bound}
The counterexample MT-M09 has a quantitative converse obstruction.
Let $\mathcal E$ be the exponential-Schwartz space of MT-M06,
let $J\subset\mathcal E$ be a closed linear subspace, and put
$Q=\mathcal E/J$. Fix a real $\sigma$ and a discrete set
$Z\subset\sigma+i\mathbb R$, with positive integer multiplicities
$m_\rho$. Suppose every functional
\[
 \delta_{\rho,k}(f)=\frac{(\mathcal Mf)^{(k)}(\rho)}{k!},
 \qquad 0\le k<m_\rho,
\]
annihilates $J$ and occurs as a coefficient functional of a
Schauder basis of $Q$ with continuous coordinate functionals.
Other basis coordinates are permitted.
Then there are constants $c>0$ and $A\ge0$, independent of
$\rho=\sigma+it$ and of the finite subset
$W\subset Z\setminus\{\rho\}$, such that
\[
       \prod_{\nu\in W}|t-\operatorname{Im}\nu|^{m_\nu}
                    \ge c(1+|t|)^{-A}.                \tag{MT-M10}
\]
In particular the distinct ordinates have a uniform polynomial
separation bound. The conclusion already follows from an ordinary
Schauder basis; absolute convergence is not needed.

\begin{proof}
The seminorms $p_{N,j}(f)=\sup_x e^{N|x|}|f^{(j)}(x)|$ make
$\mathcal E$ a Fr\'echet space: a Cauchy sequence converges uniformly,
with all derivatives, on every compact interval, and the weighted
uniform limits give the same smooth limit and all the seminorm
limits. The closed quotient $Q$ is Fr\'echet as well.
For a continuous seminorm $p$ on $\mathcal E$, write
\[
 \overline p([f])=\inf_{u\in J}p(f+u).
\]
The quotient seminorms of finite maxima of the $p_{N,j}$ generate
the quotient topology.

First fix any such $\overline p$. Let $S_n$ be the partial-sum
operators for the proposed basis, and let $R_n=S_n-S_{n-1}$ be
its rank-one coordinate projections, with $S_0=0$.
Convergence of every basis expansion makes
$\sup_n\overline p(R_nq)$ finite for each $q\in Q$.
For completeness, the uniform-boundedness step follows directly
from Baire's theorem. The closed sets
\[
 A_k=\{q:\sup_n\overline p(R_nq)\le k\},\qquad k=1,2,\ldots,
\]
cover $Q$. One contains $q_0+U$, where $U$ is a neighborhood of
zero, and $q_0$ itself may be chosen in that set. Subtraction
gives $\overline p(R_nu)\le2k$ for $u\in U$, uniformly in $n$.
Choose $U$ to contain a unit ball for a finite maximum of
quotient seminorms and rescale. This proves that for a continuous
quotient seminorm $\overline q$ and a constant $C$,
\[
             \overline p(R_nq)\le C\overline q(q)
                         \quad\hbox{for every }n,q.    \tag{1}
\]
The same conclusion holds when the seminorm on the right is
zero, by applying the bound to arbitrary scalar multiples.

Let $b_{\rho,0}$ be the basis vector with coefficient
$\delta_{\rho,0}$. For a fixed Gaussian $g_a$ with
$\mathcal Mg_a(s)=e^{as^2}$, use the test function
\[
                f_t(x)=e^{-\sigma x-itx}g_a(x).
\]
It satisfies $\delta_{\rho,0}(f_t)=1$.
Every fixed seminorm of $f_t$ is bounded by a fixed power of
$1+|t|$, by the Leibniz formula and Gaussian decay.
Apply (1) to $[f_t]$ and its coordinate projection. For each
fixed $\overline p$ there are constants $C_p,A_p$ such that
\[
       \overline p(b_{\rho,0})\le C_p(1+|t|)^{A_p}
                                      \quad(\rho\in Z). \tag{2}
\]
The constants do not depend on the enumeration position of
$b_{\rho,0}$.

We next give a uniform estimate for confluent divided differences.
For real nodes $t_0,\ldots,t_m$, repeated nodes allowed, the
divided difference of $u\mapsto e^{ixu}$ is
\[
 [t_0,\ldots,t_m]e^{ixu}
   =\int_{\Delta_m}(ix)^m
           e^{ix\sum_{j=0}^m\lambda_jt_j}\,d\lambda,
 \qquad
 \Delta_m=\{\lambda_j\ge0:\textstyle\sum\lambda_j=1\},
                                                        \tag{3}
\]
where $d\lambda$ has total mass $1/m!$.
One elementary verification suffices also for repeated nodes.
For distinct nodes, the divided difference of $u^k$ is zero
when $k<m$ and is the complete homogeneous polynomial
$h_{k-m}(t_0,\ldots,t_m)$ when $k\ge m$. This follows by
expanding the partial-fraction identity
\[
 \sum_{j=0}^m
 \frac{1}{(1-zt_j)\prod_{\ell\ne j}(t_j-t_\ell)}
          =\frac{z^m}{\prod_{j=0}^m(1-zt_j)}.
\]
On the simplex,
$\int_{\Delta_m}\lambda_0^{a_0}\cdots\lambda_m^{a_m}\,d\lambda
=\prod_j a_j!/(m+\sum a_j)!$; successive integration of
monomials proves this formula. Expanding the exponential and
using this integral gives exactly the same power series,
proving (3). Continuity as nodes coalesce gives the confluent
version, whose functionals are derivatives divided by factorials.
In particular the absolute value in (3) is at most $|x|^m/m!$.

Take $t_0=t$ and repeat each $\operatorname{Im}\nu$, $\nu\in W$,
exactly $m_\nu$ times among the other nodes. Apply the divided
difference to $u\mapsto\mathcal Mf(\sigma+iu)$.
All resulting point-jet functionals annihilate $J$.
Choose an integer $N\ge|\sigma|+1$ and put $p=p_{N,0}$.
Formula (3), Fubini's theorem and exponential decay give
\[
 \left|[t_0,\ldots,t_m]\mathcal Mf(\sigma+iu)\right|
 \le\int_{\mathbb R}\frac{|x|^m}{m!}e^{\sigma x}|f(x)|\,dx
 \le2p(f).                                             \tag{4}
\]
The last constant is independent of both the nodes and $m$,
since $\int_{\mathbb R}|x|^m e^{-|x|}\,dx=2m!$.
The same inequality holds on $Q$ with $\overline p$.

Biorthogonality says that on $b_{\rho,0}$ all the other
point jets in this divided difference vanish. The coefficient
of the simple node $t$ is
$\prod_{\nu\in W}(t-\operatorname{Im}\nu)^{-m_\nu}$.
For repeated other nodes this follows either from partial
fractions or by coalescing distinct nodes away from $t$.
Consequently (4) and (2) imply
\[
 \prod_{\nu\in W}|t-\operatorname{Im}\nu|^{-m_\nu}
       \le2\overline p(b_{\rho,0})
       \le2C_p(1+|t|)^{A_p},
\]
which is (MT-M10).
\end{proof}

For any fixed $0<r<1$, choosing all other zeros within distance
$r$ also bounds their total multiplicity by
$(A\log(1+|t|)-\log c)/|\log r|$.
The product condition is stronger than such a local count:
it controls the joint cost of arbitrarily close clusters.
Taking one other node recovers polynomial separation.
Taking the clusters in MT-M09 contradicts (MT-M10) directly.
This identifies an explicit necessary analytic input for
Masullo 9.3.2.5, PDF~425--426, even after the closed-ideal repair.
It neither asserts that arithmetic zeros violate that condition
nor proves that the condition suffices for a basis.

\subsection*{MT-M11. The same basis would bound arithmetic leading coefficients}
Fix a primitive unitary Hecke $L$-function $L(s,\eta)$ and assume
its nontrivial zeros lie on $\operatorname{Re}s=1/2$.
Suppose a closed quotient as in MT-M10 has all its nontrivial
zero jets as coefficient functionals, with their actual
multiplicities. Then there are constants $c,C,T_0>0$, depending
on the fixed function and the proposed basis, such that at every
nontrivial zero $\rho=1/2+iT$ with $|T|\ge T_0$ and order $m$,
\[
               \left|\frac{L^{(m)}(\rho,\eta)}{m!}\right|
                        \ge c(1+|T|)^{-C}.             \tag{MT-M11}
\]
In particular, at simple zeros the reciprocal of $|L'(\rho,\eta)|$
would have at most polynomial growth. The zero-location input is
$A_D$ for Dirichlet functions and $M_H$ for general Hecke functions;
the asserted quotient basis is an additional hypothesis.

\begin{proof}
We use the fixed-function analytic facts proved in NT-Z10:
polynomial growth on a fixed vertical strip away from its
finitely located poles, and $O(\log(2+|T|))$ zeros, with
multiplicity, in each bounded-height window.
The absolutely convergent Euler product gives
$|L(2+iT,\eta)|\ge c_\eta>0$: each local inverse factor
has modulus at least $(1+(\operatorname N\mathfrak p)^{-2})^{-1}$,
and the product of these lower bounds is positive.

Put $c_T=2+iT$. Choose $4\le R\le5$ so that the boundary of
$D=\{s:|s-c_T|<R\}$ contains no zero. At sufficiently large
$|T|$, it contains no pole or trivial zero either, because
the imaginary parts of their finitely many arithmetic
progressions parallel to the real axis are fixed by the infinity type.
Its zeros are therefore precisely the relevant critical-line
zeros in that disk, say $\rho_j$ with multiplicities $m_j$.
There are
$M=\sum_jm_j=O_\eta(\log(2+|T|))$ of them.
In the normalized coordinate $z=(s-c_T)/R$, put
$a_j=(\rho_j-c_T)/R$ and form the finite Blaschke product
\[
 B(s)=\prod_j
       \left(\frac{z-a_j}{1-\overline{a_j}z}\right)^{m_j}.
\]
Constant unimodular normalizations of the factors are
irrelevant. The function $G=L/B$ is holomorphic and zero-free
on the disk and its boundary. On the boundary $|B|=1$.
The polynomial growth bound and the maximum principle give
$|G(s)|\le M_T=C_1(1+|T|)^{C_2}$ throughout $D$, enlarging
$M_T$ if necessary to make it at least one.
At the center, $|B(c_T)|\le1$, so
$|G(c_T)|\ge c_\eta$.

The nonnegative harmonic function
$u(s)=\log M_T-\log|G(s)|$ satisfies the disk Harnack bound
\[
 u(\rho)\le
 \frac{R+|\rho-c_T|}{R-|\rho-c_T|}\,u(c_T)
 \le\frac{11}{5}\,u(c_T),
\]
because $|\rho-c_T|=3/2$ and $R\ge4$.
This elementary bound follows by comparing the maximum and
minimum of the Poisson kernel on each concentric boundary
circle and then letting that circle increase to radius $R$.
It implies $|G(\rho)|\ge c_2(1+|T|)^{-C_3}$.

Let $a=(\rho-c_T)/R$, so $|a|\le3/8$.
The leading coefficient of its own factor in $B$ contributes
$[R(1-|a|^2)]^{-m}\ge5^{-m}$.
Every other factor at $\rho$ has modulus
\[
 \frac{|a-a_j|}{|1-\overline{a_j}a|}
 \ge\frac{|\rho-\rho_j|}{10}.
\]
Thus, with $W$ the distinct other zeros in $D$,
\[
 \left|\lim_{s\to\rho}\frac{B(s)}{(s-\rho)^m}\right|
 \ge 10^{-M}\prod_{\rho_j\in W}|\rho-\rho_j|^{m_j}.
\]
All these points are on the same vertical line. Apply
(MT-M10) to this finite set $W$. Since
$10^{-M}\ge c_3(1+|T|)^{-C_4}$, the displayed expression
has a polynomial lower bound. Multiplying by the lower bound
for $|G(\rho)|$ proves (MT-M11).
\end{proof}

This is a necessary consequence of the proposed arithmetic
point-jet basis, not a derivative estimate obtained from GRH
alone. It converts the remaining interpolation problem into
a concrete additional test on arithmetic leading coefficients.
The work placement is the Masullo 9.3.2.5 repair, NT-Z10,
and the cote~1 trace/Fredholm programme reported in SA23
geo-H-lidskii/geo-H-genus. Positivity of the trace cannot
replace this topological prerequisite.

\subsection*{MT-M12. An exact simple-zero gate for the canonical arithmetic quotient}
Fix a primitive unitary Hecke character $\eta$ over a number field
$K$, and assume GRH for its $L$-function. Assume additionally that
all its nontrivial zeros are simple. Let
\[
 \Xi_\eta(s)=P_\eta(s)\Lambda(s,\eta)
\]
be the entire completed function with its possible poles removed.
For a non-pure-norm character take $P_\eta=1$. If
$L(s,\eta)=\zeta_K(s+i\tau)$, take
$P_\eta(s)=(s+i\tau)(s-1+i\tau)$.
The zeros of $\Xi_\eta$ are precisely the nontrivial zeros, denoted
$Z$. Put
\[
 \psi_\eta(x)=\frac1{2\pi}\int_{\mathbb R}\Xi_\eta(it)e^{-itx}\,dt,
 \qquad Q_\eta=\mathcal E/\overline{\psi_\eta*\mathcal E}.
\]
Then $\psi_\eta\in\mathcal E$, with
$\mathcal M\psi_\eta=\Xi_\eta$, and the following are equivalent:
\begin{enumerate}
\item There are $c,C>0$ such that
$|L'(\rho,\eta)|\ge c(1+|\operatorname{Im}\rho|)^{-C}$ for
every $\rho\in Z$.
\item $Q_\eta$ has an absolute Schauder basis with coefficient
functionals $[f]\mapsto\mathcal Mf(\rho)$, $\rho\in Z$.
\item $Q_\eta$ has a Schauder basis with these continuous
coefficient functionals.
\end{enumerate}
Under these conditions the closed convolution ideal is exactly
\[
 K_Z=\{f\in\mathcal E:\mathcal Mf(\rho)=0
                              \text{ for every }\rho\in Z\}. \tag{MT-M12a}
\]
Moreover $Q_\eta$ is topologically the rapidly decreasing sequence
space
\[
 s(Z)=\{(a_\rho):\sup_{\rho\in Z}
       (1+|\operatorname{Im}\rho|)^k|a_\rho|<\infty
                                  \text{ for every }k\ge0\}. \tag{MT-M12b}
\]
This equivalence uses $A_D$ for Dirichlet functions, or $M_H$
conditionally for general unitary Hecke functions. Simplicity is
an extra assumption throughout; condition (1) is the remaining
gate, not a consequence of those GRH inputs proved here.

\begin{proof}
\emph{Analytic preliminaries.}
Write $n=[K:\mathbb Q]$ and $\kappa=\pi n/4$.
The fixed-function polynomial strip bounds of NT-Z10 and
Stirling's formula for the $n$ factors $\Gamma_{\mathbb R}$
give, for every fixed $A>0$, constants $C_A,B_A$ such that
\[
 |\Xi_\eta(\sigma+it)|
       \le C_A(1+|t|)^{B_A}e^{-\kappa|t|}
                            \qquad(|\sigma|\le A).     \tag{5}
\]
On bounded $t$ the entire function is bounded on each compact
strip rectangle, so enlarging $C_A$ covers possible gamma-pole
cancellations there. Derivative bounds follow by Cauchy's formula
on a slightly larger strip. Thus $\Xi_\eta(it)$ is Schwartz.
The contour-shifting and Fourier-inversion argument of MT-M09,
using (5) in place of the Gaussian bound, shows that the displayed
inverse transform belongs to $\mathcal E$ and has Mellin transform
$\Xi_\eta$. More explicitly, shift to $\sigma=N+1$ on $x\ge0$
and to $\sigma=-N-1$ on $x\le0$; every integral with the extra
factor $(-\sigma-it)^j$ is absolutely convergent by (5).
This proves all $p_{N,j}$ bounds. The horizontal contour edges
tend to zero by the same estimate, and Fourier uniqueness then
proves the transform identity.

Assume condition (1). At a zero $\rho=1/2+iT$, differentiation
of the completed function gives
\[
 \Xi_\eta'(\rho)
  =P_\eta(\rho)q_\eta^{\rho/2}
       \prod_j\Gamma_{\mathbb R}(\rho+\mu_j)L'(\rho,\eta),
\]
where $q_\eta=|D_K|\operatorname N\mathfrak f_\eta$
is the finite conductor.
Stirling's formula and (1) imply, after enlarging constants to
cover the finitely many low zeros,
\[
        |\Xi_\eta'(\rho)|
           \ge c_1(1+|T|)^{-C_1}e^{-\kappa|T|}.         \tag{6}
\]

\emph{Polynomially bounded cardinal functions.}
Fix $a>0$. For each zero $\rho$ define the entire function
\[
 H_\rho(s)=
     \frac{\Xi_\eta(s)}
          {\Xi_\eta'(\rho)(s-\rho)}\,e^{a(s-\rho)^2}.    \tag{7}
\]
It takes the value one at $\rho$ and zero at every other zero.
For each fixed strip $|\operatorname{Re}s|\le A$, (5)--(6)
imply a bound
\[
 |H_\rho(\sigma+it)|
       \le C_A'(1+|T|)^{B_A'}
                              e^{-a(t-T)^2/2}.         \tag{8}
\]
Here are the details near the apparent denominator. If
$|s-\rho|\ge1$, use (5) and (6), together with
$\big||T|-|t|\big|\le|T-t|$ and
$1+|t|\le(1+|T|)(1+|t-T|)$.
The remaining fixed polynomial and
$e^{\kappa|t-T|}$ are absorbed by half the Gaussian.
If $|s-\rho|<1$, use
\[
 \frac{\Xi_\eta(s)}{s-\rho}
       =\int_0^1\Xi_\eta'(\rho+u(s-\rho))\,du
\]
and the Cauchy derivative bound from (5) on a slightly larger
strip. Since $|t-T|<1$, the same estimate (8) follows by
adjusting the constants.

Invert $H_\rho$ by Fourier transform and call the inverse
$h_\rho$. The same contour argument gives $h_\rho\in\mathcal E$,
$\mathcal Mh_\rho=H_\rho$, and, for each $N,j$,
\[
             p_{N,j}(h_\rho)
                    \le C_{N,j}(1+|T|)^{B_{N,j}}.      \tag{9}
\]
Indeed, on the shifted contour used for $x\ge0$ or $x\le0$,
integrating $(1+|t|)^j$ against the Gaussian in (8) costs
only another fixed power of $1+|T|$.

For any $f\in\mathcal E$, its values $\mathcal Mf(1/2+iT)$
decrease faster than every power of $|T|$. This follows by
integrating by parts arbitrarily many times in
$\int e^{x/2}f(x)e^{iTx}\,dx$; each resulting $L^1$ norm
is bounded by finitely many defining seminorms of $f$.
The local count NT-Z10 implies
$\sum_{\rho\in Z}(1+|\operatorname{Im}\rho|)^{-2}<\infty$.
Consequently
\[
              I(f)=\sum_{\rho\in Z}\mathcal Mf(\rho)h_\rho
                                                               \tag{10}
\]
converges absolutely in every seminorm of $\mathcal E$,
defines a continuous linear map, and has exactly the same
values as $f$ at all the zeros. Likewise, for every
$(a_\rho)\in s(Z)$ the series $\sum_\rho a_\rho h_\rho$
converges in $\mathcal E$ and depends continuously on that
sequence. These assertions use (9) and choose the coefficient
decay power separately for each target seminorm.

\emph{Division on vertical strips.}
We now show the ideal equality, which does not follow merely
from constructing the cardinal functions.
Condition (1) implies a uniform, fixed-strip minimum-modulus
estimate
\[
 |L(\sigma+it,\eta)|
   \ge c_A(1+|t|)^{-C_A}
            \min\{1,\operatorname{dist}(\sigma+it,Z)\},
 \qquad |\sigma|\le A,\quad |t|\ge T_A.                 \tag{11}
\]
Here is a complete proof. Put $r=A+2$ and center a disk at
$c_t=2+it$. Choose $R\in[4r+4,4r+5]$ with zero-free boundary.
At large $|t|$ there are no poles or trivial zeros in that disk.
Its $M=O_A(\log(2+|t|))$ zeros are all simple and on the
critical line. Remove them by the finite Blaschke product
$B$ for this disk, as in MT-M11, and put $G=L/B$.
Polynomial strip growth gives $|G|\le C(1+|t|)^C$ on the
whole disk, while the Euler lower bound at $c_t$ and Harnack's
inequality give
$|G|\ge c(1+|t|)^{-C}$ on the disk of radius $R/2$.
The Harnack factor on that smaller disk is at most three.
The point $s=\sigma+it$ is within radius $r<R/4$.

If the disk has no zeros, this already proves (11).
Otherwise choose a zero $\rho_0$ in it nearest to $s$.
If $|\rho_0-c_t|>R/2$, every zero in the disk is at
distance at least $R/4$ from $s$. Each Blaschke factor
therefore has modulus at least $1/8$, and
$|B(s)|\ge8^{-M}$, a polynomial lower bound.
If $|\rho_0-c_t|\le R/2$, simplicity gives
$L'(\rho_0)=G(\rho_0)B'(\rho_0)$.
In normalized coordinates $|a_0|\le1/2$, the derivative
of its own Blaschke factor is at most $4/(3R)$, and
every other factor at $\rho_0$ is at most
$2|\rho_0-\rho_j|/R$. Thus
\[
 |B'(\rho_0)|\le C_A^M
                     \prod_{\rho_j\ne\rho_0}|\rho_0-\rho_j|.
\]
The assumed lower bound for $L'$, the upper bound for $G$,
and $M=O_A(\log(2+|t|))$ imply
\[
          \prod_{\rho_j\ne\rho_0}|\rho_0-\rho_j|
                          \ge c(1+|t|)^{-C}.           \tag{12}
\]
Since $\rho_0$ is nearest to $s$,
$|\rho_0-\rho_j|\le2|s-\rho_j|$. Each denominator in a
Blaschke factor at $s$ is at most two. Hence
\[
 |B(s)|\ge(4R)^{-M}|s-\rho_0|
                       \prod_{\rho_j\ne\rho_0}|\rho_0-\rho_j|,
\]
and (12) proves (11). Zeros outside the large disk have
distance greater than one from $s$, so using the distance
to the full set $Z$ makes no difference in the minimum.
The estimate at a zero is interpreted as zero on both
sides and follows by continuity.

\emph{The closed ideal is the full zero kernel.}
The inclusion
$\overline{\psi_\eta*\mathcal E}\subset K_Z$ is immediate
from multiplication of Mellin transforms.
Take $f\in K_Z$ and write $F=\mathcal Mf$.
The quotient $F/\Xi_\eta$ is entire because all zeros
are simple and have been cancelled. On every fixed strip,
(11) and the rapid vertical decay of $F$ and $F'$ imply
rapid decay of $F/L$ for large $|t|$. If the distance to
$Z$ is at least one, divide the rapid bound for $F$
by (11). If it is less than one, choose a nearest zero
$\rho$ and use
$F(s)=\int_\rho^s F'(w)\,dw$ on the straight segment;
that segment lies in a slightly larger fixed strip and
a bounded-height window. The distance factor cancels
the corresponding factor in (11).

Inverting the gamma factors now costs at most
$C_A(1+|t|)^{C_A}e^{\kappa|t|}$ on the strip.
For any $\epsilon>0$ the entire function
\[
             J_\epsilon(s)=e^{\epsilon s^2}
                                  \frac{F(s)}{\Xi_\eta(s)}
\]
therefore decays faster than every power on every vertical
strip, in fact with a Gaussian bound. On compact rectangles
it is bounded by entireness. Contour shifting and Fourier
inversion give $j_\epsilon\in\mathcal E$ with this Mellin
transform. Uniqueness of the transform gives
\[
                  \psi_\eta*j_\epsilon=g_\epsilon*f.
\]
Finally $g_\epsilon*f\to f$ in $\mathcal E$ as
$\epsilon\downarrow0$. To verify this in each $p_{N,j}$,
write the difference as the integral of
$g_\epsilon(y)(f^{(j)}(x-y)-f^{(j)}(x))$.
For $|y|\le1$, the weighted mean-value bound is
$|y|e^{N|y|}p_{N,j+1}(f)$; outside any fixed
neighborhood of zero, use
$(e^{N|y|}+1)p_{N,j}(f)$ and the Gaussian tail.
Both integrated bounds tend to zero.
Thus $f$ belongs to the closed convolution ideal,
proving (MT-M12a).

The coordinate map $Q_\eta\to s(Z)$ is now injective.
It is continuous by the integration-by-parts bounds above,
which descend to quotient seminorms because all coordinates
annihilate the ideal. The continuous series map
$(a_\rho)\mapsto[\sum_\rho a_\rho h_\rho]$ is its inverse.
This proves (MT-M12b), and the usual coordinate vectors
give (2), with representatives $h_\rho$.
Statement (2) implies (3). Statement (3) implies (1)
by MT-M11 with $m=1$, extending its high-zero estimate
over the finitely many remaining zeros. This completes
the equivalence.
\end{proof}

Replacing $\psi_\eta$ by
$g_b*\psi_\eta$ for any $b>0$ leaves the closed ideal unchanged:
Gaussian convolution has dense range by MT-M06 and convolution
by $\psi_\eta$ is continuous.

The theorem settles a precise conditional interpolation and
spectral-synthesis interface, rather than identifying every
quotient occurring in Masullo's draft. For this canonical
completed-Hecke quotient with simple zeros, the Schauder and
absolute-Schauder existence gates are exactly the polynomial
reciprocal-derivative gate. The extension to arbitrary
multiplicities requires bounds for the whole local Hermite
interpolation data and is not asserted here.

\subsection*{MT-M13. Nuclear arithmetic convolution under the exact basis gate}
Under the equivalent conditions of MT-M12, convolution by
$f\in\mathcal E$ acts on $Q_\eta\simeq s(Z)$ as
\[
            (a_\rho)\longmapsto
                         (\mathcal Mf(\rho)a_\rho).
\]
It has an explicit nuclear representation on that Fr\'echet
space and a trace-class extension to $\ell^2(Z)$, with
\[
                   \operatorname{Tr}C_f
                         =\sum_{\rho\in Z}\mathcal Mf(\rho).
                                                               \tag{MT-M13}
\]
Here $\operatorname{Tr}$ denotes the trace of that Hilbert-space extension.
The sum converges absolutely. For
$h^\sharp(x)=e^{-x}\overline{h(-x)}$ and
$f=h*h^\sharp$, this Hilbert-space operator is positive and
\[
                   \operatorname{Tr}C_f
                         =\sum_{\rho\in Z}|\mathcal Mh(\rho)|^2
                         \ge0.
\]

\begin{proof}
Convolution descends because it preserves the closed
convolution ideal. Multiplication of Mellin transforms
gives the diagonal action. The eigenvalues decrease
faster than every power of the ordinate, so their
absolute sum is finite by NT-Z10. A diagonal operator
with these eigenvalues on the orthonormal coordinate
basis of $\ell^2(Z)$ is trace class and has the stated
trace.

To see nuclearity without an unspecified basis estimate,
write $b_\rho=[h_\rho]$ and
$\delta_\rho([g])=\mathcal Mg(\rho)$. The functionals
$\delta_\rho$ are equicontinuous: for a fixed
$N>1/2$, their absolute values are uniformly bounded
by a constant times the quotient seminorm of $p_{N,0}$.
The representation
\[
 C_f=\sum_{\rho\in Z}(1+|\operatorname{Im}\rho|)^{-2}
       \delta_\rho\otimes
       \bigl((1+|\operatorname{Im}\rho|)^2
                                  \mathcal Mf(\rho)b_\rho\bigr)
\]
has summable scalar coefficients, an equicontinuous
family of functionals, and a bounded family of vectors
in every quotient seminorm by (9) and rapid decay.
It is therefore an explicit nuclear representation.
Its diagonal sum is the same absolutely convergent
trace displayed above; the unambiguous Hilbert
trace needs no additional abstract trace convention.

The weighted involution preserves $\mathcal E$ and satisfies
\[
      \mathcal M h^\sharp(s)
                =\overline{\mathcal Mh(1-\overline s)}.
\]
At every $\rho$ on the critical line,
$1-\overline\rho=\rho$. Thus the diagonal eigenvalues
for $h*h^\sharp$ are the nonnegative squares stated,
which proves positivity and its trace formula.
\end{proof}

This is a rigorous operator model conditional on the exact
simple-zero basis gate. Its Hilbert completion is constructed
from the zero coordinates; it is not an independently
constructed geometric or prime-side operator proving GRH.
Its location in the programme is Masullo 9.3.2.5--9.3.3,
the Tate-adjoint involution, and cote~1's nuclear trace
framework as carried in SA23 geo-H-lidskii/geo-H-genus.

\subsection*{MT-M14. Bounded multiplicities and the full Hermite basis}
Retain the canonical completion and quotient of MT-M12, assume
GRH, and replace simplicity by the hypothesis
\[
                  1\le m_\rho\le m_0<\infty
                              \qquad(\rho\in Z),       \tag{13}
\]
where $m_\rho$ is the actual zero multiplicity.
Then the following are equivalent:
\begin{enumerate}
\item For some $c,C>0$,
$|L^{(m_\rho)}(\rho,\eta)/m_\rho!|
       \ge c(1+|\operatorname{Im}\rho|)^{-C}$ for all $\rho$.
\item The canonical quotient has an absolute Schauder basis
whose coefficient functionals are
$[f]\mapsto(\mathcal Mf)^{(k)}(\rho)/k!$ for all
$0\le k<m_\rho$.
\item It has an ordinary Schauder basis with precisely those
continuous coefficient functionals.
\end{enumerate}
Under these conditions the closed ideal is the full zero-jet
kernel, and the quotient is the rapid sequence space on
$\{(\rho,k):0\le k<m_\rho\}$, weighted by
$1+|\operatorname{Im}\rho|$. Assumption (13) and condition (1)
are additional arithmetic hypotheses, not consequences of GRH
established in this run.

\begin{proof}
The necessity of (1) from (3) is exactly MT-M11.
We give all changes needed for sufficiency; they also show why
the fixed bound $m_0$ is used.
Write $z=s-\rho$, $m=m_\rho$, and
\[
                  G_\rho(s)=\Xi_\eta(s)/(s-\rho)^m .
\]
The leading-coefficient hypothesis, multiplied by the gamma
and conductor factors, gives
\[
 |G_\rho(\rho)|\ge c(1+|T|)^{-C}e^{-\kappa|T|},
           \qquad T=\operatorname{Im}\rho.
\]
Cauchy's formula on a fixed larger disk around $\rho$ and
the strip bound (5) give
$|G_\rho^{(j)}(\rho)/j!|
 \le C(1+|T|)^C e^{-\kappa|T|}$ for $0\le j\le m_0$.
To see the division estimate directly, on a circle $|z|=2$
divide $\Xi_\eta(\rho+z)$ by $z^m$ and apply Cauchy's
formula inside; $m\le m_0$ keeps all constants fixed.

Fix $a>0$. Let $P_{\rho,k}(z)$ be the Taylor polynomial
through degree $m-1-k$ of the analytic germ
$e^{-az^2}/G_\rho(\rho+z)$.
The reciprocal-coefficient recursion starts with
$1/G_\rho(\rho)$ and at each subsequent degree divides
a finite sum by $G_\rho(\rho)$. The displayed upper and
lower bounds therefore bound every coefficient of
$P_{\rho,k}$ by
$C(1+|T|)^C e^{\kappa|T|}$, uniformly for
$0\le k<m\le m_0$. Only finitely many recursion steps
are involved. Define
\[
 H_{\rho,k}(s)=
      z^kP_{\rho,k}(z)e^{az^2}G_\rho(s).                \tag{14}
\]
Its Taylor expansion at $\rho$ is
$z^k+O(z^m)$, and it vanishes to the required multiplicity
at every other zero. On every fixed vertical strip its
bound has the same Gaussian form (8), with constants
depending also on $m_0$: away from $|z|<1$, use (5)
and the fixed polynomial degree in (14); inside,
use the Cauchy bound for $G_\rho$.
Inverse Mellin transforms $h_{\rho,k}\in\mathcal E$
therefore satisfy polynomial seminorm bounds (9)
uniformly over $k<m_\rho$.

The coefficients $(\mathcal Mf)^{(k)}(\rho)/k!$ are
uniformly rapidly decreasing in $|T|$ for
$k<m_0$. Indeed, they are Mellin transforms of the
fixed finite list $x^k f(x)/k!$, and the
integration-by-parts argument of MT-M12 applies to
each. The series with coefficients these jets and
vectors $h_{\rho,k}$ consequently converges absolutely
in $\mathcal E$, is continuous in $f$, and interpolates
all its jets. Conversely every rapid sequence of jets
is obtained from the corresponding convergent series.
Multiplicity is included in the local zero count,
so the same summability estimate applies.

It remains to verify the full ideal equality.
Repeat the disk argument proving (11), now with
Blaschke factors raised to the actual multiplicities.
The number of factors counted with multiplicity is still
$M=O_A(\log(2+|t|))$. If a nearest zero $\rho_0$
is in the inner half disk, the leading coefficient
identity is
\[
 \frac{L^{(m_{\rho_0})}(\rho_0)}{m_{\rho_0}!}
   =G(\rho_0)
       \lim_{s\to\rho_0}
                    \frac{B(s)}{(s-\rho_0)^{m_{\rho_0}}}.
\]
The upper bound for that leading Blaschke coefficient
is $C_A^M$ times
$\prod_{\rho_j\ne\rho_0}
          |\rho_0-\rho_j|^{m_{\rho_j}}$.
Condition (1) therefore supplies the polynomial lower
bound (12) for this multiplicity-weighted product.
Nearest-zero comparison gives
\[
 |L(s,\eta)|\ge c_A(1+|t|)^{-C_A}
                         |s-\rho_0|^{m_{\rho_0}}
\]
when the nearest distance is less than one.
When that distance is at least one, or the nearest
disk zero is outside the inner half disk, the same
argument gives a polynomial lower bound without
the distance factor. All constants are uniform
because $m_{\rho_0}\le m_0$ and $M=O_A(\log|t|)$.

If $F=\mathcal Mf$ vanishes to the required orders,
Taylor's formula at a nearest zero gives
\[
 F(s)=\frac{(s-\rho)^m}{(m-1)!}
       \int_0^1(1-u)^{m-1}
                         F^{(m)}(\rho+u(s-\rho))\,du .
\]
For $m\le m_0$ its derivative bound is uniformly
rapid in $|t|$. The distance factors cancel.
Thus $F/L$ is rapidly decreasing on each fixed
strip at large height. Exactly the Gaussian-regularized
division and approximate-identity proof in MT-M12
shows that $f$ lies in
$\overline{\psi_\eta*\mathcal E}$.
The reverse inclusion is immediate.
The interpolation series now identifies the quotient
topologically with the rapid jet-sequence space,
as in MT-M12. This proves (2), hence (3), and
completes the equivalence.
\end{proof}

No unspecified infinite Hermite estimate has been hidden in
the argument: the reciprocal-coefficient recursion has at most
$m_0$ steps. With unbounded multiplicities, the exponents of
the polynomial bounds produced by that recursion need not be
uniform. That case remains a separate interpolation gate.

\subsection*{MT-M15. Multiple zeros: positive traces and nilpotent operators}
Under the conditions of MT-M14, convolution has an explicit
nuclear representation on the rapid jet-sequence quotient
and a trace-class realization on the standard $\ell^2$
space of those jets, with
\[
          \operatorname{Tr}C_f
                    =\sum_{\rho\in Z}m_\rho\mathcal Mf(\rho).
                                                               \tag{MT-M15}
\]
For $f=h*h^\sharp$ this trace is
$\sum_\rho m_\rho|\mathcal Mh(\rho)|^2\ge0$.
However, if any $m_\rho>1$, there is no positive-definite
inner product on the full faithful jet quotient for which
every such adjoint-square convolution is a positive
operator. In particular, the trace positivity does not
upgrade to positivity of the entire nonsemisimple action.

\begin{proof}
At a zero $\rho$ let $a_j(f)=(\mathcal Mf)^{(j)}(\rho)/j!$.
The product rule in normalized jets is
\[
       a_k(f*g)=\sum_{j=0}^k a_j(f)a_{k-j}(g).
\]
Thus the block of $C_f$ at $\rho$ is a lower triangular
Toeplitz matrix with diagonal $a_0(f)$ and size $m_\rho$.
All its entries decrease rapidly with the ordinate,
uniformly in the at most $m_0$ indices.
The trace norm of a finite matrix is at most the sum
of the trace norms of its individual matrix units,
so the block trace norms are summable. This proves
the Hilbert trace-class assertion and its trace.

For nuclearity, use the rank-one expansion in these
matrix entries. The jet functionals are equicontinuous:
for a fixed $N>3/2$,
\[
 |(\mathcal Mg)^{(k)}(\rho)/k!|
 \le \frac1{k!}\int_{\mathbb R}|x|^k e^{x/2}|g(x)|\,dx
 \le2p_{N,0}(g),
\]
uniformly even before using $k<m_0$.
The interpolation vectors have polynomial seminorm
bounds, and the matrix entries have rapid decay.
Insert the summable scalar weight $(1+|T|)^{-2}$
as in MT-M13. With at most $m_0^2$ matrix units per
zero, this is an explicit nuclear representation.
The weighted involution has jet coefficients
\[
                   a_j(h^\sharp)=(-1)^j\overline{a_j(h)}.
\]
Its zeroth coefficient gives the asserted nonnegative
trace for an adjoint square.

Finally suppose $m=m_\rho>1$. The interpolating functions
in MT-M14 realize every local jet and zero jets elsewhere.
The local algebra is $\mathbb C[z]/(z^m)$, with
$z^\sharp=-z$ and conjugation on scalars.
Let $p$ denote its identity, extended by zero in every
other zero block, and let $n=z$ in this block and zero
elsewhere. Both are represented by actual functions
of $\mathcal E$, by (14). On the $m$-dimensional local
block, their convolution operators are $I$ and a
nonzero nilpotent $N$.
For every real $u$, the class $h=p+iu\,n$ satisfies
$h^\sharp=h$, and hence its adjoint square acts there as
\[
                         I+2iuN-u^2N^2.
\]
If all adjoint-square operators were positive in a
positive-definite inner product, these operators and
the corresponding operators with $-u$ would be
self-adjoint on the invariant local block.
Subtracting them for $u\ne0$ shows that $iN$ is
self-adjoint, equivalently $N^*=-N$.
A skew-adjoint operator on a finite-dimensional positive
inner-product space is normal and diagonalizable;
nilpotence forces all eigenvalues to vanish, so it is
zero. This contradicts the regular action of $z$
on $\mathbb C[z]/(z^m)$.
The argument is internal to a finite block and
does not depend on any choice of Hilbert completion.
\end{proof}

Passing to the semisimple quotient that retains only the
zeroth coordinate of each block removes the nilpotents.
Using $m_\rho$ copies of each scalar coordinate then
recovers the positive spectral trace with multiplicity.
It does not faithfully recover the discarded higher
jets. This is the precise distinction needed when
transferring Alpöge-style positivity to an orbital
jet quotient in Masullo 9.3.2.5--9.3.3, or to the
cote~1/SA23 trace programme. No assertion that the
draft explicitly claims a positive faithful Hilbert
realization of this entire quotient is made here.

\subsection*{MT-M16. The full arithmetic jet basis has one exact quantitative gate}
The bounded-multiplicity hypothesis in MT-M14 can be removed.
Fix a primitive unitary Hecke function, assume GRH, and use
the canonical $\Xi_\eta,\psi_\eta,Q_\eta$ of MT-M12.
Allow all actual multiplicities $m_\rho$, without a uniform
bound. Write $\rho=1/2+iT$ and set
\[
 P_\rho=
  \prod_{\substack{\nu\in Z,\ \nu\ne\rho\\
              |\operatorname{Im}\nu-T|\le1}}
              |T-\operatorname{Im}\nu|^{m_\nu},
\]
with empty product one. The following conditions are equivalent:
\begin{enumerate}
\item The first nonzero finite-$L$ Taylor coefficients satisfy
$|L^{(m_\rho)}(\rho,\eta)/m_\rho!|
                       \ge c(1+|T|)^{-C}$ uniformly in $\rho$.
\item The local products satisfy
$P_\rho\ge c(1+|T|)^{-C}$ uniformly in $\rho$.
\item The canonical closed convolution quotient has a
Schauder basis with the prescribed normalized Mellin
zero-jet coefficient functionals.
\item It has such an absolute Schauder basis, its closed
ideal is the full zero-jet kernel, and it is topologically
the rapid sequence space on
$\mathcal J=\{(\rho,k):0\le k<m_\rho\}$, weighted by $1+|T|$.
\end{enumerate}
Constants may differ between these equivalent conditions.
Only GRH for this fixed function is used analytically:
$A_D$ supplies it for Dirichlet functions and $M_H$
supplies it conditionally in general. The equivalent
quantitative conditions remain an extra gate; none is
being deduced here from zero location alone.

\begin{proof}
All constants below may depend on the fixed function and
on the indicated fixed strip. NT-Z10 gives both polynomial
strip growth of the finite $L$-function and a local count
including multiplicities. In particular
\[
       m_\rho=O(\log(2+|T|)),\qquad
       \sum_{\rho\in Z}m_\rho(1+|T|)^{-2}<\infty.       \tag{15}
\]
The finitely many low zeros can be included by adjusting
positive constants in every estimate.

\emph{Equivalence of the two quantitative conditions.}
Take the disk centered at $2+iT$ with a zero-free boundary
of radius $R\in[4,5]$. At large height it contains no
poles or trivial zeros. Its total zero multiplicity is
$M=O(\log(2+|T|))$. Remove its zeros by the finite
Blaschke product, with multiplicities, and write $L=GB$.
As proved in MT-M11, $G$ is bounded above by a polynomial
in $|T|$ and below by an inverse polynomial on the
inner half disk; this uses the Euler lower bound at
the center and Harnack's inequality.

The point $\rho$ lies in that inner disk. Its own leading
Blaschke coefficient is
$[R(1-|a|^2)]^{-m_\rho}$ with $a=-3/(2R)$.
Every other factor at $\rho$ is its distance from $\rho$
times a factor bounded above and below by fixed positive
constants. These statements follow directly from
$(z-a_j)/(1-\overline{a_j}z)$ and $|a|\le3/8$.
All other zeros of distance at most one are in the disk.
The distances of the remaining disk zeros are between
one and $2R$. Since there are $M$ factors with multiplicity,
all constant factors together cost at most a fixed power
of $2+|T|$. Therefore
\[
 c(1+|T|)^{-C}P_\rho
   \le\left|\frac{L^{(m_\rho)}(\rho,\eta)}{m_\rho!}\right|
   \le C(1+|T|)^C P_\rho.                              \tag{16}
\]
This proves (1) if and only if (2).

\emph{A radius adapted to the full multiplicity.}
Assume (2), and let $\delta_\rho$ be the lesser of one
and the distance to the nearest distinct zero.
If a distinct zero $\nu$ has this distance less than one,
the product $P_\nu$ contains the factor
$\delta_\rho^{m_\rho}$. Every other factor in $P_\nu$
is at most one. Hence
\[
 \delta_\rho^{m_\rho}\ge P_\nu
                         \ge c(1+|T|)^{-C}.
\]
When $\delta_\rho=1$ the same conclusion is immediate.
Put $r_\rho=\delta_\rho/4$. Equation (15) now gives
\[
                        r_\rho^{-m_\rho}
                                      \le C(1+|T|)^C. \tag{17}
\]
This is the quantitative replacement for the fixed number
of reciprocal-Taylor steps used in MT-M14.

\emph{Minimum modulus without bounded multiplicities.}
The disk proof of (11) works with all multiplicities once
(2) is available. We give the detail that makes the
constants independent of the largest multiplicity.
For a point $s=\sigma+it$ in a fixed strip, choose the
larger disk from MT-M12 and write again $L=GB$.
If that disk has no zeros, the lower bound for $G=L$
already gives the required polynomial lower bound.
If its nearest zero is outside the inner half disk,
every Blaschke factor at $s$ is at least $1/8$.
Otherwise let $\rho_0$ be a nearest disk zero in that
inner half disk. For every other zero,
$|\rho_0-\nu|\le2|s-\nu|$.
The product of the other distances, with multiplicity,
is at least $P_{\rho_0}$: omitted factors of distance
at most one only increase that product, while all
included distances greater than one are at least one.
The normalized Blaschke denominators and radii cost
only a constant to the total multiplicity
$M=O(\log(2+|t|))$. The lower bound for $G$ therefore
gives
\[
 |L(s,\eta)|\ge c_A(1+|t|)^{-C_A}
             \operatorname{dist}(s,Z)^{m_{\rho_0}}
                       \quad\hbox{if }\operatorname{dist}(s,Z)<1,
                                                               \tag{18}
\]
and a polynomial lower bound when that distance is
at least one. The choice of a nearest zero in (18)
is harmless in a tie. Zeros outside the larger disk
have distance greater than one from $s$.

\emph{Uniform Hermite interpolation.}
Put $m=m_\rho$, $z=s-\rho$, and
$G_\rho(s)=\Xi_\eta(s)/z^m$.
On the circle $|z|=r_\rho$, $\rho$ is the nearest zero.
Divide (18) by $|z|^m$ and use the fixed-strip lower
Stirling bound for the gamma factors. This gives
\[
             |G_\rho(s)|\ge
                 c(1+|T|)^{-C}e^{-\kappa|T|}
                            \quad(|z|=r_\rho).
\]
The function $G_\rho$ has no zero in that disk.
Consequently Cauchy's formula bounds the $j$th Taylor
coefficient of $e^{-az^2}/G_\rho(\rho+z)$ by
\[
       C(1+|T|)^C e^{\kappa|T|}r_\rho^{-j}.
                                                               \tag{19}
\]
For every $j<m$, (17) makes this a polynomial bound
times $e^{\kappa|T|}$, with constants independent of $j,m$.
Use exactly the Taylor truncation $P_{\rho,k}$ and
cardinal function (14) of MT-M14, now for all $k<m$.

Their fixed-strip bounds are uniform. When $|z|\ge1$,
the degree of $z^kP_{\rho,k}$ is at most $m-1$, so
\[
 |z^kP_{\rho,k}(z)/z^m|
       \le m\,C(1+|T|)^C e^{\kappa|T|}.
 \]
There is no growing positive power of $|z|$.
The factor $m$ is absorbed using (15), and the Gaussian
bound (8) follows from (5), just as before.
When $|z|<1$, the maximum principle applied to
$G_\rho$ on $|z|=2$ gives
$|G_\rho(s)|\le C(1+|T|)^C e^{-\kappa|T|}2^{-m}$.
The polynomial $z^kP_{\rho,k}$ is bounded by
$m$ times the coefficient bound (19), with (17).
This proves (8) on the remaining region as well,
after increasing a fixed constant.
Inverse transforms $h_{\rho,k}$ therefore satisfy
the polynomial seminorm bounds (9) uniformly over
every jet index $k<m_\rho$.

The normalized jets of any $f\in\mathcal E$ also have
uniform rapid height decay for \emph{all} $k\ge0$:
for each fixed $j$ and each fixed real-part strip,
\[
 \sup_{k\ge0}
 (1+|t|)^j\left|\frac{(\mathcal Mf)^{(k)}(\sigma+it)}{k!}\right|
                    \le C_{A,j}\,p_{N,j}'(f),          \tag{20}
\]
where $p_{N,j}'$ is a finite maximum of defining
seminorms, with sufficiently large fixed $N$.
Indeed, integrate by parts $j$ times in the transform
of $e^{\sigma x}x^k f(x)/k!$.
Every derivative of $x^k/k!$ is either zero or
$x^{k-\ell}/(k-\ell)!$; the remaining coefficients
depend only on $A,j$. For $N-A>1$, its weighted
$L^1$ norm is at most
$2/(N-A)^{k-\ell+1}\le2$, times the corresponding
seminorm of $f$. The region $|t|\le1$ follows from
the same integral bound without integration by parts.
This proves (20).

Equations (15), (20), and the uniform cardinal bounds
give absolute interpolation and its continuous inverse
for every rapid sequence on $\mathcal J$.
Finally, division proves the full ideal equality.
If $F=\mathcal Mf$ has all prescribed zero jets, its
Taylor remainder at a nearest zero is
\[
 F(s)=z^m\int_0^1m(1-u)^{m-1}
                  \frac{F^{(m)}(\rho+uz)}{m!}\,du .
 \]
The measure in this integral has total mass one.
Estimate (20) thus bounds the normalized remainder
uniformly in $m$, and its $|z|^m$ cancels the distance
factor in (18). The quotient $F/L$ is rapidly
decreasing on every fixed strip at large height.
The entire quotient $F/\Xi_\eta$, multiplied by
$e^{\epsilon s^2}$, is the Mellin transform of an
element of $\mathcal E$ for every $\epsilon>0$.
The Gaussian approximate-identity argument of
MT-M12 proves that $f$ lies in the closed ideal.
The reverse inclusion is immediate.
This yields the asserted topological rapid-sequence
identification and (4). Statement (4) implies (3),
and (3) implies (1) by MT-M11. All four conditions
are therefore equivalent.
\end{proof}

This removes both the simplicity assumption of MT-M12 and the
uniform multiplicity assumption of MT-M14. It does not prove the
remaining leading-coefficient or cluster-product bound.
The local $O(\log|T|)$ count is used together with that
quantitative condition; MT-M09 shows why the count alone
cannot replace it.

\subsection*{MT-M17. Nuclear trace with arbitrary zero multiplicities}
Under the equivalent conditions of MT-M16, every
$f\in\mathcal E$ acts by nuclear convolution on $Q_\eta$
and by a trace-class operator on $\ell^2(\mathcal J)$.
Its Hilbert trace is
\[
       \operatorname{Tr}C_f
                  =\sum_{\rho\in Z}m_\rho\mathcal Mf(\rho).
\]
For $f=h*h^\sharp$ this is the nonnegative sum
$\sum_\rho m_\rho|\mathcal Mh(\rho)|^2$.
The faithful-operator positivity obstruction of MT-M15
persists whenever any zero is multiple.

To prove these assertions without a uniform block-size
hypothesis, note that (15) bounds the total multiplicity
in each unit-height interval by $O(\log(2+|T|))$.
The number of Toeplitz matrix units in that interval is
therefore at most $O(\log^2(2+|T|))$.
Estimate (20) makes every matrix entry of $C_f$ uniformly
rapidly decreasing in height, independently of its jet
index. Thus the sum of the trace norms of all matrix
units is finite, giving the trace-class assertion and
the displayed diagonal trace. For nuclearity use the
equicontinuous normalized-jet functionals, the uniform
polynomial seminorm bounds for $h_{\rho,k}$, and the
scalar weight $(1+|T|)^{-2}$ for each matrix unit.
These weights are summable even with the
$O(\log^2(2+|T|))$ count, and the resulting vector
family is bounded in every seminorm by (20).
This is the same explicit nuclear decomposition as
MT-M15, now with all indices included.
The involution formula and the finite local nilpotent
argument in MT-M15 make no bounded-multiplicity
assumption, so they prove the remaining statements.

This is an exact reduction for the prescribed normalized-jet
coordinates. MT-M18 and MT-M19 below supply a different,
block-adapted absolute basis, spectral synthesis and nuclearity
without that quantitative condition. Identification with a
particular independently defined adelic quotient is a separate
map; MT-M21 begins the comparison with Meyer's construction.

\subsection*{MT-M18. Unconditional arithmetic block synthesis and an absolute adapted basis}

Fix a primitive unitary Hecke character \(\eta\) over a number field.
Use the entire completion \(\Xi=\Xi_\eta=P_\eta\Lambda(s,\eta)\)
and its inverse Mellin transform \(\psi=\psi_\eta\) from MT-M12,
but make no GRH or simplicity assumption.  The polynomial \(P_\eta\)
removes the possible pure-norm poles exactly as there.  Let \(Z\)
be the set of distinct zeros of \(\Xi\), with their actual
multiplicities \(m_\rho\).  These are the nontrivial zeros and
lie in \(0\le\operatorname{Re}s\le1\).
Put
\[
 \begin{gathered}
 \mathcal E=\{f\in C^\infty(\mathbb R,\mathbb C):
   p_{N,j}(f)=\sup_x e^{N|x|}|f^{(j)}(x)|<\infty
                    \text{ for all integers }N,j\ge0\},\\
 Q=\mathcal E/\overline{\psi*\mathcal E}.
 \end{gathered}
\]
All constants below concern this one fixed Hecke function; no
uniformity in its field, conductor, or infinity type is asserted.

There are real numbers
\[
 n+\tfrac14<h_n<n+\tfrac34\quad(n\in\mathbb Z)
\]
not equal to any zero ordinate, with the following properties.
Let \(\Omega_n\) be the interior of the rectangle
\([-1,2]+i[h_n,h_{n+1}]\), and orient its boundary
\(\Gamma_n\) positively.  Let \(A(\Omega_n)\) be the Banach
algebra of functions holomorphic in \(\Omega_n\) and continuous
on its closure, with the supremum norm.  Write
\[
 \begin{split}
 I_n&=\{F\in A(\Omega_n):
        F^{(j)}(\rho)=0\text{ for }\rho\in Z\cap\Omega_n,\
                                  0\le j<m_\rho\},\\
 J_n&=A(\Omega_n)/I_n,\qquad
 \|a\|_n=\inf_{F\bmod I_n=a}\sup_{\Gamma_n}|F|,\qquad
 d_n=\dim_{\mathbb C}J_n=\sum_{\rho\in Z\cap\Omega_n}m_\rho .
 \end{split}                                                    \tag{M18.1}
\]
An empty block has \(J_n=0\).  Then
\[
 d_n=O(\log(2+|n|)),\qquad
 s(J)=\{(a_n):q_k(a):=\sup_n(2+|n|)^k\|a_n\|_n<\infty
                                  \text{ for every }k\ge0\}      \tag{M18.2}
\]
is a Fr\'echet space, and the Mellin block restriction gives a
topological isomorphism
\[
 Q\simeq s(J).                                                    \tag{M18.3}
\]
In particular
\[
 \overline{\psi*\mathcal E}
 =\{f\in\mathcal E:(\mathcal Mf)^{(j)}(\rho)=0
                         \text{ for every }\rho,\ 0\le j<m_\rho\}.
                                                                  \tag{M18.4}
\]
The block expansion in \(Q\) converges absolutely in every
continuous seminorm.  Choosing an Auerbach basis in each
\((J_n,\|\cdot\|_n)\) gives an absolute Schauder basis of \(Q\)
with block-adapted continuous coefficient functionals.

Neither (M18.3) nor this last basis assertion prescribes the
individual zero jets as its coordinate functionals.  Thus the
result does not assert the quantitative condition of MT-M16.
The block partition and its adapted bases are choices; the
closed quotient in (M18.3) is the canonical one.

\begin{proof}
\emph{Analytic inputs and the transform topology.}
NT-Z10 proves unconditionally, including for a fixed
infinite-order unitary character, the polynomial upper bound
for the finite \(L(s,\eta)\) on every fixed vertical strip and
the bound \(O(\log(2+|T|))\) for the number of nontrivial
zeros, with multiplicities, in any fixed-height window.
It also gives
\[
                   |L(2+iT,\eta)|\ge\zeta_K(2)^{-1}.              \tag{M18.5}
\]
The standard completed Hecke functional equation locates all
the zeros of the entire completion in the fixed critical strip.
Exact classical source locators are Ramakrishnan--Valenza,
\emph{Fourier Analysis on Number Fields}, \S7.4,
Theorem 7-19, printed p.~279, and \S7.7, Theorem 7-28,
printed p.~289; the fixed-strip and local-count deductions
are proved in NT-Z10.  We use these unconditional facts,
not a zero-location consequence of either GRH input.

Put \(\kappa=\pi[K:\mathbb Q]/4\).  Stirling's formula and
the preceding strip bound give, for every fixed \(A>0\),
\[
 |\Xi(\sigma+it)|\le C_A(2+|t|)^{B_A}e^{-\kappa|t|}
                         \quad(|\sigma|\le A).                  \tag{M18.6}
\]
Cauchy's formula gives the same type of bound for \(\Xi'\)
on a smaller strip.  For large \(|t|\), the factor
\(\Xi(s)/L(s,\eta)\) is nonzero and has upper and lower
bounds by \(e^{-\kappa|t|}\) times fixed positive or negative
powers of \(2+|t|\), uniformly on each fixed strip.  Its
finitely located gamma exceptions are immaterial at large
height.

For clarity, the transform fact used below is elementary.
If an entire function \(H\) decreases faster than every
power of \(|t|\), uniformly on each strip \(|\sigma|\le A\),
its inverse Fourier integral on \(\sigma=0\) is the Mellin
transform inverse of an element of \(\mathcal E\).
For \(x\ge0\), shift the inverse contour to
\(\sigma=N+1\), and for \(x\le0\) shift it to
\(\sigma=-N-1\).  The horizontal edges tend to zero by the
uniform strip decay.  The \(j\)th derivative is bounded
by the integral of \(|s|^j|H(s)|\) on the shifted line,
times \(e^{-(N+1)|x|}\); these integrals are finite.
This proves all \(p_{N,j}\) bounds and continuity from the
corresponding strip seminorms.  Fourier inversion proves
the transform identity.  Conversely, integration by parts
in \(\mathcal Mf(\sigma+it)=\int e^{(\sigma+it)x}f(x)\,dx\)
shows rapid vertical decay, with each strip seminorm
bounded by finitely many \(p_{N,j}\).  Entireness follows
by exponential domination.  In particular (M18.6) proves
\(\psi\in\mathcal E\) and \(\mathcal M\psi=\Xi\), without GRH.

\emph{Selecting all horizontal boundaries at once.}
For \(n\in\mathbb Z\), let \(M_n\) count, with multiplicities,
the zeros whose ordinates lie in \([n-2,n+3]\).  Thus
\(M_n\le C\log(2+|n|)\).  On
\(I_n^*=(n+1/4,n+3/4)\), set
\[
 U_n(h)=\sum_{\substack{\rho\in Z\\
                    \operatorname{Im}\rho\in[n-2,n+3]}}
       m_\rho\log^+\frac1{|h-\operatorname{Im}\rho|}.
\]
Every summand has integral at most
\(\int_{-1}^1-\log|u|\,du=2\) over \(I_n^*\).
Averaging over its length \(1/2\) therefore permits a
choice \(h_n\in I_n^*\), avoiding every zero ordinate,
with \(U_n(h_n)\le 5M_n\); if \(M_n=0\), use its midpoint.
(The harmless slack from \(4M_n\) to \(5M_n\) gives a
positive-measure set when \(M_n>0\).)
Every zero with \(|h_n-\operatorname{Im}\rho|<1\)
is included in this sum.  Consequently
\[
 \prod_{\rho\in Z}
   \min\{1,|h_n-\operatorname{Im}\rho|\}^{m_\rho}
      \ge e^{-5M_n}\ge c(2+|n|)^{-C}.                            \tag{M18.7}
\]
Only finitely many factors in this product differ from one.
Also \(1/2<h_{n+1}-h_n<3/2\), so the rectangles have
uniformly bounded perimeter and cover all zero ordinates.

We claim, for every fixed \(A>0\), that
\[
 \min_{|\sigma|\le A}|\Xi(\sigma+i h_n)|
       \ge c_A(2+|n|)^{-C_A}e^{-\kappa|n|}.                      \tag{M18.8}
\]
Here is a complete minimum-modulus argument.  For large
\(|h_n|\), center a disk at \(c_n=2+i h_n\) and take a
zero-free boundary of radius
\(R\in[4(A+3)+4,4(A+3)+5]\).  Such a radius exists because
there are only finitely many zeros in that compact annulus.
At sufficiently large height the disk contains no finite-\(L\)
pole or trivial zero, since their possible ordinates are fixed
by the infinity type.  Its total zero multiplicity is
\(M=O_A(\log(2+|n|))\), by NT-Z10.  Remove its zeros by the
finite disk Blaschke product \(B\), with multiplicities,
and write \(L=GB\).  The factor \(G\) is holomorphic and
zero-free on the disk and its boundary.  Since \(|B|=1\)
on the boundary, (M18.6)'s finite-\(L\) antecedent and the
maximum principle give \(|G|\le C(2+|n|)^C\) there and
inside.  At the center, \(|B(c_n)|\le1\), so (M18.5)
gives \(|G(c_n)|\ge c>0\).
Apply Harnack's inequality to the nonnegative harmonic
function \(\log(C(2+|n|)^C)-\log|G|\).  On the inner
half disk its value is at most three times its value
at the center.  Thus
\[
                        |G(s)|\ge c(2+|n|)^{-C}
\]
on that inner disk.  This form of Harnack follows by
bounding the ratios of the Poisson kernels on concentric
circles; the ratio is \((R+r)/(R-r)\le3\) for \(r\le R/2\).

For \(s=\sigma+i h_n\), \(|\sigma|\le A\), the normalized
coordinate \(z=(s-c_n)/R\) has \(|z|<1/4\).
For a disk zero \(\rho\), with \(a_\rho=(\rho-c_n)/R\),
\[
 \left|\frac{z-a_\rho}{1-\overline{a_\rho}z}\right|
    \ge\frac{|s-\rho|}{2R}
    \ge\frac{\min\{1,|h_n-\operatorname{Im}\rho|\}}{2R}.
\]
Products of omitted factors in (M18.7) are at most one.
It follows that
\[
 |B(s)|\ge(2R)^{-M}
   \prod_{\rho\in Z}
      \min\{1,|h_n-\operatorname{Im}\rho|\}^{m_\rho}
             \ge c_A(2+|n|)^{-C_A}.
\]
Multiplication by \(G\) and by the lower Stirling bound
for \(\Xi/L\) proves (M18.8) at large height.
There are only finitely many remaining \(n\); each
horizontal segment is compact and zero-free for \(\Xi\),
so reducing \(c_A>0\) covers them.

The same disk proof, now centered at \(2+it\), gives a
polynomial lower bound for \(L(\sigma+it,\eta)\) on either
fixed vertical line \(\sigma=-A\) or \(\sigma=A\), for
every \(A\ge2\) and large \(|t|\).  Indeed their distances
from all disk zeros in the critical strip are at least one,
so every normalized Blaschke factor is at least \(1/(2R)\).
There is no logarithmic-distance product to control.
The lower Stirling bound and the compact low-height range
give the corresponding lower bound for \(\Xi\).
The same argument applies to the line \(\sigma=-1\)
and to \(\sigma=2\).  Combining it with (M18.8) yields
\[
 \inf_{w\in\Gamma_n}|\Xi(w)|
       \ge c(2+|n|)^{-C}e^{-\kappa|n|}.                          \tag{M18.9}
\]
All constants here are independent of separations between
zeros inside a rectangle.

\emph{The finite block algebras and an entire interpolation kernel.}
All zeros of a block are interior points.  Each of its
finitely many jet functionals is continuous on
\(A(\Omega_n)\), by Cauchy's formula on a disk contained
in \(\Omega_n\).  The map to these jets is onto, by
ordinary polynomial Hermite interpolation.  Its kernel
is \(I_n\), a closed ideal, proving that (M18.1) is a
genuine finite-dimensional quotient norm, with dimension
\(d_n\).  The maximum principle identifies the supremum
over the closure with that over \(\Gamma_n\).
The quotient norm is submultiplicative because the
supremum norm is.  The local zero count gives (M18.2)'s
dimension bound.

Fix \(a>0\).  For \(F\in A(\Omega_n)\) define
\[
 H_nF(z)=\frac1{2\pi i}\int_{\Gamma_n}
    F(w)\frac{\Xi(z)-\Xi(w)}{\Xi(w)(z-w)}
                      e^{a(z-w)^2}\,dw .                       \tag{M18.10}
\]
The apparent singularity at \(z=w\) is removable.
Consequently this is an entire function of \(z\).
If \(z\) is inside the rectangle, Cauchy's formula
for the second term gives
\[
 H_nF(z)=F(z)+\Xi(z)\frac1{2\pi i}
             \int_{\Gamma_n}
               \frac{F(w)e^{a(z-w)^2}}{\Xi(w)(z-w)}\,dw .
                                                                  \tag{M18.11}
\]
The integral on the right is holomorphic in the interior.
If \(z\) is outside the closed rectangle, that second
Cauchy integral is zero and \(H_nF(z)\) is divisible
locally by \(\Xi(z)\).
These statements follow also for continuous boundary
values by applying the formulas on interior rectangles
and taking limits.  Thus \(H_nF\) matches every
prescribed block jet of \(F\), and vanishes to the
full multiplicity at every zero outside the block.

Moreover, if all the block jets of \(F\) vanish, then
\(F/\Xi\) is holomorphic in \(\Omega_n\) and continuous
on its closure.  For exterior \(z\), both terms in
(M18.10) then integrate to zero by Cauchy's theorem.
Entireness implies \(H_nF=0\) everywhere.  It follows
that \(H_n\) depends only on \(a_n=F\bmod I_n\).
Denote this entire function by \(H_n a_n\).

For every fixed \(A>0\) there are constants \(C_A,B_A'>0\)
and \(c_a>0\) such that
\[
 |H_n a_n(\sigma+it)|
       \le C_A(2+|n|)^{B_A'}\|a_n\|_n
                         e^{-c_a(t-n)^2},
                 \qquad |\sigma|\le A.                         \tag{M18.12}
\]
To check the estimate before taking the quotient infimum,
bound \(|F|\) by its boundary supremum in (M18.10).
If \(|z-w|\ge1\), use (M18.6) and (M18.9) to bound
the difference quotient by
\[
 C_A(2+|n|)^C(2+|t|)^C
                    e^{\kappa\,\lvert t-\operatorname{Im}w\rvert}
       +1 .
\]
Here \(|\operatorname{Im}w-n|<2\), while
\[
 |e^{a(z-w)^2}|
       \le C_A e^{-a(t-\operatorname{Im}w)^2}.
\]
A fixed polynomial and the fixed linear exponential
are absorbed by a smaller Gaussian, and
\(2+|t|\le(2+|n|)(2+|t-n|)\).
If \(|z-w|<1\), the identity
\[
 \frac{\Xi(z)-\Xi(w)}{z-w}
      =\int_0^1\Xi'(w+u(z-w))\,du
\]
and the strip derivative bound give a polynomial in
\(2+|n|\) after division by (M18.9).
In this region \(|t-n|<3\), so it too satisfies (M18.12)
after a constant adjustment.  The perimeter is bounded
independently of \(n\).  Taking the infimum over all
representatives of \(a_n\) proves (M18.12).

The transform fact above now gives a unique
\(\mathcal T_n(a_n)\in\mathcal E\) whose Mellin transform is
\(H_n a_n\).  It is linear in \(a_n\), and for every
defining seminorm
\[
 p_{N,j}(\mathcal T_n(a_n))
       \le C_{N,j}(2+|n|)^{B_{N,j}}\|a_n\|_n .                  \tag{M18.13}
\]
Indeed on either shifted vertical line, the integral
of \((1+|t|)^j\) against (M18.12)'s Gaussian costs
only another fixed power of \(2+|n|\).
In particular no individual reciprocal derivative
or Hermite interpolation matrix norm has entered
these estimates.

\emph{Spectral synthesis for the closed ideal.}
Write \(J\) for the right side of (M18.4).
It is a closed subspace, since each Mellin jet is
continuous.  Multiplication of Mellin transforms
immediately gives
\(\overline{\psi*\mathcal E}\subseteq J\).
For the reverse inclusion take \(f\in J\) and
\(F=\mathcal Mf\).  The quotient \(F/\Xi\) is entire:
every zero, including its full multiplicity, has
been cancelled.

Fix \(A\ge2\).  Use the wider rectangles
\([-A,A]+i[h_n,h_{n+1}]\).  On their horizontal
edges (M18.8) applies; on their vertical edges
the vertical-line argument following (M18.8)
applies.  The rapid fixed-strip decay of \(F\)
and these lower bounds show, on the whole boundary,
\[
 \left|\frac{F(s)}{\Xi(s)}\right|
       \le C_{A,f}(2+|n|)^{C_A}e^{\kappa|n|}.
\]
The maximum principle for the entire quotient
extends this bound through the rectangle.
Since their heights have gaps between \(1/2\)
and \(3/2\), it follows that on every fixed strip
\[
 \left|\frac{F(\sigma+it)}{\Xi(\sigma+it)}\right|
       \le C_{A,f}(2+|t|)^{C_A}e^{\kappa|t|}.                  \tag{M18.14}
\]
For every \(\epsilon>0\), the entire function
\[
                   e^{\epsilon s^2}F(s)/\Xi(s)
\]
therefore has rapid decay on every fixed strip.
It is the Mellin transform of some
\(j_\epsilon\in\mathcal E\).
With \(g_\epsilon(x)=(4\pi\epsilon)^{-1/2}
e^{-x^2/(4\epsilon)}\), Fourier uniqueness gives
\[
                         \psi*j_\epsilon=g_\epsilon*f.
\]
The right side tends to \(f\) in \(\mathcal E\).
For completeness, on \(|y|\le1\) the weighted
mean-value bound for
\(f^{(j)}(x-y)-f^{(j)}(x)\) is
\(|y|e^{N|y|}p_{N,j+1}(f)\).
On the complement of any fixed neighborhood of
zero use
\((e^{N|y|}+1)p_{N,j}(f)\).
Integration against \(g_\epsilon(y)\) makes both
bounds tend to zero, by its first moment and
Gaussian tails.  This proves the reverse inclusion
and hence (M18.4), independently of the interpolation
construction.

\emph{The topological decomposition and the adapted basis.}
The restriction map
\[
 R:\mathcal E\longrightarrow s(J),\qquad
        (Rf)_n=(\mathcal Mf)|_{\Omega_n}\bmod I_n
\]
is continuous.  Indeed the quotient norm is at most
\(\sup_{\Gamma_n}|\mathcal Mf|\), and these boundary
suprema decrease faster than every power of \(|n|\),
with each such bound controlled by finitely many
seminorms of \(f\).  Define
\[
                    S(a)=\sum_{n\in\mathbb Z}\mathcal T_n(a_n).
\]
For each seminorm \(p_{N,j}\), (M18.13) and
\(q_{B_{N,j}+2}(a)\), with the exponent rounded up
if necessary, bound the sum of its term seminorms
by a constant times \(\sum_n(2+|n|)^{-2}\).
Thus this series converges absolutely in every
seminorm of \(\mathcal E\), and defines a continuous
linear map \(S:s(J)\to\mathcal E\).
Completeness of \(s(J)\) follows directly from the
complete finite-dimensional norms in its weighted
supremum seminorms; completeness of \(\mathcal E\)
follows from convergence of all weighted derivative
seminorms on compact intervals.

Every Mellin jet of \(S(a)\) may be computed
termwise by continuity.  The interpolation property
of (M18.11) gives \(RS=\mathrm{id}\).
Also \(\ker R=J=\overline{\psi*\mathcal E}\).
Therefore \(R\) induces the continuous bijection
(M18.3), with continuous inverse
\(a\mapsto[S(a)]\).
The block operators \(\mathcal T_n R_n\) on \(\mathcal E\)
are finite-rank continuous projections, mutually
annihilating for distinct \(n\).
Their absolute sum is the continuous projection
\(SR\); on \(Q\) their sum is the identity.

Finally every finite-dimensional complex normed
space has an Auerbach basis \(v_{n,1},\ldots,v_{n,d_n}\):
its vectors and dual coordinate functionals
\(\phi_{n,j}\) all have norm one.
Here is the elementary existence argument.  In
any fixed coordinates maximize the absolute
determinant of \(d_n\) vectors in the compact unit
ball.  A maximum is positive.  Replacing its
\(j\)th column by any unit-ball vector multiplies
the determinant by the corresponding dual
coordinate.  Maximality makes that coordinate
have modulus at most one; its value on its own
column is one.  This proves the claimed norms.

For \(a=(a_n)\in s(J)\), its resulting scalar
coefficients satisfy
\(|\phi_{n,j}(a_n)|\le\|a_n\|_n\).
Writing \(e_{n,j}\) for \(v_{n,j}\) in its block
and zero elsewhere, one has
\[
 \sum_{n,j}q_k\bigl(\phi_{n,j}(a_n)e_{n,j}\bigr)
 \le\sum_n d_n(2+|n|)^k\|a_n\|_n
 \le q_{k+3}(a)\sum_n d_n(2+|n|)^{-3}<\infty .
\]
This proves an absolute expansion in every
seminorm, in any enumeration of these vectors.
The coordinate functionals are continuous and
unique by the block isomorphism and the finite
basis in each block.  Transport through (M18.3)
proves the asserted absolute Schauder basis.
\end{proof}

The spectral block \(J_n\), as an algebra, is the
finite-length complex algebra
\[
 \mathbb C[s]\Big/
       \prod_{\rho\in Z\cap\Omega_n}(s-\rho)^{m_\rho}
 \simeq
 \prod_{\rho\in Z\cap\Omega_n}
                  \mathbb C[z]/(z^{m_\rho}).
\]
Its stable norm is the quotient holomorphic supremum
norm in (M18.1), not a norm obtained by declaring its
individual jet coordinates orthonormal.
Thus the construction is a finite-length spectral
divisor decomposition compatible with the Mellin
convolution algebra.  It directly refines the repair
of Masullo's Lemma 9.3.2.5, original PDF pp.~425--426:
closure of the convolution ideal is essential, and
the adapted basis has different coordinates from
that lemma's prescribed jets.  The inherited programme
locators are cote~1's nuclear-trace framework and
SA23 geo-H-lidskii/geo-H-genus, as carried in the
attached returns; no new inspection of original
cote scans is claimed here.  Nuclearity and the
trace action of these finite spectral algebras
are treated separately in MT-M19.

No first-priority claim is made for spectral synthesis in related
function spaces; the explicit block norms and bounds above supply
the precise topology and interpolation statement used here.



\subsection*{MT-M19. Unconditional nuclear convolution and its exact positivity criterion}
Use the canonical quotient, stable block algebras \(J_n\), and
topological isomorphism \(Q\simeq s(J)\) of MT-M18. No GRH,
simplicity, or reciprocal-derivative bound is assumed initially.
For every \(f\in\mathcal E\), convolution \(C_f\) is nuclear on
\(Q\). There is a Hilbert completion \(\mathcal H_B\) obtained
by making an Auerbach basis in each block orthonormal, on which
every \(C_f\) is trace class, with
\[
 \operatorname{Tr}_{\mathcal H_B} C_f
       =\sum_{\rho\in Z}m_\rho\mathcal Mf(\rho).       \tag{M19.1}
\]
The series is absolutely convergent and the trace is independent
of all block and Auerbach choices. Nuclearity is uniform when
\(f\) ranges over a bounded subset of \(\mathcal E\).

Define \(h^\sharp(x)=e^{-x}\overline{h(-x)}\). Then
\[
 \operatorname{Tr}_{\mathcal H_B} C_{h*h^\sharp}\ge0
          \quad\hbox{for every }h\in\mathcal E
 \quad\Longleftrightarrow\quad
 \text{GRH for }L(s,\eta).                              \tag{M19.2}
\]
Under GRH the trace equals
\(\sum_\rho m_\rho|\mathcal Mh(\rho)|^2\).
This concerns the trace; the chosen block Hilbert norm need
not make the operator \(C_{h*h^\sharp}\) positive.

\begin{proof}
Put \(F=\mathcal Mf\) and \(b_n(f)=\sup_{w\in\Gamma_n}|F(w)|\).
These numbers decrease faster than every power of \(2+|n|\),
with each decay bound controlled by finitely many defining
seminorms of \(f\). This follows from the integration-by-parts
estimate in MT-M18 on the one fixed strip containing all
\(\Gamma_n\). Convolution preserves the closed ideal: the
identity \(\mathcal M(f*g)=\mathcal Mf\,\mathcal Mg\) holds
by absolute integrability, and convolution is continuous on
\(\mathcal E\). Explicitly its seminorms are bounded by
\[
 p_{N,j}(f*g)
 \le p_{N,j}(f)\int_{\mathbb R}e^{N|y|}|g(y)|\,dy
 \le 2p_{N,j}(f)p_{N+1,0}(g).
\]
On \(J_n\), the induced map \(B_n(f)\) is multiplication by
the restriction of \(F\). The quotient Banach-algebra norm
therefore gives
\[
                      \|B_n(f)\|\le b_n(f).            \tag{M19.3}
\]

Let \(v_{n,j},\phi_{n,j}\), \(1\le j\le d_n\), be the
Auerbach vectors and dual functionals of MT-M18, all of norm
one. Regard each vector as supported in its own block and
write \(\ell_{n,j}(a)=\phi_{n,j}(a_n)\). The family
\(\ell_{n,j}\) is equicontinuous, since
\(|\ell_{n,j}(a)|\le q_0(a)\). Set
\[
 \lambda_{n,j}=(2+|n|)^{-2},\qquad
 u_{n,j}=(2+|n|)^2 B_n(f)v_{n,j}.
\]
The sum of \(|\lambda_{n,j}|\) is finite because
\(d_n=O(\log(2+|n|))\). For each \(k\),
\[
 q_k(u_{n,j})\le(2+|n|)^{k+2}b_n(f),
\]
so \(u_{n,j}\) is a bounded family in \(s(J)\). The series
\[
 C_f a=\sum_{n,j}\lambda_{n,j}\ell_{n,j}(a)u_{n,j}       \tag{M19.4}
\]
is thus an explicit nuclear decomposition, absolutely convergent
in every seminorm. Its value is the block multiplication map
because the finite Auerbach expansion in each block is exact.
The same estimates are uniform on a bounded set of \(f\)'s,
which proves the asserted uniformity.

Give the scalar Auerbach coordinates their ordinary square-sum
norm, and let \(\mathcal H_B\) be its completion. For each
nonempty block this Euclidean norm lies between
\(d_n^{-1/2}\|a_n\|_n\) and
\(d_n^{1/2}\|a_n\|_n\), since the vectors and coordinate
functionals have norm one. Consequently the inclusion
\(s(J)\to\mathcal H_B\) is continuous and injective, and
has dense image because it contains all finitely supported
block vectors. Every matrix entry of \(B_n(f)\) in these
orthonormal coordinates has modulus at most \(b_n(f)\),
by (M19.3). Expressing it as a sum of rank-one matrix units
therefore gives
\[
       \|B_n(f)\|_1\le d_n^2b_n(f),\qquad
       \sum_n\|B_n(f)\|_1<\infty.                     \tag{M19.5}
\]
The block direct sum is trace class and extends \(C_f\).

Algebraically, polynomial Hermite interpolation identifies
\(J_n\) with
\(\prod_{\rho\in Z\cap\Omega_n}\mathbb C[z]/(z^{m_\rho})\).
In the monomial basis of a local factor, multiplication by
\(F(\rho+z)\) is triangular with its constant term
\(F(\rho)\) on every diagonal entry. Its finite-dimensional
trace is \(m_\rho F(\rho)\). Trace is invariant under a
change of basis, so (M19.5) proves (M19.1). It also proves
the choice-independence stated there. Absolute convergence
follows directly from
\(\sum_n d_n b_n(f)<\infty\). All traces in the statement
are the ordinary Hilbert traces just constructed; no
unjustified trace on an arbitrary locally convex space is used.

Direct substitution gives
\[
 \mathcal Mh^\sharp(s)
       =\overline{\mathcal Mh(1-\bar s)}.              \tag{M19.6}
\]
The unitary Hecke functional equation implies that the zero
involution \(\rho\mapsto\rho^\dagger=1-\bar\rho\)
preserves \(Z\) and its multiplicities. Indeed it relates
\(\Lambda(s,\eta)\) to \(\Lambda(1-s,\bar\eta)\),
while complex conjugation relates the latter character to
\(\eta\). The pole-removal polynomial introduces no zeros
into \(Z\). Formula (M19.6) therefore preserves the full
zero-jet ideal, so \(\sharp\) descends to an antilinear
algebra involution of \(Q\). If GRH holds then \(\rho^\dagger=\rho\),
and (M19.1), (M19.6) give the stated nonnegative sum.

Conversely, suppose \(\rho\ne\rho^\dagger\) is a pair
off the critical line, of common multiplicity \(m\).
They have the same ordinate and hence lie in the same
block. In that block prescribe the constant jet 1 at
\(\rho\), the constant jet \(-1\) at \(\rho^\dagger\),
and the zero jet at every other zero; prescribe zero in
all other blocks. This is a finitely supported element
of \(s(J)\). The right inverse in MT-M18 realizes it as
\(\mathcal Mh\) for some \(h\in\mathcal E\). Formula
(M19.1) then has just two nonzero contributions for
\(h*h^\sharp\), both equal to \(-m\). Its trace is
\(-2m<0\), contradicting the left side of (M19.2).
This proves the equivalence.
\end{proof}

This is a precise zero-side Weil positivity criterion in
the canonical quotient, with an explicit unconditional
nuclear and trace-class realization. It supplies no proof
of positivity independent of the zeros. The trace is the
sum of traces on finite-length spectral algebras; multiple
zeros contribute their lengths. In the programme of
cote~1 and SA23 geo-H-lidskii/geo-H-genus, MT-M21--MT-M22 identify the specified Meyer adelic summation
quotient with this model, and MT-M24 identifies the actual
Masullo smooth quotient. MT-M25 treats the actual line Hilbert
completion and its lost jets. A cohomological realization and
a positivity argument defined independently of the zeros remain separate. Adopted \(A_D\)
supplies (M19.2)'s positive conclusion for Dirichlet
functions; \(M_H\) supplies it conditionally in general.

\subsection*{MT-M20. A weaker positive Hilbert completion and its exact obstruction}
Keep the unconditional canonical quotient of MT-M18.
If all nontrivial zeros are simple, individual Mellin
evaluation induces a continuous injective map with dense
image
\[
       Q\longrightarrow\ell^2(Z),\qquad
       [g]\longmapsto(\mathcal Mg(\rho))_{\rho\in Z}.   \tag{M20.1}
\]
It therefore defines a positive-definite continuous Hilbert
norm on \(Q\). Every convolution extends diagonally to a
trace-class operator on this completion. Under GRH,
\(C_{h^\sharp}=C_h^*\), and \(C_{h*h^\sharp}\) is a
positive operator. No reciprocal-derivative bound is
needed for this weaker Hilbert norm.

More precisely, the following are equivalent:
\begin{enumerate}
\item GRH holds for \(L(s,\eta)\), and every nontrivial
zero is simple.
\item There is a positive-definite continuous inner
product on \(Q\) such that every convolution extends to
a bounded operator on its Hilbert completion and every
\(C_{h*h^\sharp}\) is positive there.
\end{enumerate}
The inner product is required to be faithful on the
\emph{full jet quotient}. Semisimplifying its multiple-zero
local algebras changes this requirement.

\begin{proof}
Point evaluation at any zero in \(\Omega_n\) has norm at
most one on the quotient \(J_n\), by the maximum principle
and the definition of its quotient norm. Thus
\[
 \sum_{\rho\in Z}|\mathcal Mg(\rho)|^2
 \le\sum_n d_n\|(Rg)_n\|_n^2
 \le q_2(Rg)^2\sum_n d_n(2+|n|)^{-4}<\infty.           \tag{M20.2}
\]
This proves continuity. When all zeros are simple, zero
values mean all zero jets, and MT-M18 identifies the kernel
exactly with the closed ideal; hence the map is injective.
Every finitely supported zero-value vector lies in its
image, since it specifies elements of finitely many
\(J_n\)'s and the interpolation right inverse realizes
them. This proves density and (M20.1).

Convolution by \(f\) has diagonal entries
\(\mathcal Mf(\rho)\). Their absolute sum is finite by
the local zero count and rapid strip decay, so its extension
is trace class. Under GRH (M19.6) makes the entries for
\(h^\sharp\) the complex conjugates of those for \(h\).
The adjoint and positivity statements follow. This proves
the implication (1) to (2).

For the converse, finite-support elements of the block
algebras are actual elements of \(Q\), by MT-M18, and
convolution realizes their algebra multiplication. If
\(\rho\ne\rho^\dagger\), let \(e,e^\dagger\) be the
local idempotents equal to 1 in the full local algebras
at those two zeros and zero in every other local algebra.
The involution interchanges them. For \(h=e-e^\dagger\)
one has
\[
                  h*h^\sharp=-(e+e^\dagger).
\]
On the nonzero finite-dimensional invariant subspace
\((e+e^\dagger)Q\), convolution by this element is
\(-\mathrm{id}\). Its quadratic form is strictly negative
on every nonzero vector in any positive-definite inner
product. This contradicts (2), and proves GRH.

Now suppose a critical-line zero has multiplicity
\(m\ge2\). In its local algebra
\(\mathbb C[z]/(z^m)\), the involution is
\(z^\sharp=-z\), by (M19.6). Let \(e\) be that local
idempotent and \(n\) the element represented by \(z\),
both zero in all other local factors. They are realized
by elements of \(Q\). On the finite-dimensional invariant
space \(eQ\), write \(N=C_n\). It is a nonzero nilpotent
operator. For real \(u\), the element \(h_u=e+iun\)
satisfies \(h_u^\sharp=h_u\), and
\[
 C_{h_u*h_u^\sharp}|_{eQ}
                     =I+2iuN-u^2N^2.                \tag{M20.3}
\]
By (2) these operators are positive, in particular
self-adjoint, for both \(u=1\) and \(u=-1\).
Subtracting them shows that \(iN\) is self-adjoint
in the restricted positive-definite inner product.
A self-adjoint nilpotent on a finite-dimensional
inner-product space is zero: diagonalize the Hermitian
matrix, whose real eigenvalues must all vanish.
This contradicts \(N\ne0\). Thus all zeros are simple,
completing the converse.
\end{proof}

The Hilbert norm in (M20.1) can be strictly weaker than
the original Fr\'echet topology. Its existence therefore
does not make point evaluations a Schauder coordinate
system for that topology: MT-M16 still gives the exact
extra leading-coefficient condition for that assertion.
With multiple zeros, evaluation kills the higher-jet
nilpotent directions. A semisimple Hilbert model with
\(m_\rho\) copies of each scalar eigenvalue retains
the multiplicity-weighted trace, but is not a faithful
Hilbert completion of this full jet quotient. No
simplicity consequence of \(A_D\) or \(M_H\) is claimed.


\subsection*{MT-M21. Identification with Meyer's zeta-summation quotient over \(\mathbb Q\)}
For the Riemann zeta function, the canonical quotient in MT-M18
is naturally isomorphic, as a topological convolution module,
to the zero-space in Meyer's classical spectral realization.
This identification is unconditional and includes all zero
multiplicities.

Here are the precise source and conventions. In Ralf Meyer,
\emph{A spectral interpretation for the zeros of the Riemann
zeta function}, G\"ottingen Seminars 2004--05, pp.~117--137,
let \(H_- =\mathcal O(\mathbb R_+^\times)\) be the strong
Schwartz space of Section~2, let \(H_+\) be the even Schwartz
space on \(\mathbb R\), and let
\[
 H_\cap=\{u\in H_+:u(0)=\widehat u(0)=0\},\qquad
 (Zu)(x)=\sum_{k\ge1}u(kx).
\]
His Theorem~3.3, printed p.~124 (PDF p.~8), proves the
closed-image assertion for the zeta summation map; Section~4,
printed pp.~125--126, defines
\(H_-^0=H_-/ZH_\cap\).
Theorem~4.1, printed p.~126 (PDF p.~10), characterizes its
Mellin-transform kernel by two half-plane division conditions.
Under the logarithmic coordinate map
\[
       \mathcal L:\mathcal E\longrightarrow H_-,
       (\mathcal Lf)(x)=f(\log x),
\]
one has the exact equality of closed subspaces
\[
          \mathcal L\bigl(\overline{\psi_\zeta*\mathcal E}\bigr)
                         = ZH_\cap.                  \tag{M21.1}
\]
Consequently \(Q_\zeta\simeq H_-^0\), and the adapted
absolute basis, nuclear convolution, and trace constructions
of MT-M18--MT-M20 apply to this independently defined
zeta-summation quotient with their stated positivity hypotheses.

\begin{proof}
The map \(\mathcal L\) is a topological isomorphism:
the invariant derivatives \((x\,d/dx)^j\) become
\(d^j/dt^j\) for \(t=\log x\), and requiring rapid
decay after every real power weight \(x^\sigma\)
is equivalent to the defining seminorms
\(\sup_t e^{N|t|}|f^{(j)}(t)|\). Polynomial factors in
\(t\), if included in the Schwartz definition, are
absorbed by increasing \(N\). Its Mellin transform is
exactly \(\mathcal Mf(s)\), since \(d^\times x=dt\).
Multiplicative convolution becomes additive convolution.

Let \(F=\mathcal Mf\). Meyer's Theorem~4.1 says that
\(\mathcal Lf\in ZH_\cap\) precisely when
\[
 t\longmapsto F(\sigma+it)/\zeta(\sigma+it)
          \text{ is Schwartz for every }\sigma\ge\tfrac12,
                                                               \tag{M21.2}
\]
and
\[
 t\longmapsto F(\sigma+it)/\zeta(1-\sigma-it)
          \text{ is Schwartz for every }\sigma\le\tfrac12.
                                                               \tag{M21.3}
\]
The conditions include holomorphic removal of the corresponding
zeros on those lines. If they hold, every nontrivial zero
\(\rho\) of real part at least \(1/2\) is a zero of \(F\)
to at least its full multiplicity, by (M21.2). For zeros in
the other half-plane, the zeta functional equation says that
\(\zeta(1-s)\) has the same nontrivial zero divisor as
\(\zeta(s)\); condition (M21.3) gives the same conclusion.
Thus \(f\) belongs to the full zero-jet kernel in MT-M18.

Conversely, suppose \(F\) vanishes on this entire divisor.
The function \(F/\zeta\) is holomorphic on
\(\operatorname{Re}s>-2\). All nontrivial zero singularities
are removed, the pole of \(\zeta\) at 1 gives a removable
zero of its reciprocal, and its first trivial zero is at
\(-2\), outside this open half-plane.

Fix \(A\ge2\). Use rectangles
\([-1,A+1]+i[h_n,h_{n+1}]\), with the horizontal levels
constructed in MT-M18 for \(\zeta\). The finite-\(L\)
Blaschke estimates in that proof give a fixed inverse
polynomial lower bound for \(|\zeta|\) on the horizontal
edges at large height. The two vertical edges are a fixed
positive distance from every nontrivial zero, so the same
disk argument gives their polynomial lower bound. They
also avoid all trivial zeros. The finitely many low-height
boundaries cause no problem: the levels \(h_n\) are never
zero, and both vertical lines are zero-free.

For every \(N\), rapid fixed-strip decay of \(F\)
therefore bounds \(|F/\zeta|\) on these boundaries by
\(C_{A,N,f}(2+|n|)^{-N}\), after increasing the decay
power used for \(F\). The maximum principle for the
holomorphic quotient gives the same bound inside each
rectangle. Since adjacent heights have uniformly bounded
gaps, \(F/\zeta\) decreases rapidly on the whole strip
\([-1,A+1]\). Cauchy's formula on small fixed disks
inside that strip gives rapid decay of all its vertical
derivatives on \([1/2,A]\). This proves (M21.2).

Apply this same argument to \(F^\vee(s)=F(1-s)\).
It is an entire strip-Schwartz function and vanishes on
the full nontrivial divisor, because that divisor is
invariant under \(s\mapsto1-s\). The result for
\(F^\vee/\zeta\), followed by \(s\mapsto1-s\),
is exactly (M21.3). Meyer's image criterion now gives
\(\mathcal Lf\in ZH_\cap\).

We have proved that \(\mathcal L^{-1}(ZH_\cap)\)
is precisely the full zero-jet kernel. MT-M18 identifies
that kernel with \(\overline{\psi_\zeta*\mathcal E}\),
proving (M21.1). Both are closed subspaces of the
topologically identified test-function spaces, so their
quotients, with quotient topology, are naturally
topologically isomorphic. The logarithmic identification
commutes with scaling and integrated convolution, giving
the stated module assertion.
\end{proof}

This settles a specified arithmetic comparison map, beyond
a zero-built abstract model. It is a comparison with
Meyer's Poisson/zeta-summation construction, not an
identification with an unspecified adelic quotient in
Masullo's draft or with an \(\ell\)-adic cohomology group.
The classical spectral realization itself is Meyer's
result; the proof here supplies the explicit closed-ideal
and stable-block crosswalk used in this investigation.
The full number-field idele-class construction is located
in Meyer, \emph{On a representation of the idele class
group related to primes and zeros of L-functions},
Duke Math. J.~127 (2005), 519--595, author version
\texttt{math/0311468v3}, Theorem~5.8 (pp.~43--44) and
Theorem~5.11 (p.~47). A sectorwise topological comparison
with that construction requires its precise image
criterion and is not inferred merely from equal spectra.


\subsection*{MT-M22. The fixed Hecke sector of Meyer's adelic zero-space}

For every fixed primitive unitary Hecke \(L(s,\eta)\), the
canonical quotient \(Q_\eta\) of MT-M18 is isomorphic to the
corresponding fixed compact-character sector of Meyer's
idele-class zero-space.  The isomorphism is unconditional
and compatible with convolution, including all zero
multiplicities.  The sector is first given its quotient
Fr\'echet topology from \(H_-^\chi\); the proof also
identifies the quotient bornology by an explicit splitting.

We fix the character convention explicitly.  Write
\[
 L(s,\eta)=\prod_{\mathfrak p}
       (1-c_{\mathfrak p}(N\mathfrak p)^{-s})^{-1},
\]
where \(c_{\mathfrak p}=0\) at a ramified place.  Let
\(\beta\) be the associated unitary idele-class character
in Meyer's convention: for his \(S\)-prime element
\(p_v\in C_S\), at every unramified \(v\notin S\),
\[
                      \beta(p_v)^{-1}=c_{\mathfrak p_v}.
                                                               \tag{M22.1}
\]
Thus the component is parameterized by
\(\omega_s=\beta|\cdot|^s\).
This convention is dictated by the inverse in Meyer's
Euler factor, rather than by a choice of the name
\(\eta\) or \(\eta^{-1}\) for its idele realization.
For example, \(L(s,\eta)=\zeta_K(s+i\tau)\) corresponds
to \(\beta=|\cdot|^{i\tau}\).
If an idele character has instead been labeled by
its direct values \(c_{\mathfrak p_v}\) on Meyer's
\(p_v\), it must be inverted to obtain \(\beta\).

Let \(C=C_K\), \(C^1=\ker|\cdot|\), and choose a
continuous smooth splitting \(a:\mathbb R\to C\)
with \(|a(t)|=e^t\).
Normalize Haar measure on \(C^1\) to one, so
\(d^\times x=dc\,dt\) in \(x=ca(t)\).
Put \(\chi=\beta|_{C^1}\).
Use Meyer's spaces and embeddings
\[
 H_-=\mathcal S(C)_{\mathbb R},\qquad
 H_+=\mathcal S(\mathbb A_K)/K^\times,\qquad
 H_-^0=H_-/(H_+\cap H_-),
\]
where the intersection uses his embeddings into
\(\mathcal S(C)_{]1,\infty[}\oplus
  \mathcal S(C)_{]-\infty,0[}\).
Superscript \(\chi\) denotes the subspace on which
left translation by \(c\in C^1\) acts by \(\chi(c)\).
Then
\[
 (\mathcal L_\beta f)(x)=\beta(x)^{-1}f(\log|x|)
\]
induces an isomorphism
\[
              Q_\eta \ \simeq\ (H_-^0)^\chi.                  \tag{M22.2}
\]
It is topological for the specified quotient topology.
Norm translation by \(t\) corresponds to idele
translation by \(a(t)\), with the explicit unitary
factor \(\beta(a(t))\).

The primary source is Ralf Meyer, \emph{On a representation
of the idele class group related to primes and zeros of
\(L\)-functions}, Duke Math. J.~127 (2005), 519--595,
author version \texttt{math/0311468v3}.
The exact input is the weighted spaces in Definition 4.1
and (22), author pp.~24--25; the embeddings in Lemma 5.4,
pp.~40--41; and the image criterion in \S5.7, pp.~46--47,
following Lemma 5.10.  Theorem 5.1, pp.~38--39, supplies
the finite-level description used there.
Theorems 5.8, pp.~43--44, and 5.11, p.~47, are the
classical summability and spectral-multiplicity statements.
The comparison below uses the image criterion,
rather than inferring an isomorphism from equal spectra.

\begin{proof}
\emph{The exact test-function sector.}
A norm splitting exists by varying a positive component
at one archimedean place with the appropriate local
norm exponent.  In the resulting product
\(C=C^1\times a(\mathbb R)\), every \(\chi\)-homogeneous
function has the form \(\chi(c)^{-1}u(t)\).
Since \(\beta\) is unitary,
\(\beta(a(t))=e^{i\tau t}\) for a fixed real \(\tau\).
Multiplication by this fixed modulation preserves
every exponential-Schwartz seminorm, in both
directions, by the Leibniz formula.

Consequently \(\mathcal L_\beta\) is a topological
isomorphism from \(\mathcal E\) onto \(H_-^\chi\).
Indeed the norm-direction invariant derivative becomes
\(d/dt\); derivatives in the compact direction act
on its one fixed smooth character by constants.
All norm weights are \(e^{\alpha t}\), and the
requirement for every real \(\alpha\) is equivalent
to the seminorms \(e^{N|t|}|f^{(j)}(t)|\).
Polynomial factors in \(t\) are absorbed by increasing
\(N\).  The character is trivial on a suitable
compact subgroup with elementary compact quotient,
so these functions occupy one fixed Bruhat level.
There is no limit over varying compact characters.

Haar normalization as in Meyer's \S3.1, p.~16,
gives exactly
\[
 \widehat{\mathcal L_\beta f}(\omega_s)
       =\int_C\beta(x)^{-1}f(\log|x|)
                      \beta(x)|x|^s\,d^\times x
       =\mathcal Mf(s).                                        \tag{M22.3}
\]
The transform is zero on all other compact-character
components by orthogonality on \(C^1\).
A direct convolution calculation, using
\(\beta(y)^{-1}\beta(y^{-1}x)^{-1}=\beta(x)^{-1}\),
shows
\[
  (\mathcal L_\beta f)*(\mathcal L_\beta g)
                         =\mathcal L_\beta(f*g).
\]
Also
\[
 \lambda_b\mathcal L_\beta f
   =\beta(b)\mathcal L_\beta
         \bigl(f(\,\cdot-\log|b|)\bigr),\qquad b\in C.            \tag{M22.4}
\]
These identities specify the representation convention.

\emph{Meyer's image condition on this component.}
Choose once and for all a sufficiently large finite
set \(S\), containing the archimedean places, all
ramification of \(\beta\), and the places needed for
his normalized additive character and differential
idele.  His multiplier \(\mathcal L_S\) has transform
\[
 \widehat{\mathcal L_S}(\omega)
       =\prod_{v\notin S}(1-\omega(p_v)^{-1})^{-1}.
\]
By (M22.1), on \(\omega_s\) this is the partial finite
Hecke function
\[
 L_\eta^S(s)
       =L(s,\eta)D_{\eta,S}(s),\qquad
 D_{\eta,S}(s)=
        \prod_{v\in S_{\mathrm{fin}}}
             (1-c_{\mathfrak p_v}(N\mathfrak p_v)^{-s}).         \tag{M22.5}
\]
A ramified factor in \(D_{\eta,S}\) is one.
The reflected quasi-character is
\[
 \omega_s^*=|\cdot|\omega_s^{-1}
                    =\beta^{-1}|\cdot|^{1-s},
\]
and therefore
\[
              \widehat{\mathcal L_S}(\omega_s^*)
                       =L_{\bar\eta}^S(1-s).                   \tag{M22.6}
\]
Here \(\bar\eta\) denotes the dual unitary Hecke
character, with conjugate Euler coefficients.

For an entire strip-Schwartz function \(F(s)\)
on this component, the equivalent reduced
image criterion printed on Meyer p.~47 is
\[
 \begin{array}{ll}
 F(\alpha+it)/L_\eta^S(\alpha+it)
        \text{ is Schwartz in }t,&0<\alpha\le1,\\[2mm]
 F(it)/L_{\bar\eta}^S(1-it)
        \text{ is Schwartz in }t.&
 \end{array}                                                    \tag{M22.7}
\]
All apparent zero singularities must be removable.
The source's remaining conditions for
\(\alpha>1\) and \(\alpha<0\) follow from \(F\)
being the transform of \(H_-\).
Zero extension to other compact-character
components causes no additional Schwartz condition.
Thus (M22.7) characterizes membership of
\(\mathcal L_\beta f\) in \(\Sigma(H_+)\),
where \(F=\mathcal Mf\).

This is also the relevant intersection in
the quotient, rather than only its first-coordinate
image.  Meyer's Poisson formula gives
\[
 \Sigma u-J\Sigma\mathfrak F^*u
                         =c_0-c_1|\cdot|^{-1}.
\]
The second term lies in
\(\mathcal S(C)_{]-\infty,0[}\).
If \(\Sigma u\in H_-\), so does the right side.
Its decay as \(|x|\to0\) forces \(c_0=c_1=0\):
these distinct powers cannot decay after every
negative weight.  Hence
\(i_+u=i_-(\Sigma u)\).
Conversely, an element of the intersection has
its first coordinate in \(\Sigma(H_+)\cap H_-\).
Compact-character projection commutes with the
maps.  Therefore (M22.7) characterizes
\((H_+\cap H_-)^\chi\).

\emph{Removing the image criterion's growth gate.}
If (M22.7) holds, every zero of \(\Xi_\eta\) in
\(0<\operatorname{Re}s\le1\) is a zero of \(F\)
with at least its full multiplicity, by its first
condition.  The omitted finite Euler factors have
no zeros in that right half-plane, and the gamma
factors have neither zeros nor poles there.
The possible pure-norm pole is cancelled in
\(\Xi_\eta\) and is not a zero of it.
If a nontrivial zero is on \(\operatorname{Re}s=0\),
the functional equation identifies its multiplicity
with that of \(L_{\bar\eta}(1-s)\); the second
condition gives the same conclusion.
Thus membership implies all the zero-jet
vanishings of MT-M18.

Conversely suppose \(F=\mathcal Mf\) vanishes
on the full divisor of \(\Xi_\eta\).
Set \(G=F/\Xi_\eta\), which is entire.
The rectangle proof of (M18.14) gives,
for every \(A,N>0\), the sharpened estimate
\[
 \sup_{|\sigma|\le A}|G(\sigma+it)|
       \le C_{A,N,f}e^{\kappa|t|}(2+|t|)^{-N},
 \qquad \kappa=\pi[K:\mathbb Q]/4.                              \tag{M22.8}
\]
Indeed the boundary lower bound for \(\Xi_\eta\)
costs \(e^{\kappa|n|}(2+|n|)^{B_A}\).
Use the decay power \(N+B_A\), rounded up and
increased if needed, for \(F\) on that boundary.
The maximum principle and bounded block heights
give (M22.8), with no zero-spacing hypothesis.

Write the standard completion as
\[
 \Xi_\eta(s)=P_\eta(s)q_\eta^{s/2}
                  \Gamma_\eta(s)L(s,\eta).
\]
On each closed strip in \(\operatorname{Re}s>0\),
the unitary archimedean gamma product is
holomorphic and bounded above by
\(e^{-\kappa|t|}\) times a fixed power
of \(2+|t|\), by Stirling's formula.
Consequently
\[
        F(s)/L(s,\eta)
          =G(s)P_\eta(s)q_\eta^{s/2}\Gamma_\eta(s)
\]
is holomorphic and rapidly decreasing there,
by (M22.8).  The reciprocal of a finite-\(L\)
pole is a zero, so pure-norm poles cause
no exception.

The completed functional equation gives
\[
 \Xi_\eta(s)
   =\epsilon_\eta P_\eta(s)q_\eta^{(1-s)/2}
            \Gamma_{\bar\eta}(1-s)L(1-s,\bar\eta),
                    \qquad|\epsilon_\eta|=1.
\]
It follows that \(F(s)/L(1-s,\bar\eta)\)
is holomorphic and rapidly decreasing
on each closed strip in
\(\operatorname{Re}s<1\).
The reflected gamma arguments have positive
real part there.  Thus neither half-plane
division asks \(F\) to cancel trivial
finite-\(L\) zeros outside the appropriate half.

For \(\delta>0\), every nontrivial omitted
Euler factor satisfies
\[
 |1-c_{\mathfrak p}(N\mathfrak p)^{-s}|
                         \ge1-(N\mathfrak p)^{-\delta}
       \quad(\operatorname{Re}s\ge\delta).
\]
The ramified factors are one.
Thus \(D_{\eta,S}^{-1}\) is bounded on each
such right half-plane, as is
\(D_{\bar\eta,S}(1-s)^{-1}\) when
\(\operatorname{Re}s\le1-\delta\).
Their derivatives are bounded on smaller
fixed strips.  Therefore
\[
 \begin{split}
 F/L_\eta^S&\text{ is strip-Schwartz on }\operatorname{Re}s>0,\\
 F/L_{\bar\eta}^S(1-s)&\text{ is strip-Schwartz on }
                                      \operatorname{Re}s<1.
 \end{split}                                                    \tag{M22.9}
\]
Cauchy's formula on slightly larger strips
gives every vertical derivative estimate.
These statements imply (M22.7).
We have proved
\[
 \mathcal L_\beta^{-1}\bigl((H_+\cap H_-)^\chi\bigr)
   =\{f:(\mathcal Mf)^{(j)}(\rho)=0,\
          \Xi_\eta(\rho)=0,\ j<m_\rho\}
   =\overline{\psi_\eta*\mathcal E},                             \tag{M22.10}
\]
where the last equality is MT-M18.

\emph{The precise quotient topology and bounded lifting.}
Give \(H_-^\chi\simeq\mathcal E\) its
just-identified Fr\'echet topology and its
bounded-set bornology.  At this fixed character
these are precisely the source's weighted
Schwartz structure: it lies at one Bruhat
level, and the countable family of real
weight endpoints in (22) generates its
seminorms and bounded sets.
Here bounded and precompact sets coincide,
so the convention of using the precompact
bornology gives the same result.  Indeed,
for a bounded sequence in \(\mathcal E\),
the bounds on one additional derivative
give equicontinuity of every derivative
on each compact interval.  Diagonal
Arzel\`a--Ascoli gives a subsequence
converging with every derivative on every
compact interval.  For a prescribed
\(p_{N,j}\), the uniform \(p_{N+1,j}\)
bound makes the weighted tails outside
\([-R,R]\) at most \(Ce^{-R}\).
The subsequence is therefore Cauchy in
every defining seminorm and converges
in the complete space \(\mathcal E\).
This proves relative compactness of every
bounded set in its metrizable topology;
the converse is immediate.
The subspace in (M22.10) is closed.
Hence (M22.10) immediately gives a
topological isomorphism from \(Q_\eta\)
onto the sector quotient furnished with
the quotient topology from \(H_-^\chi\).

Here is why the original quotient bornology
agrees as well.  The maps \(i_+\) and \(i_-\)
are bornological embeddings by the cited
source.  Therefore the bornology of their
intersection equals the subspace bornology
inherited from \(H_-^\chi\): a subset of the
intersection is bounded in either embedded
space exactly when it is bounded in the
common ambient space.  No topology on
\(H_+^\chi\), or on its quotient, is being
assumed in this argument.

Compact-character averaging is a bounded
projection preserving \(H_+\) and \(H_-\).
At each finite level it is ordinary integration
over a compact group against its smooth
character, and these projections agree under
the level inclusions.
It preserves their intersection and splits
both spaces into its image and kernel.
The algebraic quotient and its quotient
bornology therefore split accordingly.
In particular \((H_-^0)^\chi\) is the
bornological quotient
\(H_-^\chi/(H_+\cap H_-)^\chi\);
representatives can be projected before
passing to the quotient.

Finally the continuous linear section
\(S:s(J)\to\mathcal E\) in MT-M18 supplies
an explicit lifting map for this quotient,
after transport by \(\mathcal L_\beta\).
Every bounded set in the quotient
Fr\'echet space lifts to a bounded set
through this section.  Conversely the
quotient map is continuous, so it takes
bounded sets to bounded sets.
The same two assertions hold for
precompact sets, using continuity of
the section and quotient map.
The quotient itself has bounded sets
precompact, since its continuous section
maps them to bounded, hence precompact,
subsets of \(\mathcal E\), and its
quotient map sends these back to the
original sets.
Thus the source's quotient bornology
is exactly the bornology of the specified
Fr\'echet quotient.  This is an explicit
splitting argument, not an inference
from equality of spectra.
Together with (M22.3)--(M22.4) and
multiplicativity, it proves all parts
of (M22.2).
\end{proof}

The stable blocks and nuclear traces in
MT-M18--MT-M19 therefore apply to this
independently defined adelic summation
quotient in each fixed Hecke sector.
All constants may depend on the character
and its infinity type; no uniform
decomposition over all characters is
asserted.  The trace positivity and
faithful-operator distinctions of
MT-M19--MT-M20 retain their exact GRH
and multiplicity hypotheses.
This identifies Meyer's specified
quotient, rather than an unspecified
quotient in Masullo's draft.
Its geometric location is the norm
direction of the idele-class action,
Tate's Poisson summation, and cote~1's
nuclear-trace framework as carried in
SA23 geo-H-lidskii/geo-H-genus.
The primary checks cover the named
normalization, image, and embedding
passages; no complete-paper audit
is claimed.


\subsection*{MT-M23. Every nonabelian Galois top field obstructs weak Mertens and faithful positivity}

Let $L/K$ be a nonabelian Galois extension of number fields. The
classical theorem of Hu--Kaneko--Martin--Schildkraut (HKMS),
\emph{Order of zeros of Dedekind zeta functions}, Theorem1.1,
author version2107.03269v2 p.2, proves unconditionally that
$\zeta_L$ has infinitely many nontrivial zeros of multiplicity
at least two. The complete first proof is \S3 p.5; a second
proof occupies \S4 pp.5--6. Their Theorem1.2, p.2 with proof
\S5 pp.6--8, gives infinitely many zeros of multiplicity at
least three whenever $\operatorname{Gal}(L/K)$ has an irreducible
representation of degree at least three. These are cited classical
results; no new multiplicity theorem is claimed here.

Set
\[
 M_L(x)=\sum_{N\mathfrak a\le x}\mu_L(\mathfrak a),\qquad
 I_L(X)=\int_1^X|M_L(x)/x|^2\,dx.
\]
The following consequences hold without either adopted GRH input:
\[
 \limsup_{X\to\infty}\frac{I_L(X)}{(\log X)^3}>0,
 \qquad
 \limsup_{X\to\infty}\frac{I_L(X)}{\log X}=+\infty.
 \tag{MT-M23a}
\]
If the group has an irreducible representation of degree at least
three, the exponent $3$ in the first assertion can be replaced by
$5$. The canonical full-jet quotient of $\zeta_L$ has no faithful
continuous Hilbert completion on which every convolution extends
boundedly and every $C_{h*h^\sharp}$ is a positive operator.
Nevertheless its nuclear trace construction and adapted absolute
basis exist unconditionally by MT-M18--MT-M19, and $M_H$ makes
the traces of these adjoint-square convolutions nonnegative.

\paragraph{Complete proof of the consequences.}
Choose a multiple zero supplied by HKMS, of actual multiplicity
$m\ge2$, or $m\ge3$ in the stronger case. If $\zeta_L$ has any
off-critical zero, the functional equation and conjugation give
a zero with real part greater than $1/2$. The last assertion of
NT-L16 then gives infinite limsup after division by every fixed
nonnegative power of $\log X$; all displayed conclusions follow.
Otherwise all nontrivial zeros are critical, and our chosen zero
$\rho$ satisfies the exact NT-L16 bound
\[
 \limsup_{X\to\infty}\frac{I_L(X)}{(\log X)^{2m-1}}
 \ge
 \frac{1}{(2m-1)((m-1)!)^2
       |\rho\,\zeta_L^{(m)}(\rho)/m!|^2}>0.
\]
When $2m-1$ equals the claimed exponent this is the first assertion.
When it is greater, take an unbounded sequence on which the
displayed ratio is bounded below by a fixed positive constant;
multiplying by the remaining positive power of $\log X$ gives
infinite limsup. The same argument with exponent one proves
the second assertion. This dichotomy uses no unproved zero
location. Finally MT-M20 proves that a faithful positive Hilbert
realization with the stated properties forces every nontrivial
zero to be simple. The actual multiple zeros contradict that
necessary condition. MT-M18--MT-M19 give the separate basis and
trace statements, and their GRH equivalence gives the conditional
trace positivity under $M_H$.

For the programme this is a finite-monodromy Artin-motive
obstruction to discarding nilpotent spectral blocks. It sharpens
the Kummer examples NT-L17 to all nonabelian Galois top fields.
The cote1/SA23 \texttt{geo-H-lidskii} trace route must retain the
difference between positivity of a scalar trace and positivity
of every operator in a faithful full-multiplicity Hilbert model.
HKMS's general multiplicity prediction for representations of
degree at least four remains their Conjecture6.2, not a theorem
used here. A non-Galois field is covered only when it is itself
the top field of a nonabelian Galois extension of some subfield.


\subsection*{MT-ART09. A smaller relative-zero gate for solvable Artin holomorphy}

Let $L/K$ be a finite solvable Galois extension, with group $G$,
and let $s_0\ne1$. Write
\[
 q(s_0)=\operatorname{ord}_{s=s_0}
             \frac{\zeta_L(s)}{\zeta_K(s)}\ge0.
\]
The nonnegativity follows from the classical Aramata--Brauer
theorem. When $|G|$ has at least two distinct prime divisors,
let $p_2$ be the second smallest and put
\[
                 B(G)=\max\{2,p_2-2\}.
\]
Then every Artin $L(s,\chi,L/K)$ is holomorphic at $s_0$ if
$q(s_0)\le B(G)$. This is the combination of Peng-Jie Wong,
\emph{Zeros of Dedekind zeta functions and holomorphy of Artin
$L$-functions}, Theorems1.3 and1.4, printed p.254, with their
complete proofs on pp.259 and260. The solvability hypothesis
is essential to the statements as cited here. In a $p$-group
all irreducible characters are monomial, so Artin holomorphy
already follows globally from Hecke holomorphy.

For $4\nmid |G|$ the sufficient bound $q(s_0)\le3$ holds,
without adding a solvability assumption: Wong Corollary5.3,
printed pp.259--260, includes the group-theoretic implication
that these groups are solvable. His Corollary5.1, p.259, is the
odd-degree special case. One may therefore replace $B(G)$ by
$\max\{B(G),3\}$ when applicable. The stronger Sylow-structure
variants in his Theorems5.2 and5.4 are separately located there;
no unstated prime-factor hypothesis from them is being used.

\paragraph{Complete transfer to the adopted Hecke hypothesis.}
Under $M_H$, a possible pole of an Artin function in
$0<\Re s<1$ must lie on $\Re s=1/2$: Brauer induction writes
that function as a finite quotient of Hecke functions, whose
possible nontrivial denominator zeros are all on that line.
This is the MT-ART01 confinement, and does not assert Artin
holomorphy. The contrapositive of the just-cited local theorem
now proves that a possible pole must satisfy
\[
 \Re s_0=\tfrac12,\qquad q(s_0)\ge B(G)+1.
 \tag{MT-ART09a}
\]
If $4\nmid |G|$, use the strengthened bound at least four.
Thus, if the relative zeta function has order at most $B(G)$
at every point of the critical line, all Artin functions are
holomorphic throughout the open critical strip. This last
zero-order condition is an extra input, not supplied by $M_H$.
For an unconditional local conclusion at a particular $s_0$,
only the corresponding relative-order inequality is needed.

This is a strictly finer criterion than measuring the order
of $\zeta_L$ alone: by the definition of a meromorphic order,
\[
 q(s_0)=\operatorname{ord}_{s_0}\zeta_L
                 -\operatorname{ord}_{s_0}\zeta_K.
\]
The base contribution is removed exactly, including when the
base has a multiple zero. It is not necessary to posit any
noncoincidence of zeros. The argument is a complete input
transfer to the cited classical local theorem, not a new proof
of that theorem or of general Artin holomorphy.

Geometrically, $\zeta_L/\zeta_K$ is the $L$-function of the
augmentation Artin motive attached to the finite etale cover
after deleting ramified primes, with the standard inertia
factors restored. The relative virtual permutation class is
$[\mathbb Q[G]]-[\mathbb Q]$ in its representation $K_0$.
This belongs to the cote1 trace and SGA5 finite-correspondence
framework already fixed in AG-C1, and sharpens the
solvable part of Q-ARTIN-POS beyond MT-ART02. It does not turn a
virtual-character inequality into coefficientwise positivity.


\subsection*{MT-M24. The specified Masullo smooth quotient is the Meyer quotient}
This is a comparison of the actual test spaces and summation maps. It does
not assume Masullo's global converse theorem, his spectral basis claims,
or $M$. Fix the primitive unitary Hecke function and the convention
$\beta(\mathfrak p_v)^{-1}=c_{\mathfrak p_v}$ from MT-M22; put
$C=C_K$, $C^1=\ker|\cdot|$, and $\chi=\beta|_{C^1}$.
Use the norm splitting and volume-one measure on $C^1$ chosen there.
Write $\mathcal E$ for the exponential Schwartz space of MT-M18,
$\mathcal I_\eta=\overline{\psi_\eta*\mathcal E}$ for its closed
completed-Hecke ideal, and $Q_\eta=\mathcal E/\mathcal I_\eta$.
For a real number $a$, let $M_a u(g)=|g|^a u(g)$, and define
\[
 \Theta_\chi=\Phi_{K,\chi}\bigl(\mathcal S_0(\mathbf A_K)\bigr),
 \qquad
 Q^{\mathrm{Mas}}_{\chi,a}
 =\mathcal S_{\mathrm E}(C)(\chi)\big/
       \overline{M_a\Theta_\chi}^{\,\mathcal S_{\mathrm E}(C)(\chi)}.
 \tag{M24.1}
\]
Here the symbols on the right have their literal meanings in Masullo:
\[
 \mathcal S_0(\mathbf A_K)
  =\{f\in\mathcal S(\mathbf A_K):f(0)=\widehat f(0)=0\},
 \qquad
 \Phi_K(f)(g)=\sum_{\xi\in K^\times}f(\xi g),
 \tag{M24.2}
\]
and $\Phi_{K,\chi}$ is compact-character averaging of $\Phi_K$.
The closure in (M24.1) does not enlarge $M_a\Theta_\chi$.
There is an explicit topological isomorphism
\[
 Q_\eta\xrightarrow{\ \sim\ }Q^{\mathrm{Mas}}_{\chi,a},
 \qquad
 [f]\longmapsto
 \bigl[g\longmapsto |g|^a\beta(g)^{-1}f(\log|g|)\bigr].
 \tag{M24.3}
\]
For $a=0$ this identifies the quotient in Masullo (7.1.4.7),
original PDF p.~327. For $a=1/2$ it identifies the centered quotient
in (1.8.3.1), original PDF p.~45. These two source displays use
different twists; (M24.3) keeps the twist explicit.

\begin{proof}
First verify the test-space topology rather than just its elements.
Masullo Definition D.3.3.1, original PDF p.~804, requires
$|\cdot|^b u\in\mathcal S(C)$ for every real $b$.
Definition D.3.3.3, p.~805, gives each fixed compact-invariance
level the projective topology of all these weighted Schwartz spaces,
and gives their union the locally convex inductive topology.
On the $\chi$-sector the map
\[
 \mathcal L_\beta f(g)=\beta(g)^{-1}f(\log|g|)
 \tag{M24.4}
\]
is a linear bijection from $\mathcal E$ onto that sector.
Indeed, the $\chi$-eigencondition for left translation says
$u(cg)=\chi(c)^{-1}u(g)$ for $c\in C^1$; its restriction to
the norm splitting therefore determines $u$. The unitary value
of $\beta$ on that splitting is $e^{i\tau x}$ for a fixed real
$\tau$, so removing it only changes derivatives by fixed finite
linear combinations.

Every function in this sector is invariant under $\ker\chi$.
The compact quotient $C^1/\ker\chi$ is a closed subgroup of the
circle, hence is either finite or a circle, and this is one permitted
elementary compact-invariance level. At this level the weighted
Schwartz seminorms, after $x=\log|g|$, are equivalent to
\[
 p_{N,j}(f)=\sup_x e^{N|x|}|f^{(j)}(x)|,
 \qquad N,j\ge0.
 \tag{M24.5}
\]
To see the equivalence directly, the weights $b=N$ and $b=-N$
bound (M24.5); conversely every polynomial factor in $x$ occurring
in a weighted Schwartz seminorm is absorbed by one larger
exponential weight. Derivatives on the compact factor multiply
this fixed character by constants. Inclusion of this level in the
inductive limit is continuous. The inverse of (M24.4) on the
sector is also continuous: compact-character averaging followed
by restriction to the norm splitting is continuous into
$\mathcal E$ on each elementary level, by these same seminorm
estimates. It is therefore continuous on the inductive limit.
This establishes the stated sector topology directly, without
using a topological claim about a different quotient.

Next identify the denominator as a set of actual functions.
Use Meyer's common ambient space and embeddings from MT-M22:
\[
 i_+(v)=(\Sigma v,J\Sigma\mathcal F_*v),
 \qquad i_-(u)=(u,u),
 \qquad H_+=\mathcal S(\mathbf A_K)/K^\times,
 \quad H_-=\mathcal S_{\mathrm E}(C).
\]
The summation maps agree term by term:
$\Sigma f(g)=\Phi_K(f)(g)$ by (M24.2).
Their Poisson identity is
\[
 J\Sigma\widehat f-\Sigma f
      =f(0)-|\cdot|^{-1}\widehat f(0).
 \tag{M24.6}
\]
For $f\in\mathcal S_0(\mathbf A_K)$ the two entries agree.
The first entry has all the right-half-plane Schwartz weights,
and the second all the left-half-plane weights, by the defining
Meyer embeddings. Their common value has all real weights:
between any two admitted weights the weighted derivative is
bounded by the two endpoint weighted derivatives, and polynomial
Schwartz factors are treated in the same way. Thus this common
value belongs to $H_-$, and
$\Phi_K(\mathcal S_0)\subset H_+\cap H_-$.

Conversely, if $(u,u)=i_+(v)$ lies in this intersection, choose
an adelic Schwartz representative $f$ of the coinvariant class
$v$. Equation (M24.6) then vanishes identically. The functions
$1$ and $|\cdot|^{-1}$ are linearly independent because the norm
map is onto $\mathbf R_{>0}$, so $f(0)=\widehat f(0)=0$.
Consequently
\[
       \Phi_K(\mathcal S_0(\mathbf A_K))=H_+\cap H_-.
 \tag{M24.7}
\]
This argument uses an actual representative of the coinvariant
class, not a claim that every completion class has an adelic
representative. $H_+$ here is Meyer's quotient of the adelic
Schwartz space.

Compact-character averaging preserves both embedded spaces
and their intersection. On functions, the volume-one projector is
\[
 P_\chi u(g)=\int_{C^1}u(cg)\chi(c)\,dc.
\]
Masullo (2.2.5.2), p.~95, is precisely this formula applied to
$\Phi_K f$, except that his chosen compact Haar measure can have
volume $c_K\ne1$. This multiplies the operator by the fixed
nonzero scalar $c_K$ and leaves its image unchanged. Hence
\[
       \Theta_\chi=(H_+\cap H_-)^\chi.
 \tag{M24.8}
\]
One can also verify surjectivity in (M24.8) directly: an element
of the intersection in the $\chi$-sector already has an
$\mathcal S_0$ representative by (M24.7), and applying
$\Phi_{K,\chi}$ to that representative returns $c_K$ times the
element. No choice of a compact subgroup lift in the idele group
is required.

MT-M22, whose image criterion and topology were proved using
Meyer's specified maps, now gives
\[
 \mathcal L_\beta^{-1}\Theta_\chi
 =\{f\in\mathcal E:(\mathcal Mf)^{(j)}(\rho)=0
         \text{ for all }0\le j<m_\rho\}
 =\mathcal I_\eta.
 \tag{M24.9}
\]
The middle space is closed, since each indicated evaluation is
continuous. Thus $\Theta_\chi$ is closed in the actual sector
topology just checked. Multiplication by $|\cdot|^a$ is an
automorphism of that topology, with inverse multiplication by
$|\cdot|^{-a}$; this follows either from (M24.5), using an
integer larger than $|a|$, or directly by translating the family
of all real weights. This proves the closure assertion and the
continuous two-sided quotient isomorphism (M24.3).
\end{proof}

The action is also explicit. For $b\in C$,
\[
 \lambda_b M_a\mathcal L_\beta f
  =\beta(b)|b|^{-a}M_a\mathcal L_\beta
       \bigl(f(\,\cdot-\log|b|)\bigr).
 \tag{M24.10}
\]
For an idele-class convolution test $k$, set
\[
 k_\beta(y)=\int_{C^1}k(c\tau(e^y))\beta(c\tau(e^y))\,dc.
\]
Changing variables in the Haar integral gives
\[
 C_k M_a\mathcal L_\beta f
   =M_a\mathcal L_\beta\bigl((e^{-a\cdot}k_\beta)*f\bigr).
 \tag{M24.11}
\]
In particular the spectral multiplier is evaluated at $\rho-a$,
with the compact-character factor $\beta$ retained. At $a=0$ this
is precisely the MT-M22 module comparison; at $a=1/2$ it is the
centered version, not an assertion that the same convolution
operator is unchanged by the twist.
An alternative Haar normalization rescales the scalar convolution
integral by the same fixed positive constant; it changes neither
the denominator nor the quotient topology.

\paragraph{Exact source scope and the remaining completion question.}
The original Masullo PDF pages used here are: the theta-sum
definition (2.2.3.1), pp.~81--82, and its underlying sum (2.2.2.1),
p.~77; the polar hyperplanes (2.2.4.23), p.~91;
character averaging (2.2.5.2), p.~95; the actual smooth quotient
(7.1.4.7), p.~327; its centered introductory version (1.8.3.1),
p.~45; and Definitions D.3.3.1/D.3.3.3, pp.~804--805.
The original p.~327 quotient formula was rendered and read;
the other defining formulas were checked against extraction
from the original PDF. The Meyer inputs are his summation map
and Lemma~5.3, pp.~39--40; the embeddings in Lemma~5.4,
pp.~40--41; and Poisson formula (31), p.~41, with the complete
provided arguments already read for MT-M22. No new proof of
the full Masullo draft is claimed.

Masullo Definition~10.1.1.1, p.~434, instead closes
$M_a\Theta_\chi$ in a weighted Hilbert norm, and
(10.2.1.1)--(10.2.1.2), p.~439, and (10.4.1.3), p.~456,
form the resulting Hilbert quotients. Those are different
topological completions of the identified smooth quotient.
Their kernel, retained jets and trace-class properties must be
proved in that norm; MT-M25 and MT-M26 address those maps.
The quotient identification itself is unconditional. Any
subsequent use of universal Hecke GRH is still the conditional
input $M$.

The geometric crosswalk is now an actual Tate--Poisson
summation quotient, rather than an unspecified space sharing
its zero multiplicities. It advances the inherited cote~1 /
SA23 geo-H-lidskii and geo-H-genus interfaces at the precise
test-space and closed-denominator level. It makes no claim
to construct an $\ell$-adic or motivic realization.

\subsection*{MT-M25. Exact line-completion kernel and the strict jet cutoff}
Fix the Hecke function, $\mathcal E$, $\Xi=\Xi_\eta$,
$\mathcal I=\mathcal I_\eta$ and $Q=\mathcal E/\mathcal I$
of MT-M18. For $s\ge0$ and $a\in\mathbf R$, put
\[
 H_s=L^2\bigl(\mathbf R,(1+x^2)^{s/2}\,dx\bigr),
 \qquad T_af(x)=e^{ax}f(x),
 \qquad J_{a,s}=\overline{T_a\mathcal I}^{\,H_s},
 \quad Z_{a,s}=H_s/J_{a,s}.
 \tag{M25.1}
\]
These are exactly the fixed-character versions of Masullo's
line norm (9.1.2.3), original PDF p.~399, his closure
(10.1.1.3), p.~434 with $I=\{0\}$, and his quotient
(10.2.1.2), p.~439, by MT-M24. The fixed unitary norm
modulation has absolute value one and is harmless for this norm.
Let $m_\rho$ be the multiplicity of a zero of $\Xi$, and set
\[
       q_s=\max\{0,\lceil(s-1)/2\rceil\}.
 \tag{M25.2}
\]
Then
\[
 J_{a,s}=
 \bigcap_{\substack{\rho=a+i\gamma:\Xi(\rho)=0\\
                    0\le k<\min(m_\rho,q_s)}}
 \ker\left[u\longmapsto\int_{\mathbf R}
                  u(x)x^k e^{i\gamma x}\,dx\right].
 \tag{M25.3}
\]
Every functional displayed here is continuous on $H_s$.
The natural map
\[
 q_{a,s}:Q\longrightarrow Z_{a,s},\qquad[f]\longmapsto[T_af]
 \tag{M25.4}
\]
is continuous with dense image, and its kernel consists exactly
of the classes whose jets $F^{(k)}(\rho)$, $F=\mathcal Mf$,
vanish at all the indices in (M25.3). Thus the line completion
retains only the zeros on $\Re\rho=a$ and only their derivatives
with the \emph{strict} inequality
\[
                         2k<s-1.
 \tag{M25.5}
\]
It is faithful on the full smooth quotient if and only if all
zeros lie on that line and $m_\rho\le q_s$ for every zero.
No GRH or simplicity is assumed in these assertions.

\begin{proof}
The continuous bilinear dual of $H_s$ can be written as
$H_{-s}$ under $\ell_g(u)=\int u(x)g(x)\,dx$.
This is just the Riesz representation theorem after multiplying
by $(1+x^2)^{s/4}$; complex conjugation can be absorbed into
the representing $g$.
For a polynomial $P$ of degree $k$ with nonzero leading
coefficient,
\[
 P(x)e^{i\gamma x}\in H_{-s}
 \quad\Longleftrightarrow\quad
 \int_{|x|\ge1}|x|^{2k-s}\,dx<\infty
 \quad\Longleftrightarrow\quad 2k<s-1.
 \tag{M25.6}
\]
This proves continuity of the displayed functionals and the
necessity of their strict numerical threshold. It also shows
that the number of admissible nonnegative integers $k$ is
exactly $q_s$, including $0\le s\le1$ when there are none.
Each such functional vanishes on $T_a\mathcal I$ by MT-M18's
full-jet ideal equality, hence on its Hilbert closure.

We prove there are no other annihilators. The elementary
inequality
\[
 1+(x+y)^2\le2(1+x^2)(1+y^2)
\]
implies, for either sign of $s$,
\[
 \|u(\,\cdot-y)\|_{H_{\pm s}}
       \le 2^{s/4}(1+y^2)^{s/4}\|u\|_{H_{\pm s}}.
 \tag{M25.7}
\]
For the negative sign use the same inequality with $x$ and
$x-y$ interchanged before taking reciprocals. Minkowski's
integral inequality therefore gives
\[
 \|k*u\|_{H_{\pm s}}
 \le2^{s/4}\|u\|_{H_{\pm s}}
       \int|k(y)|(1+y^2)^{s/4}\,dy.
 \tag{M25.8}
\]
In particular convolution by a Schwartz function is bounded
on both spaces. Smooth compactly supported functions are
dense in both weighted $L^2$ spaces: truncate first, and then
approximate on the remaining compact interval, where the
weight and its inverse are bounded. These facts also prove
strong continuity of translation, first on that dense subspace
and then by (M25.7).

Let $g\in H_{-s}$ annihilate $T_a\mathcal I$.
It defines a tempered distribution, since Cauchy--Schwarz
bounds $|\int ug|$ by the continuous Schwartz seminorm
$\|u\|_{H_s}\|g\|_{H_{-s}}$.
Use the Fourier convention
$\mathcal F_+u(t)=\int u(x)e^{itx}\,dx$ and its inverse
$(2\pi)^{-1}\int\phi(t)e^{-itx}\,dt$.
Define the tempered distribution $D$ by
$D(\phi)=\ell_g(\mathcal F_+^{-1}\phi)$.
Since $\psi_a=T_a\psi_\eta\in\mathcal E$ and
$\mathcal F_+\psi_a(t)=\Xi(a+it)$, we obtain
\[
                         \Xi(a+it)D=0.
 \tag{M25.9}
\]
Here is why the identity holds on every Schwartz test,
not just an entire Fourier transform. For $u\in\mathcal E$,
$\psi_a*u\in T_a\mathcal I$ because
$T_a(f_1*f_2)=T_af_1*T_af_2$ and $T_a\mathcal E=\mathcal E$.
Thus $\ell_g(\psi_a*u)=0$. The space $\mathcal E$ contains
$C_c^\infty$ and is dense in the usual Schwartz space;
convolution by $\psi_a$ is continuous there, and
$\ell_g$ is continuous there. This extends the equality to
all Schwartz $u$, which is precisely (M25.9).

The zeros of $t\mapsto\Xi(a+it)$ on the real line form a
locally finite set. Away from them, a compactly supported
smooth test can be divided by this nonvanishing multiplier,
so (M25.9) shows $\operatorname{supp}D$ is contained in that
set. Near a zero $\gamma$ of order $m_\rho$, multiplication
by a smooth invertible function reduces (M25.9) to
$(t-\gamma)^{m_\rho}D=0$. A distribution supported there
with this property depends only on the first $m_\rho$ jets
of a test: subtract its Taylor polynomial of degree
$m_\rho-1$ times a cutoff equal to one near $\gamma$;
the remainder is divisible by $(t-\gamma)^{m_\rho}$ there.
It is therefore a linear combination of
$\delta_\gamma^{(k)}$, $0\le k<m_\rho$.

Choose a smooth frequency cutoff supported near this single
zero and equal to one near it. Multiplying $D$ by that
cutoff corresponds to convolution of $g$ with a Schwartz
kernel, with an inessential reflection dictated by the
chosen Fourier convention. By (M25.8) the resulting
function still belongs to $H_{-s}$. Its inverse transform
is a polynomial times $e^{i\gamma x}$. Formula (M25.6)
therefore excludes every nonzero coefficient of degree
$k$ with $2k\ge s-1$. Locally only the functionals listed
in (M25.3) occur.

Finally these local pieces exhaust the annihilator in its
Hilbert norm. Choose $\vartheta\in C_c^\infty(\mathbf R)$
equal to one near zero and put
$\vartheta_R(t)=\vartheta(t/R)$, $R\ge1$.
The associated convolution kernels have the form
$Rk(Rx)$, possibly reflected, where $k$ is Schwartz and
$\int k=1$. Their weighted $L^1$ norms in (M25.8) are
uniformly bounded, because
\[
 \int R|k(Ry)|(1+y^2)^{s/4}\,dy
       =\int|k(v)|(1+v^2/R^2)^{s/4}\,dv
       \le\int|k(v)|(1+v^2)^{s/4}\,dv.
\]
Translation continuity and dominated convergence in the
same integral show $Rk(R\cdot)*g\to g$ in $H_{-s}$.
The distribution $\vartheta_RD$ is supported on finitely
many zeros, so each such convolution is a finite linear
combination of the permitted polynomial waves. Hence
the full annihilator of $T_a\mathcal I$ is the Hilbert
closure of the span of precisely the functionals in
(M25.3). Taking its annihilator proves (M25.3).

The inclusion $T_a:\mathcal E\to H_s$ is continuous,
since one exponential seminorm with weight greater than
$|a|$ absorbs the polynomial factor in (M25.1). Its range
is dense because $T_a\mathcal E=\mathcal E$ contains
$C_c^\infty$. It vanishes on $\mathcal I$ after passing
to the quotient, and so induces (M25.4), continuously
with dense image. The functional indexed by $(\rho,k)$
evaluates $T_af$ to
\[
       \int e^{ax}f(x)x^ke^{i\gamma x}\,dx
                    =(\mathcal Mf)^{(k)}(\rho).
\]
This gives its exact kernel. If all zero jets are retained,
MT-M18 makes that kernel zero. If any zero jet is omitted,
the finite-block interpolation of MT-M18 realizes a class
with just that jet nonzero and all other zero jets zero.
It is a nonzero class in the kernel of (M25.4). This proves
the stated necessary and sufficient condition for faithfulness.
\end{proof}

\paragraph{A literal boundary correction with an arithmetic witness.}
The original PDF (10.2.1.6), p.~440, displays
\[
 N_s(\rho)=\min\{N(\rho),[(s-1)/2+1]\}.
\]
At $s=3$ this allows the first derivative at any double zero.
But its proposed representing function is
$x e^{i\gamma x}$, whose squared $H_{-3}$ norm contains
\[
              \int_{|x|\ge1}\frac{x^2}{(1+x^2)^{3/2}}\,dx
                         =+\infty.
 \tag{M25.10}
\]
Thus the endpoint inclusion is false as written; the
correct retained multiplicity is
$\min\{m_\rho,q_s\}$. Both the source norm (9.1.2.3),
p.~399, and the displayed cutoff (10.2.1.6), p.~440, were
rendered and visually read from the original PDF.
This is not a hypothetical objection requiring an unknown
double zero of the Riemann zeta function. By NT-L17,
the actual degree-six field
$K=\mathbf Q(\sqrt[3]{2},\sqrt{-3})$ has nontrivial zeros
of $\zeta_K$ of order at least two. Choose such a zero
$\rho$, put $a=\Re\rho$, and use its trivial Hecke
character. Then (M25.10) is a concrete counterexample
to the stated membership at its permitted $a$ and $s=3$.
Under $M$, one may take $a=1/2$.

More generally, for every fixed finite $s$, choose an odd
prime $r$ with $r-1>q_s$. The Kummer field of NT-L17
has zeros of multiplicity at least $r-1$. Its line quotient
at $a=1/2$ cannot faithfully retain the full smooth
quotient: under $M$ it loses the higher jets, and without
$M$ either that loss occurs or an off-line zero is lost.
The strict cutoff is a property of this particular
Hilbert completion, not a failure of the unconditional
adapted-block synthesis in MT-M18.

This gives the precise analytic map at the inherited
cote~1 / SA23 geo-H-lidskii and geo-H-genus interface:
the Hilbert norm can discard part of the finite-length
spectral divisor. It does not prove or refute the
conditional endpoint $M$, and it invokes none of the
draft's unverified general basis or nuclearity arguments.

\subsection*{MT-ART10. An integral Galois-module class forces a central zero}

Let $N/\mathbb Q$ be a tamely ramified $Q_8$-extension and let
$\chi$ be the unique irreducible complex character of degree two.
Fr\"ohlich's classical theorem identifies the two signs
\[
 W(\chi)=+1\quad\Longleftrightarrow\quad
       \mathcal O_N\text{ is free over }\mathbb Z[Q_8].
 \tag{MT-ART10a}
\]
The exact primary statement is Fr\"ohlich, \emph{The Galois
module structure of algebraic integer rings in fields with
generalised quaternion group} (1974), Theorem1, printed p.82,
with the locally free class invariant defined on pp.81--82.
The original proof is his \emph{Artin-root numbers and normal
integral bases for quaternion fields}, Inventiones17(1972),
143--166. Its full proof was not retrieved here; (MT-ART10a)
is used as the classical theorem, not claimed as a fresh proof.
Martinet's \emph{Fonctions z\^eta au point $1/2$}, \S VI,
printed114--115, gives the equivalent normal-basis formulation.

Consequently, if $\mathcal O_N$ is not free then
\[
 L(\tfrac12,\chi)=0,
 \qquad \operatorname{ord}_{s=1/2}\zeta_N(s)\ge2.
 \tag{MT-ART10b}
\]
These assertions are unconditional and are compatible with
both adopted GRH inputs.

\paragraph{Complete short deduction.}
Choose a cyclic subgroup $C_4\subset Q_8$ and a faithful
degree-one character $\lambda$ of $C_4$. Its induction is
$\chi$: its values are $2$ at the identity, $-2$ at the central
involution, and zero elsewhere, which is the irreducible
degree-two character. Artin induction-invariance therefore
identifies $L(s,\chi)$ with a nontrivial finite-order Hecke
function over $N^{C_4}$, so it is entire. Since $\chi$ is
selfdual, its completed function satisfies
\[
                 \Lambda(s,\chi)=W(\chi)\Lambda(1-s,\chi).
\]
Nonfreeness and (MT-ART10a) give $W(\chi)=-1$. Evaluation at
$s=1/2$ forces $\Lambda(1/2,\chi)=0$. The archimedean factors
and conductor power are finite and nonzero there, so the
finite-part $L$-function also vanishes. If $K=N^{\{1,-1\}}$
is the biquadratic subfield, the regular-character identity
gives the all-prime factorization
\[
                         \zeta_N(s)=\zeta_K(s)L(s,\chi)^2.
\]
The first factor is holomorphic at $1/2$, so the order is at
least two. More precisely the order of $L(s,\chi)$ there is
positive and odd, by the same functional equation. This proves
(MT-ART10b) without a general Artin-holomorphy hypothesis.

\paragraph{A completely specified classical example.}
Martinet, \emph{Modules sur l'alg\`ebre du groupe quaternionien}
(1971), TheoremIV.1 printed406 and Example3 printed407, gives
\[
 K=\mathbb Q(\sqrt5,\sqrt{41}),\qquad
 M=\frac{5+\sqrt5}{2}\frac{41+\sqrt{205}}{2},\qquad
 N=K(\sqrt M).
 \tag{MT-ART10c}
\]
This is a totally real tame $Q_8$-extension, ramified only at
$5$ and $41$, with nonfree integer ring. In his trace criterion
the two values are $205$ and $(1+5+41+205)/4=63$; they differ
modulo four, so the nonfree case applies. The radicand was
checked directly in the original page407 image, where the
second square root is $\sqrt{5\cdot41}$.
Thus this particular field has $\zeta_N(1/2)=0$ and satisfies
(MT-ART10b). For a convenient polynomial description, multiplying
$X^2-M$ over the four conjugates of $M$ gives
\[
 X^8-205X^6+13940X^4-378225X^2+3404025.
 \tag{MT-ART10d}
\]
An exact rational coefficient calculation of this norm is
recorded separately; the Galois group, ramification and integral
module classification are the cited classical example.

The programme bridge here is precise: tameness makes
$\mathcal O_N$ a locally free integral Galois module, and its
class in the locally free class group of $\mathbb Z[Q_8]$
controls one functional-equation sign. Rationalizing erases
this integral class while retaining the usual regular
representation. This is an example of the integral $K_0$
information distinguished in AG-LEFSCH-LATTICE and
AG-LEFSCH-DEFECT, not a general formula for all arithmetic
$K$-groups. It also shows that a residual condition such as
the nonvanishing of a Dedekind zeta function at $1/2$ cannot
be discarded merely because $M_H$ holds. It is not an example
inside the particular $S_5$ family of MT-S502; that family's
central nonvanishing gate remains its own arithmetic question.


\subsection*{MT-M26. Trace class on the specified strip Hilbert quotient}
Fix one primitive unitary Hecke function and the notation of
MT-M18 and MT-M24. Let $s\ge0$, $a\in\mathbf R$, and let
$I=[u,v]$ be a compact interval with $u<v$. Define
\[
 \begin{split}
 w_{I,s}(x)&=\max\{e^{2ux},e^{2vx}\}(1+x^2)^{s/2},\\
 H_{I,s}&=L^2(\mathbf R,w_{I,s}(x)\,dx),\\
 J_{a,I,s}&=\overline{T_a\mathcal I_\eta}^{\,H_{I,s}},
 \qquad Z_{a,I,s}=H_{I,s}/J_{a,I,s},
 \qquad T_af(x)=e^{ax}f(x).
 \end{split}                                                    \tag{M26.1}
\]
Assume there is a closed strip $[c_0,c_1]$ containing the real
parts of all zeros of $\Xi_\eta$, with
\[
                 [c_0-a,c_1-a]\subset(u,v).
                                                               \tag{M26.2}
\]
Then the natural map $Q_\eta\to Z_{a,I,s}$ is continuous,
injective and has dense image. More precisely,
\[
 J_{a,I,s}
 =\{h\in H_{I,s}:(\mathcal Mh)^{(k)}(\rho-a)=0
       \text{ for all zeros }\rho\text{ and }0\le k<m_\rho\}.
                                                               \tag{M26.3}
\]
The Mellin transforms and derivatives in (M26.3) are defined
by absolutely convergent integrals inside the strip $I$.
For every $g\in\mathcal E$, convolution induces a trace-class
operator on this \emph{actual} Hilbert quotient, and
\[
       \operatorname{Tr}_{Z_{a,I,s}}C_g
          =\sum_{\Xi_\eta(\rho)=0}
                        m_\rho(\mathcal Mg)(\rho-a),
                                                               \tag{M26.4}
\]
with absolute convergence. The trace-class estimates are
uniform when $g$ ranges over a bounded subset of $\mathcal E$.

Unconditionally one may take $[c_0,c_1]=[0,1]$, so the result
holds whenever $[0,1]\subset a+\operatorname{int}I$.
Under GRH for this one function, one may take $c_0=c_1=1/2$.
In particular, under the conditional input $M$, the result
holds at $a=1/2$ for every $I=[-\alpha,\alpha]$ with
$\alpha>0$, for each fixed Hecke character separately.

\begin{proof}
We first specify the adapted blocks needed inside the available
strip. Choose fixed real numbers
\[
        a+u<b_-<c_0\le c_1<b_+<a+v.
\]
Use the selected horizontal heights $h_n$ from MT-M18 and
rectangles $\Omega_n=(b_-,b_+)\times(h_n,h_{n+1})$.
The proof of MT-M18 gives the same block theorem for these
vertical edges, as the following estimates make explicit.
The horizontal minimum-modulus estimate there was proved
on every fixed horizontal segment, so it applies on
$[b_-,b_+]$. On either vertical edge the distance from every
zero is at least
$\delta=\min(c_0-b_-,b_+-c_1)>0$.
In the Blaschke-product proof of MT-M18, each factor on these
edges is therefore at least
$\min(1,\delta)/(2R)$, where $R$ is the fixed disk radius
used there. There are $O(\log(2+|n|))$ such factors.
Their product has a polynomial lower bound. The same
zero-free factor estimate and the two-sided Stirling bound
then give
\[
 \inf_{w\in\partial\Omega_n}|\Xi_\eta(w)|
    \ge c(2+|n|)^{-B}e^{-\kappa|n|}.
                                                               \tag{M26.5}
\]
At bounded height the edges contain no zeros and compactness
supplies the bound. The fixed width and bounded height of the
rectangles give uniformly bounded perimeter. The entire
difference-kernel construction in MT-M18, its Gaussian estimate,
and its quotient-infimum argument consequently give finite
Banach algebras $J_n$ and lifts $\mathcal T_n:J_n\to\mathcal E$
with
\[
 \begin{gathered}
 d_n=\dim J_n=O(\log(2+|n|)),\qquad
 R_m\mathcal T_n=\begin{cases}\operatorname{id}_{J_n},&m=n,\\0,&m\ne n,
                         \end{cases}\\
 p_{N,j}(\mathcal T_n z)
       \le C_{N,j}(2+|n|)^{B_{N,j}}\|z\|_{J_n}.
 \end{gathered}                                                 \tag{M26.6}
\]
Here $J_n$ is the quotient of the holomorphic supremum-norm
algebra on the closed rectangle by its full block jets.
The absolute synthesis, closed-ideal equality and two-sided
block isomorphism of MT-M18 also hold: their proof uses
(M26.5), bounded perimeter and inclusion of every zero in
exactly one rectangle, all just verified. In particular the
lifts are not unstable individual Hermite interpolation
vectors. This establishes the precise version of MT-M18
being used below.

For $h\in H_{I,s}$ its Mellin integral is holomorphic on
$u<\Re z<v$. If $\sigma$ ranges over a compact subinterval
of $(u,v)$, then for every integer $k\ge0$,
\[
 \begin{split}
 |(\mathcal Mh)^{(k)}(\sigma+it)|
 &\le\|h\|_{H_{I,s}}
   \left(\int_{\mathbf R}
     \frac{|x|^{2k}e^{2\sigma x}}{w_{I,s}(x)}\,dx\right)^{1/2}\\
 &\le C_k\|h\|_{H_{I,s}},
 \end{split}                                                    \tag{M26.7}
\]
uniformly in $t$. The integral converges by the positive
distance to both endpoints, uniformly in that compact
subinterval. The same estimate justifies differentiation
under the integral on compact subsets of the strip.
Consequently
\[
 R_n^H(h)=\bigl[w\longmapsto(\mathcal Mh)(w-a)\bigr]\in J_n
 \quad\hbox{satisfies}\quad
       \|R_n^H(h)\|_{J_n}\le C\|h\|_{H_{I,s}},
                                                               \tag{M26.8}
\]
where $C$ is independent of $n$. This follows by bounding
the quotient norm by the supremum of this representative
on the contour. For $h=T_af$ the representative is
$\mathcal Mf(w)$, so $R_n^H$ vanishes on $T_a\mathcal I_\eta$,
and by continuity on its Hilbert closure. It induces a
bounded map $r_n:Z_{a,I,s}\to J_n$ with norm at most $C$.

Let $\pi:H_{I,s}\to Z_{a,I,s}$ be the quotient map, and put
$s_n=\pi T_a\mathcal T_n$. The map $T_a:\mathcal E\to H_{I,s}$
is continuous: a seminorm $p_{N,0}$ with
$N>|a|+\max(|u|,|v|)+1$ bounds its norm by an integrable
exponential times the fixed polynomial weight. Thus (M26.6)
implies, for some fixed constants $C,B$,
\[
 \|s_n\|_{J_n\to Z_{a,I,s}}\le C(2+|n|)^B,
 \qquad r_m s_n=\delta_{mn}\operatorname{id}_{J_n}.
                                                               \tag{M26.9}
\]
The last identity follows directly from the jets of the
lifts. In particular each block is actually realized in
the Hilbert quotient.

Convolution by $g\in\mathcal E$ is bounded on $H_{I,s}$.
Indeed, with $B_I=\max(|u|,|v|)$, the elementary polynomial
weight inequality used in MT-M25 gives
\[
 \|h(\,\cdot-y)\|_{H_{I,s}}
 \le2^{s/4}e^{B_I|y|}(1+y^2)^{s/4}\|h\|_{H_{I,s}}.
                                                               \tag{M26.10}
\]
Minkowski's inequality integrates this bound against
$|g(y)|$, and the resulting integral is finite.
The subspace $T_a\mathcal I_\eta$ is invariant under this
convolution because
\[
           g*T_af=T_a\bigl((e^{-a\cdot}g)*f\bigr)
                                                               \tag{M26.11}
\]
and $\mathcal I_\eta$ is a convolution ideal in $\mathcal E$.
Its Hilbert closure is therefore invariant as well. Hence
$C_g$ induces a bounded operator on $Z_{a,I,s}$.

Let $M_{g,n}$ be multiplication in $J_n$ by the block class
of $w\mapsto\mathcal Mg(w-a)$. Submultiplicativity of the
quotient algebra norm shows
\[
 \|M_{g,n}\|\le b_n(g):=
       \sup_{w\in\partial\Omega_n}|\mathcal Mg(w-a)|.
                                                               \tag{M26.12}
\]
For each $A>0$, $b_n(g)\le C_A(g)(2+|n|)^{-A}$,
with $C_A(g)$ bounded on bounded subsets of $\mathcal E$,
by its strip-Schwartz Mellin transform. These constants
are bounded by finitely many $\mathcal E$ seminorms:
integrating by parts sufficiently many times on the
fixed strip bounds the high ordinates, and direct
integration bounds the remaining compact set. The finite-rank
operators $s_nM_{g,n}r_n$ have rank at most $d_n$ and
satisfy the Hilbert trace-norm bound
\[
 \|s_nM_{g,n}r_n\|_1
       \le C d_n(2+|n|)^B b_n(g).
                                                               \tag{M26.13}
\]
Here we used the elementary inequality
$\|T\|_1\le\operatorname{rank}(T)\|T\|$ for a finite-rank
operator on a Hilbert space. Thus the series
\[
                     \sum_n s_nM_{g,n}r_n
                                                               \tag{M26.14}
\]
converges in trace norm, uniformly for $g$ in bounded test
sets. These same estimates show that
$g\mapsto C_g$ is continuous into the trace-class norm.

It equals the actual induced convolution. To verify this
without assuming spectral completeness in the Hilbert
space, first take a vector $\pi T_af$ with $f\in\mathcal E$.
Equation (M26.11) and the absolute block synthesis in the
\emph{already proved smooth quotient} give the identity
with (M26.14) on this vector. Applying $\pi T_a$ is
legitimate because it is continuous from that smooth
quotient into the Hilbert quotient; no injectivity has
been assumed. These vectors are dense: $T_a\mathcal E$
contains $C_c^\infty$, which is dense in the weighted
$L^2$ space by truncation and compact approximation.
Both operators are bounded, so the identity holds on
all of $Z_{a,I,s}$.

We can now prove the exact closure statement (M26.3).
One inclusion follows already from (M26.7). For the other,
if every zero jet of $h\in H_{I,s}$ vanishes, then
$r_n[ h ]=0$ for every $n$. The operator identity just
proved gives $C_g[h]=0$ for every $g\in\mathcal E$.
Take a smooth compactly supported approximate identity
$g_\epsilon$ of integral one. Translation is strongly
continuous on $H_{I,s}$: this holds first for compactly
supported smooth functions, then for all functions by
their density and (M26.10). That estimate is uniformly
bounded on small $y$, so
$g_\epsilon*h\to h$ in $H_{I,s}$.
Passing to the quotient gives $[h]=0$, proving (M26.3).
In particular $Q_\eta\to Z_{a,I,s}$ is injective, since
all zero jets of a smooth vector still determine its
class by MT-M18. Continuity and density were proved above.

Finally the finite-rank trace identity and (M26.9) give
\[
 \operatorname{Tr}(s_nM_{g,n}r_n)
       =\operatorname{Tr}_{J_n}(M_{g,n}r_ns_n)
       =\operatorname{Tr}_{J_n}M_{g,n}.
\]
In the local algebra at a zero $\rho$ of multiplicity
$m_\rho$, multiplication by a holomorphic function is
triangular in the powers of $w-\rho$, with that function's
value at $\rho$ repeated $m_\rho$ times on the diagonal.
This computes the trace as
$\sum_{\rho\in\Omega_n}m_\rho\mathcal Mg(\rho-a)$.
Summing is legitimate by (M26.13), and the scalar sum
is also absolutely convergent since $d_n=O(\log(2+|n|))$
and the Mellin transform is rapidly decreasing on the
fixed strip. This proves (M26.4).
\end{proof}

The injection from the smooth quotient has proper dense
image. Indeed $\Xi_\eta$ has infinitely many nontrivial
zeros, hence infinitely many nonzero $J_n$. The rapid
block space of MT-M18 is not normable: if one norm gave
its topology, it would be bounded above by a constant
times some $q_K$, while $q_{K+1}$ would be bounded by
a constant times that norm. A unit vector in a single
block at height $n$, scaled by $(2+|n|)^{-K}$, contradicts
these two inequalities as $|n|\to\infty$. Were the
continuous injection onto the Hilbert quotient surjective,
the Fr\'echet open mapping theorem would make the smooth
quotient normable. Thus this is a faithful Hilbert
completion, not an identification of its norm topology
with the stronger smooth topology.

At $a=1/2$ put $h^*(x)=\overline{h(-x)}$. Under GRH,
(M26.4) yields
\[
 \operatorname{Tr}C_{h*h^*}
      =\sum_\rho m_\rho|\mathcal Mh(\rho-1/2)|^2\ge0.
                                                               \tag{M26.15}
\]
On any strip satisfying the unconditional sufficient
condition $[0,1]\subset1/2+\operatorname{int}I$, positivity
of these traces for every $h\in\mathcal E$ is equivalent
to GRH for the fixed function. For if
$\rho^\dagger=1-\overline\rho\ne\rho$ is an off-line pair,
the unitary functional equation gives equal multiplicities.
MT-M18 finite-block interpolation, followed by $T_{1/2}$,
produces $h$ with Mellin values $1,-1$ at
$\rho-1/2,\rho^\dagger-1/2$, respectively, and zero at
all other zeros. Formula (M26.4) then gives
$\operatorname{Tr}C_{h*h^*}=-2m_\rho$.
No positivity of the operator $C_{h*h^*}$ itself is claimed;
the faithful full-jet obstruction of MT-M20 still applies.

\paragraph{Literal source interface and remaining scope.}
Masullo's norm (9.2.2.31), original PDF p.~416, becomes
(M26.1) after $x=\log|g|$ in the fixed compact-character
sector; the supremum over subintervals is attained at
the endpoints of the compact interval $I$. His closure
Definition~10.1.1.1, p.~434, becomes $J_{a,I,s}$ by
MT-M24. For $I=[-\alpha,\alpha]$ and $a>0$, his stated
parameter range, the resulting space
is his specified quotient (10.4.1.3), original PDF p.~456,
in that sector. The convolution convention is exactly
(M24.11); an idele-class test uses the scalar function
$k_\beta$, so its trace is evaluated at the correspondingly
shifted quasi-character. Compactly supported smooth tests
are included in $\mathcal E$.

This independently supplies fixed-character spectral
trace class on those literal quotients; it does not invoke
Theorem~10.3.2.1 or the draft's general Schauder claims.
All constants may depend on the fixed field and character.
A simultaneous summation over varying compact characters
would require additional uniform conductor and archimedean
parameter estimates; that is not inferred from a
fixed-character result. Neither the trace formula nor
the space identification proves $M$. The remaining
positivity input is still conditional.

In the inherited cote~1 / SA23 geo-H-lidskii and geo-H-genus
crosswalk, the map from the completed-Hecke spectral
divisor to a specified adelic Hilbert quotient is now
explicit and faithful whenever the containing strip has
interior room. The line quotient MT-M25 explains why
that interior condition cannot simply be discarded.
